\documentclass[11pt, a4paper]{ctexart}
\usepackage{../common/ncd-style}

\hypersetup{
  pdftitle={新概念丹麦语 3 · DU3 Modul 3 全真模考四套卷},
  pdfauthor={Jaehaerys},
  pdfsubject={Ny Koncept Dansk 3 · 4 Komplette Modultests til DU3 Modul 3}
}

\begin{document}

% =====================================================================
% Book Cover Page
% =====================================================================
\begin{titlepage}
  \centering
  \vspace*{2.2cm}
  {\color{NCDRed}\rule{\linewidth}{2.5pt}}\\[1.1cm]
  {\fontsize{30}{36}\selectfont\bfseries\color{NCDNavy} 新概念丹麦语 3 · 全真模考四套卷}\\[0.5cm]
  {\LARGE\itshape\color{NCDTeal} Ny Koncept Dansk 3 · 4 Komplette Modultests til DU3 Modul 3}\\[0.8cm]
  {\large\bfseries\color{NCDGray} 冲刺实战篇（四套完整真题卷） \quad|\quad 对标 SIRI / Modultest.dk (2025/2026) 最新考纲 (B1)}\\[1.1cm]
  {\color{NCDRed}\rule{\linewidth}{2.5pt}}\\[1.8cm]

  \begin{tcolorbox}[
    colback=NCDBgBlue,
    colframe=NCDNavy,
    width=0.90\linewidth,
    arc=4pt,
    boxrule=0.8pt,
    left=14pt, right=14pt, top=12pt, bottom=12pt
  ]
    \textbf{\color{NCDNavy}本册使用定位：}效仿剑桥雅思真题集“一册四套全真卷”体例，严格依据丹麦官方 \textbf{SIRI / Gyldendal \textit{Modultest.dk} (2025/2026 现行 DU3 Modul 3 考纲)} 编写四套零重复全真模考试卷（\textit{Modultest 1--4}），与官方真题卷 A--C 及示范题构成完整互补。\\[4pt]
    每套严格覆盖：\textbf{阅读五部分（50 分钟 / 24 分制 / 15 分及格）} $\rightarrow$ \textbf{书面双任务（Opgave 1 投诉/推荐信 4 要点 \& Opgave 2 私人回信 $\ge$ 90 词 / 30 分钟）} $\rightarrow$ \textbf{口语 Mindmap 陈述与双人四图偏好互动} $\rightarrow$ \textbf{官方 Rettenøgle 评分卡 + 逐题精析 + 高频屈折词表 + B1/B1+ 双档范文与口语互动逐字稿}。
  \end{tcolorbox}

  \vfill
  {\large\color{NCDGray} 建议节奏：每 2--3 天限时完成 1 套 $\rightarrow$ 对照评分卡算分 $\rightarrow$ 精读 Del 4 语篇逻辑解析与双档范文 $\rightarrow$ 与搭档/跟读演练四图对话}
\end{titlepage}

% Table of Contents
\tableofcontents
\clearpage

% =====================================================================
% Official Exam Blueprint & Scoring Overview
% =====================================================================
\phantomsection
\addcontentsline{toc}{section}{DU3 Modul 3 官方考纲、评分标准与考场时间分配总览}
\markboth{DU3 Modul 3 考纲与评分总览}{DU3 Modul 3 考纲与评分总览}
\section*{DU3 Modul 3 官方考纲、评分标准与考场时间分配总览}

在正式计时刷四套模考卷之前，请先确认 \textbf{DU3（Danskuddannelse 3）Modul 3 (B1)} 的官方考场规则、题型结构与评分底线（依据 SIRI 与 Gyldendal \textit{Modultest.dk} 2025/2026 现行标准）：

\begin{tipbox}{三大模块考核总表（自测计时与算分基准）}
\small
\renewcommand{\arraystretch}{1.24}
\begin{tabularx}{\linewidth}{@{} >{\bfseries\color{NCDNavy}\raggedright\arraybackslash}p{2.3cm} >{\raggedright\arraybackslash}p{4.5cm} >{\raggedright\arraybackslash}p{3.6cm} >{\raggedright\arraybackslash}X @{}}
\toprule
\textbf{考核模块} & \textbf{题型与题量规格} & \textbf{计时与及格线} & \textbf{本套四卷备考建议} \\
\midrule
\textbf{Del 1}\newline Læsning\newline（阅读理解） &
五部分合订发放（禁用词典/手机）：\newline
· Opgave 1：6 空完形（0 为例词 + 6 答案 + 6 干扰，共 13 词）\newline
· Opgave 2A \& 2B (NY)：两段对话气泡补全（每段 5 条 X 发言，Z 为例句，A--F 选 3 项、3 项多余，各 3 分）\newline
· Opgave 3 (NY)：5 段短文各自 4 选 1 补句（每段 A--D，共 5 分）\newline
· Opgave 4 (NY)：3 人自述 A/B/C 匹配 7 题（0 为例题，共 7 分） &
\textbf{总计时 50 分钟}\newline
满分 \textbf{24 分}\newline
及格：\textbf{15--24 分}\newline
（0--14 分 = Ikke bestået）\newline
建议配速：8 + 10 + 14 + 13 min（留 5 min 复查） &
严格设置 50 分钟倒计时。重点关注 Modul 3 新增的 \textbf{Opgave 2A/2B}（跨回合指代与语气衔接）与 \textbf{Opgave 3}（段内因果/转折句前瞻后顾，留意 \textit{For / Men / Derfor} 信号词）。建议四套均刷到 \textbf{19 分以上}。 \\
\midrule
\textbf{Del 2}\newline Skrivning\newline（书面表达） &
· \textbf{Opgave 1}（正式/半正式应用文）：投诉信（\textit{Klage}）或好评推荐（\textit{Anbefaling}），必须覆盖题目指定的 \textbf{4 个要点}。\newline
· \textbf{Opgave 2}（私人回信）：阅读朋友邮件并回答全部问题，需得体开头与落款。 &
· Opgave 2 明确要求 \textbf{minimum 90 ord}，计时 \textbf{30 minutter}。\newline
· 建议两篇均写到 \textbf{105--135 词}，兼顾 B1 句法复杂度（整体评估 Bestået/Ikke bestået）。 &
主动嵌入 NCD 3 核心武器：\ding{172} \textit{-s} 被动 vs \textit{blive} 被动；\ding{173} 介词 + 形式宾语 \textit{det} + \textit{at/hv-} 从句；\ding{174} 无 \textit{hvis} 条件倒装与 \textit{desværre/ikke} 从句前置；\ding{175} 投诉时对机构统一用 \textit{I / jer / jeres}。 \\
\midrule
\textbf{Del 3}\newline Mundtlig\newline（口语交际） &
· \textbf{Opgave 1 (Emneopgave)}：自带 2 张纯 \textit{stikord} Mindmap（工作/学习/健康/环保/数字化/家务等），抽签 1 张陈述 + 考官追问（需简要说明理由）。\newline
· \textbf{Opgave 2 (Paropgave + Individuel)}：Del 1 两人围绕抽签话题与 4 张情景图互动讨论偏好与原因；Del 2 考官单独追问 3 个深入问题。 &
· Opgave 1：Del 1 陈述 \textbf{1--2 min} + Del 2 追问 \textbf{3--4 min}。\newline
· Opgave 2：Del 1 双人约 \textbf{4 min} + Del 2 个人约 \textbf{3 min}。\newline
（均为整体评估 \textit{Bestået}） &
\textbf{官方红线提醒}：\textit{Min bolig}、\textit{Min familie}、\textit{Min hverdag} 复杂度过低，\textbf{不可用于 Modul 3 Opgave 1}！Opgave 2 双人环节务必对\textbf{全部 4 张图}均明确表达偏好并用 \textit{fordi / derfor / synes} 简要论证。 \\
\bottomrule
\end{tabularx}
\end{tipbox}

\begin{grammatikbox}{本册四套模考卷全景导航（与官方样卷 A--C 零重复主题矩阵）}
\small
\renewcommand{\arraystretch}{1.22}
\begin{tabularx}{\linewidth}{@{} >{\bfseries\color{NCDNavy}\raggedright\arraybackslash}p{1.9cm} >{\raggedright\arraybackslash}p{5.4cm} >{\raggedright\arraybackslash}p{3.9cm} >{\raggedright\arraybackslash}X @{}}
\toprule
\textbf{试卷编号} & \textbf{阅读五题主题（Opg 1, 2A/2B, 3, 4）} & \textbf{写作双任务（Opg 1 \& 2）} & \textbf{口语话题卡（Opg 1 \& 2）} \\
\midrule
Modultest 1 & 转行故事 / 项目延期与借工具聊天 / 购买电动车 / 三人谈居家办公 & Opg 1：向 Fitness Nord 投诉\newline Opg 2：回信 Emil（骑车与环保习惯） & Opg 1：\textit{Mit arbejde}, \textit{Mine grønne vaner}\newline Opg 2：\textit{Transport i hverdagen}（4 图） \\
Modultest 2 & 断网一周实验 / 合买礼物与膝伤休跑聊天 / 加入美食社团 / 三人谈屏幕时间 & Opg 1：向 Hotel Kystblik 投诉\newline Opg 2：回信 Astrid（夜校提升与平衡） & Opg 1：\textit{Mine sunde/usunde vaner}, \textit{Mine digitale vaner}\newline Opg 2：\textit{Fritid og motion}（4 图） \\
Modultest 3 & 音乐节志愿者 / 小组报告与代收快递聊天 / 开设自行车咖啡馆 / 三人谈社团生活 & Opg 1：好评推荐 Byens Cykler\newline Opg 2：回信 Henrik（丹麦求职面试） & Opg 1：\textit{Min uddannelse}, \textit{Mit frivillige arbejde}\newline Opg 2：\textit{Indkøb og forbrug}（4 图） \\
Modultest 4 & 搬往小村庄 / 周末改期与孩子生病请假聊天 / 首次企业实习 / 三人谈家庭分工 & Opg 1：向 NordBus 客运投诉\newline Opg 2：回信 Line（丹麦适应与文化对比） & Opg 1：\textit{Arbejdsdeling i min familie}, \textit{Lære dansk}\newline Opg 2：\textit{Ferie og rejser}（4 图） \\
\bottomrule
\end{tabularx}
\end{grammatikbox}

% =====================================================================
% Four Complete Mock Exams
% =====================================================================
\ncdmodultestmthree{1}{Arbejdsliv, Hverdag og Forbrug}{全真模考卷一：职场转型、健身房投诉与环保习惯}{职场转型完形 / 项目延期与借工具聊天 / 购买电动车 / 三人谈居家办公 / Fitness Nord 投诉信与环保回信 / 工作与绿色习惯 Mindmap 及日常通勤四图口语互动}{\textit{-s} 被动与 \textit{blive} 被动 · 介词 + 形式宾语 \textit{det} + \textit{at/om} 从句 · 省略 \textit{hvis} 条件倒装 · 篇章连接词与同义转述 · 投诉语域 \textit{I/jer/jeres} 一致性}

% =====================================================================
% DEL 1: LÆSNING (50 minutter · 24 point)
% =====================================================================
\ncdsubsec{Del 1: Læsning（阅读理解｜50 minutter · 24 point）}

\begin{tipbox}{考场规则与配速建议（Læsning｜15--24 point = Bestået）}
\small
Læsning Opgave 1, 2A, 2B, 3 og 4 udleveres samtidig. Samlet tid: \textbf{50 minutter}. Kursisten må \textbf{ikke} bruge hjælpemidler（禁用任何词典、笔记或手机）. Der gives 1 point for hvert korrekt svar. \textbf{Bestået：15--24 point（及格）｜Ikke bestået：0--14 point}。建议配速：Opgave 1（8 min）$\rightarrow$ Opgave 2A/2B（10 min）$\rightarrow$ Opgave 3（14 min）$\rightarrow$ Opgave 4（13 min），最后留 5 分钟复查。
\end{tipbox}

% ---------------------------------------------------------------------
\subsubsection*{Opgave 1 (6 point): \textit{Klara skifter branche}}

\begin{emailbox}{Kursistark --- Instruktion（Opgave 1）}
\begin{itemize}
  \item Læs teksten \textit{Klara skifter branche}.
  \item Skriv de ord (1--6), der mangler. Du skal kun bruge hvert ord \textbf{én gang}. Der er \textbf{seks ord, du ikke skal bruge}. Se eksemplet (0).
\end{itemize}

\vspace{0.35em}
\hrule
\vspace{0.45em}

\textbf{\color{NCDNavy}Klara skifter branche}

\vspace{0.3em}
I ti år arbejdede Klara som laborant på en stor medicinalfabrik ved Odense. Hun kunne \textbf{(0) \underline{godt}} lide sine kollegaer, men efterhånden følte hun, at arbejdet var blevet alt for ensformigt. Hver dag lavede hun de samme prøver uden vinduer, og hun savnede noget mere konkret og grønt.

\vspace{0.25em}
En aften så hun et opslag på nettet om et mindre firma, \textbf{(1)}\ \rule{2.1cm}{0.4pt}\ hun ikke kendte i forvejen, men som leverer økologiske måltidskasser i elvarebiler. Firmaet søgte en medarbejder til at planlægge ruter og sikre, at maden kom frisk frem. Klara tøvede ikke længe, \textbf{(2)}\ \rule{2.1cm}{0.4pt}\ opgaverne lød spændende for hende.

\vspace{0.25em}
Allerede ugen efter blev hun kaldt til samtale. Chefen fortalte, at stillingen krævede overblik og god kontakt med både chauffører og lokale gårde. Klara havde \textbf{(3)}\ \rule{2.1cm}{0.4pt}\ ingen erfaring med logistik, men hun forklarede, at hun var vant til at overholde faste deadlines. Tre dage senere ringede chefen og spurgte, \textbf{(4)}\ \rule{2.1cm}{0.4pt}\ hun ville begynde den første i næste måned.

\vspace{0.25em}
Nu har Klara været i sit nye job i et halvt år. I begyndelsen var \textbf{(5)}\ \rule{2.1cm}{0.4pt}\ mange nye systemer at lære, og nogle dage havde hun travlt om morgenen. I dag klarer hun det hurtigere og har meget mere energi. Klara tjener lidt mindre end før, men hun vil \textbf{(6)}\ \rule{2.1cm}{0.4pt}\ tilbage til sit gamle arbejde.
\end{emailbox}

\begin{grammatikbox}{Ordboks（选词框 · 共 13 词：1 例词 + 6 答案 + 6 多余干扰词）}
\centering
\large\bfseries\color{NCDNavy}
\textcolor{NCDGray}{\textit{godt} {\small (0)}} \qquad fordi \qquad hvis \qquad som \qquad kunne \qquad desværre \qquad mens\\[0.25em]
der \qquad selvom \qquad om \qquad hvor \qquad aldrig \qquad men
\end{grammatikbox}

% ---------------------------------------------------------------------
\subsubsection*{Opgave 2A (NY · 3 point): \textit{Mikkel skriver til sin kollega Sara}}

\begin{emailbox}{Kursistark --- Instruktion（Opgave 2A）}
\textbf{Situation:} Mikkel skriver til sin kollega Sara.
\begin{itemize}
  \item Læs Mikkels og Saras chat. Der er 5 beskeder fra Mikkel.
  \item Find de svar fra Sara, der passer (A--F).
  \item Skriv bogstavet på det svar, der passer. Der er \textbf{tre svar, du ikke skal bruge}. Se eksemplet (0).
\end{itemize}
\end{emailbox}

\small
\renewcommand{\arraystretch}{1.22}
\begin{longtable}{@{\extracolsep{\fill}} >{\raggedright\arraybackslash}p{8.2cm} >{\raggedright\arraybackslash}p{8.2cm} @{}}
\toprule
\textbf{\color{NCDNavy}Chat mellem Mikkel og Sara} & \textbf{\color{NCDNavy}Saras svar (Z = Eksempel · vælg 3 af A--F)} \\
\midrule
\endfirsthead
\toprule
\textbf{\color{NCDNavy}Chat mellem Mikkel og Sara（续表）} & \textbf{\color{NCDNavy}Saras svar（续表）} \\
\midrule
\endhead
\bottomrule
\endfoot
\begin{tcolorbox}[enhanced, colback=NCDBgGray, colframe=NCDGray!60, boxrule=0.6pt, arc=3pt, left=7pt, right=7pt, top=5pt, bottom=5pt]
  \textbf{\color{NCDNavy}Mikkel:} Hej Sara! Har du set mailen fra vores leder om det grønne projekt?
\end{tcolorbox}
\vspace{0.2em}
\begin{tcolorbox}[enhanced, colback=NCDBgTeal, colframe=NCDTeal, boxrule=0.6pt, arc=3pt, left=7pt, right=7pt, top=4pt, bottom=4pt]
  \textbf{\color{NCDRed}0.}\ \textbf{\color{NCDTeal}Sara:} \hfill \fbox{\textbf{\ Z\ }} \textcolor{NCDGray}{\small [Eksempel]}
\end{tcolorbox}
&
\begin{tcolorbox}[enhanced, colback=NCDBgGray, colframe=NCDGray!60, boxrule=0.6pt, arc=3pt, left=7pt, right=7pt, top=5pt, bottom=5pt]
  \textbf{\color{NCDGray}Z:} \textcolor{NCDGray}{Ja, det har jeg lige læst. Vi skal aflevere rapporten to dage tidligere, ikke?} \ {\small\color{NCDRed}[Eksempel]}
\end{tcolorbox}
\\
\midrule
\begin{tcolorbox}[enhanced, colback=NCDBgGray, colframe=NCDGray!60, boxrule=0.6pt, arc=3pt, left=7pt, right=7pt, top=5pt, bottom=5pt]
  \textbf{\color{NCDNavy}Mikkel:} Jo, desværre allerede på onsdag i stedet for fredag. Tror du, vi kan nå at blive færdige?
\end{tcolorbox}
\vspace{0.2em}
\begin{tcolorbox}[enhanced, colback=NCDBgTeal, colframe=NCDTeal, boxrule=0.6pt, arc=3pt, left=7pt, right=7pt, top=4pt, bottom=4pt]
  \textbf{\color{NCDRed}1.}\ \textbf{\color{NCDTeal}Sara:} \hfill \fbox{\phantom{\textbf{M}}}
\end{tcolorbox}
&
\begin{tcolorbox}[enhanced, colback=NCDBgGold, colframe=NCDGold!70, boxrule=0.6pt, arc=3pt, left=7pt, right=7pt, top=4pt, bottom=4pt]
  \textbf{\color{NCDNavy}A:} Ja, jeg har allerede sendt rapporten til vores leder i fredags.
\end{tcolorbox}
\vspace{0.15em}
\begin{tcolorbox}[enhanced, colback=NCDBgGold, colframe=NCDGold!70, boxrule=0.6pt, arc=3pt, left=7pt, right=7pt, top=4pt, bottom=4pt]
  \textbf{\color{NCDNavy}B:} Jeg tager gerne budgettet, for jeg har allerede alle tallene i et skema.
\end{tcolorbox}
\\
\midrule
\begin{tcolorbox}[enhanced, colback=NCDBgGray, colframe=NCDGray!60, boxrule=0.6pt, arc=3pt, left=7pt, right=7pt, top=5pt, bottom=5pt]
  \textbf{\color{NCDNavy}Mikkel:} God idé. Vil du skrive den sidste del om budgettet, eller skal jeg tage den?
\end{tcolorbox}
\vspace{0.2em}
\begin{tcolorbox}[enhanced, colback=NCDBgTeal, colframe=NCDTeal, boxrule=0.6pt, arc=3pt, left=7pt, right=7pt, top=4pt, bottom=4pt]
  \textbf{\color{NCDRed}2.}\ \textbf{\color{NCDTeal}Sara:} \hfill \fbox{\phantom{\textbf{M}}}
\end{tcolorbox}
&
\begin{tcolorbox}[enhanced, colback=NCDBgGold, colframe=NCDGold!70, boxrule=0.6pt, arc=3pt, left=7pt, right=7pt, top=4pt, bottom=4pt]
  \textbf{\color{NCDNavy}C:} Nej, jeg har desværre ikke set mailen endnu, så jeg ved ikke, hvad der står i den.
\end{tcolorbox}
\vspace{0.15em}
\begin{tcolorbox}[enhanced, colback=NCDBgGold, colframe=NCDGold!70, boxrule=0.6pt, arc=3pt, left=7pt, right=7pt, top=4pt, bottom=4pt]
  \textbf{\color{NCDNavy}D:} Det bliver stramt, men hvis vi deler opgaverne op i dag, kan det nok lade sig gøre.
\end{tcolorbox}
\\
\midrule
\begin{tcolorbox}[enhanced, colback=NCDBgGray, colframe=NCDGray!60, boxrule=0.6pt, arc=3pt, left=7pt, right=7pt, top=5pt, bottom=5pt]
  \textbf{\color{NCDNavy}Mikkel:} Perfekt! Skal vi mødes online i morgen formiddag kl. 9 for at tjekke det hele igennem?
\end{tcolorbox}
\vspace{0.2em}
\begin{tcolorbox}[enhanced, colback=NCDBgTeal, colframe=NCDTeal, boxrule=0.6pt, arc=3pt, left=7pt, right=7pt, top=4pt, bottom=4pt]
  \textbf{\color{NCDRed}3.}\ \textbf{\color{NCDTeal}Sara:} \hfill \fbox{\phantom{\textbf{M}}}
\end{tcolorbox}
&
\begin{tcolorbox}[enhanced, colback=NCDBgGold, colframe=NCDGold!70, boxrule=0.6pt, arc=3pt, left=7pt, right=7pt, top=4pt, bottom=4pt]
  \textbf{\color{NCDNavy}E:} Klokken 9 har jeg et tandlægebesøg, men jeg er tilbage på kontoret klokken 10.30.
\end{tcolorbox}
\vspace{0.15em}
\begin{tcolorbox}[enhanced, colback=NCDBgGold, colframe=NCDGold!70, boxrule=0.6pt, arc=3pt, left=7pt, right=7pt, top=4pt, bottom=4pt]
  \textbf{\color{NCDNavy}F:} Det passer mig fint kl. 9, for jeg har slet ingen andre møder i morgen formiddag.
\end{tcolorbox}
\\
\midrule
\begin{tcolorbox}[enhanced, colback=NCDBgGray, colframe=NCDGray!60, boxrule=0.6pt, arc=3pt, left=7pt, right=7pt, top=5pt, bottom=5pt]
  \textbf{\color{NCDNavy}Mikkel:} Intet problem! Så holder vi mødet kl. 10.30 i stedet. Vi tales ved i morgen.
\end{tcolorbox}
&
\textcolor{NCDGray}{\small\textit{（第 5 条发言用于收尾并验证第 3 题逻辑衔接）}}
\\
\end{longtable}
\normalsize

% ---------------------------------------------------------------------
\subsubsection*{Opgave 2B (NY · 3 point): \textit{Nadia skriver til sin nabo Henrik}}

\begin{emailbox}{Kursistark --- Instruktion（Opgave 2B）}
\textbf{Situation:} Nadia skriver til sin nabo Henrik.
\begin{itemize}
  \item Læs Nadias og Henriks chat. Der er 5 beskeder fra Nadia.
  \item Find de svar fra Henrik, der passer (A--F).
  \item Skriv bogstavet på det svar, der passer. Der er \textbf{tre svar, du ikke skal bruge}. Se eksemplet (0).
\end{itemize}
\end{emailbox}

\small
\renewcommand{\arraystretch}{1.22}
\begin{longtable}{@{\extracolsep{\fill}} >{\raggedright\arraybackslash}p{8.2cm} >{\raggedright\arraybackslash}p{8.2cm} @{}}
\toprule
\textbf{\color{NCDNavy}Chat mellem Nadia og Henrik} & \textbf{\color{NCDNavy}Henriks svar (Z = Eksempel · vælg 3 af A--F)} \\
\midrule
\endfirsthead
\toprule
\textbf{\color{NCDNavy}Chat mellem Nadia og Henrik（续表）} & \textbf{\color{NCDNavy}Henriks svar（续表）} \\
\midrule
\endhead
\bottomrule
\endfoot
\begin{tcolorbox}[enhanced, colback=NCDBgGray, colframe=NCDGray!60, boxrule=0.6pt, arc=3pt, left=7pt, right=7pt, top=5pt, bottom=5pt]
  \textbf{\color{NCDNavy}Nadia:} Hej Henrik! Er du hjemme her i eftermiddag? Min bogreol i stuen er gået lidt løs.
\end{tcolorbox}
\vspace{0.2em}
\begin{tcolorbox}[enhanced, colback=NCDBgTeal, colframe=NCDTeal, boxrule=0.6pt, arc=3pt, left=7pt, right=7pt, top=4pt, bottom=4pt]
  \textbf{\color{NCDRed}0.}\ \textbf{\color{NCDTeal}Henrik:} \hfill \fbox{\textbf{\ Z\ }} \textcolor{NCDGray}{\small [Eksempel]}
\end{tcolorbox}
&
\begin{tcolorbox}[enhanced, colback=NCDBgGray, colframe=NCDGray!60, boxrule=0.6pt, arc=3pt, left=7pt, right=7pt, top=5pt, bottom=5pt]
  \textbf{\color{NCDGray}Z:} \textcolor{NCDGray}{Hej Nadia. Ja, jeg er hjemme fra kl. 16. Har du brug for hjælp?} \ {\small\color{NCDRed}[Eksempel]}
\end{tcolorbox}
\\
\midrule
\begin{tcolorbox}[enhanced, colback=NCDBgGray, colframe=NCDGray!60, boxrule=0.6pt, arc=3pt, left=7pt, right=7pt, top=5pt, bottom=5pt]
  \textbf{\color{NCDNavy}Nadia:} Ja tak, eller må jeg låne din boremaskine? Jeg skal nemlig hænge hylden fast, før jeg får gæster i aften.
\end{tcolorbox}
\vspace{0.2em}
\begin{tcolorbox}[enhanced, colback=NCDBgTeal, colframe=NCDTeal, boxrule=0.6pt, arc=3pt, left=7pt, right=7pt, top=4pt, bottom=4pt]
  \textbf{\color{NCDRed}1.}\ \textbf{\color{NCDTeal}Henrik:} \hfill \fbox{\phantom{\textbf{M}}}
\end{tcolorbox}
&
\begin{tcolorbox}[enhanced, colback=NCDBgGold, colframe=NCDGold!70, boxrule=0.6pt, arc=3pt, left=7pt, right=7pt, top=4pt, bottom=4pt]
  \textbf{\color{NCDNavy}A:} Det haster ikke, så du kan bare aflevere den i morgen efter arbejde.
\end{tcolorbox}
\vspace{0.15em}
\begin{tcolorbox}[enhanced, colback=NCDBgGold, colframe=NCDGold!70, boxrule=0.6pt, arc=3pt, left=7pt, right=7pt, top=4pt, bottom=4pt]
  \textbf{\color{NCDNavy}B:} Jeg skal bruge den hele dagen i morgen, så jeg har den desværre ikke selv i dag.
\end{tcolorbox}
\\
\midrule
\begin{tcolorbox}[enhanced, colback=NCDBgGray, colframe=NCDGray!60, boxrule=0.6pt, arc=3pt, left=7pt, right=7pt, top=5pt, bottom=5pt]
  \textbf{\color{NCDNavy}Nadia:} Tusind tak, det er virkelig sødt af dig! Hvad tid passer det bedst, at jeg banker på?
\end{tcolorbox}
\vspace{0.2em}
\begin{tcolorbox}[enhanced, colback=NCDBgTeal, colframe=NCDTeal, boxrule=0.6pt, arc=3pt, left=7pt, right=7pt, top=4pt, bottom=4pt]
  \textbf{\color{NCDRed}2.}\ \textbf{\color{NCDTeal}Henrik:} \hfill \fbox{\phantom{\textbf{M}}}
\end{tcolorbox}
&
\begin{tcolorbox}[enhanced, colback=NCDBgGold, colframe=NCDGold!70, boxrule=0.6pt, arc=3pt, left=7pt, right=7pt, top=4pt, bottom=4pt]
  \textbf{\color{NCDNavy}C:} Selvfølgelig! Du kan også låne nogle skruer, hvis du mangler nogen.
\end{tcolorbox}
\vspace{0.15em}
\begin{tcolorbox}[enhanced, colback=NCDBgGold, colframe=NCDGold!70, boxrule=0.6pt, arc=3pt, left=7pt, right=7pt, top=4pt, bottom=4pt]
  \textbf{\color{NCDNavy}D:} Det er en flot bogreol! Hvor mange gæster får du til middag i aften?
\end{tcolorbox}
\\
\midrule
\begin{tcolorbox}[enhanced, colback=NCDBgGray, colframe=NCDGray!60, boxrule=0.6pt, arc=3pt, left=7pt, right=7pt, top=5pt, bottom=5pt]
  \textbf{\color{NCDNavy}Nadia:} Super, så kommer jeg kl. 16.30. Vil du have den tilbage i aften eller i morgen?
\end{tcolorbox}
\vspace{0.2em}
\begin{tcolorbox}[enhanced, colback=NCDBgTeal, colframe=NCDTeal, boxrule=0.6pt, arc=3pt, left=7pt, right=7pt, top=4pt, bottom=4pt]
  \textbf{\color{NCDRed}3.}\ \textbf{\color{NCDTeal}Henrik:} \hfill \fbox{\phantom{\textbf{M}}}
\end{tcolorbox}
&
\begin{tcolorbox}[enhanced, colback=NCDBgGold, colframe=NCDGold!70, boxrule=0.6pt, arc=3pt, left=7pt, right=7pt, top=4pt, bottom=4pt]
  \textbf{\color{NCDNavy}E:} Jeg vil meget gerne have den tilbage allerede kl. 18, før dine gæster kommer.
\end{tcolorbox}
\vspace{0.15em}
\begin{tcolorbox}[enhanced, colback=NCDBgGold, colframe=NCDGold!70, boxrule=0.6pt, arc=3pt, left=7pt, right=7pt, top=4pt, bottom=4pt]
  \textbf{\color{NCDNavy}F:} Kom bare omkring kl. 16.30, så ligger værktøjet klar i entréen.
\end{tcolorbox}
\\
\midrule
\begin{tcolorbox}[enhanced, colback=NCDBgGray, colframe=NCDGray!60, boxrule=0.6pt, arc=3pt, left=7pt, right=7pt, top=5pt, bottom=5pt]
  \textbf{\color{NCDNavy}Nadia:} Perfekt, så kommer jeg med den i morgen ved femtiden. Vi ses om lidt!
\end{tcolorbox}
&
\textcolor{NCDGray}{\small\textit{（第 5 条发言用于收尾并验证第 3 题归还时间）}}
\\
\end{longtable}
\normalsize

% ---------------------------------------------------------------------
\subsubsection*{Opgave 3 (NY · 5 point): \textit{Kaspers nye elbil}}

\begin{emailbox}{Kursistark --- Instruktion（Opgave 3）}
\begin{itemize}
  \item Læs teksten \textit{Kaspers nye elbil}. Der er fem afsnit.
  \item I hvert afsnit mangler der en sætning. Under hvert afsnit står der fire sætninger.
  \item Skriv bogstavet på den sætning, der passer i afsnittet. Der er \textbf{tre sætninger, du ikke skal bruge}.
\end{itemize}
\end{emailbox}

\vspace{0.25em}
\begin{center}
\textbf{\large\color{NCDNavy}Kaspers nye elbil}
\end{center}

\begin{emailbox}{Afsnit 1}
\textbf{1.} Kasper arbejder som tekniker i Vejle, men bor i en lille by tyve kilometer væk. I mange år har han kørt på arbejde i en gammel benzinbil, som efterhånden bruger meget brændstof og tit skal på værksted. \underline{\hspace{2.8cm}} For det koster næsten lige så meget i reparationer hver måned, som en ny bil ville koste. Derfor begynder han at undersøge, om en brugt elbil kunne være et bedre valg for ham.
\vspace{0.35em}
\begin{itemize}[leftmargin=1.6em, itemsep=0.15em]
  \item[\textbf{A.}] Kasper er faktisk rigtig glad for sin gamle bil.
  \item[\textbf{B.}] Det er Kasper blevet ret træt af.
  \item[\textbf{C.}] Så han vil gerne beholde bilen i ti år mere.
  \item[\textbf{D.}] Men han tager alligevel toget hver morgen.
\end{itemize}
\hfill \textbf{Svar:} \rule{1.2cm}{0.4pt}
\end{emailbox}

\begin{emailbox}{Afsnit 2}
\textbf{2.} På nettet finder Kasper en tre år gammel elbil hos en forhandler i Horsens, og han bestiller en prøvetur om lørdagen. Bilen er både stille og hurtig, og sædet er meget behageligt. Kasper er dog lidt bekymret for, hvordan han skal lade bilen op i hverdagen, fordi han bor i lejlighed uden egen indkørsel. \underline{\hspace{2.8cm}} Han fortæller nemlig, at der lige er sat fire nye ladestandere op på parkeringspladsen ved supermarkedet to hundrede meter fra Kaspers gade.
\vspace{0.35em}
\begin{itemize}[leftmargin=1.6em, itemsep=0.15em]
  \item[\textbf{A.}] Heldigvis kan sælgeren berolige ham med en god nyhed.
  \item[\textbf{B.}] Derfor beslutter Kasper straks, at han slet ikke vil have en elbil.
  \item[\textbf{C.}] Men sælgeren ved ingenting om opladning i området.
  \item[\textbf{D.}] Og Kasper har allerede en ladestander i sin egen garage.
\end{itemize}
\hfill \textbf{Svar:} \rule{1.2cm}{0.4pt}
\end{emailbox}

\begin{emailbox}{Afsnit 3}
\textbf{3.} Efter prøveturen tager Kasper hjem for at regne på økonomien. En elbil er dyrere at købe end en almindelig benzinbil, så han bliver nødt til at bruge hele sin opsparing og låne lidt i banken. \underline{\hspace{2.8cm}} Strøm er nemlig meget billigere end benzin, og den grønne ejerafgift er også lav. På sigt vil Kasper faktisk kunne spare penge hver måned, og samtidig kører han mere miljøvenligt.
\vspace{0.35em}
\begin{itemize}[leftmargin=1.6em, itemsep=0.15em]
  \item[\textbf{A.}] Derfor synes hans bankrådgiver, at han skal købe en dieselbil.
  \item[\textbf{B.}] Så prisen på strøm er højere end prisen på benzin.
  \item[\textbf{C.}] Men han har slet ikke lyst til at køre miljøvenligt.
  \item[\textbf{D.}] Alligevel viser hans regnestykke, at det kan betale sig.
\end{itemize}
\hfill \textbf{Svar:} \rule{1.2cm}{0.4pt}
\end{emailbox}

\begin{emailbox}{Afsnit 4}
\textbf{4.} Mandag ringer Kasper til forhandleren i Horsens og skriver kontrakten under. Fjorten dage senere kan han endelig hente sin nye bil. Allerede den første weekend vil han teste rækkevidden og besøge sin søster, der bor på Fyn. På motorvejen glemmer han dog at holde øje med batteriet, fordi han nyder den stille kørsel og lytter til en spændende lydbog. \underline{\hspace{2.8cm}} Heldigvis ligger der en hurtiglader ved den næste afkørsel, så han når frem uden problemer.
\vspace{0.35em}
\begin{itemize}[leftmargin=1.6em, itemsep=0.15em]
  \item[\textbf{A.}] Derfor løber batteriet tør midt på motorvejen, og han må ringe efter hjælp.
  \item[\textbf{B.}] Han har nemlig ladet bilen helt op hjemme i stuen.
  \item[\textbf{C.}] Pludselig opdager han, at der kun er otte procent strøm tilbage.
  \item[\textbf{D.}] Men turen til Fyn tager kun ti minutter i alt.
\end{itemize}
\hfill \textbf{Svar:} \rule{1.2cm}{0.4pt}
\end{emailbox}

\begin{emailbox}{Afsnit 5}
\textbf{5.} Nu har Kasper haft elbilen i tre måneder, og hverdagen fungerer rigtig godt. Han har downloadet en app, som viser de billigste timer til opladning om aftenen. Mens bilen lader ved supermarkedet, klarer han som regel ugens indkøb på samme tid. \underline{\hspace{2.8cm}} Flere af hans kollegaer på arbejdet er endda begyndt at spørge ham til råds, fordi de også overvejer at skifte benzinbilen ud med en elbil.
\vspace{0.35em}
\begin{itemize}[leftmargin=1.6em, itemsep=0.15em]
  \item[\textbf{A.}] På den måde sparer han både tid og penge i en travl hverdag.
  \item[\textbf{B.}] Men han synes stadig, det er alt for besværligt at finde en ladestander.
  \item[\textbf{C.}] Derfor savner han sin gamle benzinbil rigtig meget.
  \item[\textbf{D.}] Og han må nu gå fem kilometer hver dag for at handle.
\end{itemize}
\hfill \textbf{Svar:} \rule{1.2cm}{0.4pt}
\end{emailbox}

% ---------------------------------------------------------------------
\subsubsection*{Opgave 4 (NY · 7 point): \textit{Tre personer fortæller om at arbejde hjemmefra}}

\begin{emailbox}{Kursistark --- Instruktion（Opgave 4）}
\begin{itemize}
  \item Læs de tre tekster, og læs spørgsmålene (1--7).
  \item Find den person, der passer til hvert af de syv spørgsmål.
  \item Sæt kryds (\textbf{X}) ud for den person (A, B eller C), der passer til spørgsmålet. Se eksemplet (0).
\end{itemize}
\end{emailbox}

\vspace{0.2em}
\small
\renewcommand{\arraystretch}{1.3}
\begin{longtable}{@{\extracolsep{\fill}} >{\raggedright\arraybackslash}p{10.2cm} >{\centering\arraybackslash}p{1.75cm} >{\centering\arraybackslash}p{1.75cm} >{\centering\arraybackslash}p{1.75cm} @{}}
\toprule
\textbf{Spørgsmål} & \textbf{A. Freja} & \textbf{B. Omar} & \textbf{C. Mads} \\
\midrule
\endfirsthead
\toprule
\textbf{Spørgsmål（续表）} & \textbf{A. Freja} & \textbf{B. Omar} & \textbf{C. Mads} \\
\midrule
\endhead
\bottomrule
\endfoot
0. Hvem arbejder i en offentlig institution? \textcolor{NCDGray}{\small (Eksempel)} & \textbf{X} & B & C \\
\midrule
1. Hvem spiser bedre frokost på arbejdspladsen end derhjemme? & A \rule{0.55cm}{0.4pt} & B \rule{0.55cm}{0.4pt} & C \rule{0.55cm}{0.4pt} \\
\midrule
2. Hvem har fået ondt i kroppen af at arbejde hjemme? & A \rule{0.55cm}{0.4pt} & B \rule{0.55cm}{0.4pt} & C \rule{0.55cm}{0.4pt} \\
\midrule
3. Hvem må kun arbejde hjemmefra én fast dag om ugen? & A \rule{0.55cm}{0.4pt} & B \rule{0.55cm}{0.4pt} & C \rule{0.55cm}{0.4pt} \\
\midrule
4. Hvem bliver nogle gange forstyrret af sin familie derhjemme? & A \rule{0.55cm}{0.4pt} & B \rule{0.55cm}{0.4pt} & C \rule{0.55cm}{0.4pt} \\
\midrule
5. Hvem savner den uformelle snak med sine kollegaer? & A \rule{0.55cm}{0.4pt} & B \rule{0.55cm}{0.4pt} & C \rule{0.55cm}{0.4pt} \\
\midrule
6. Hvem bruger den sparede transporttid på at dyrke motion? & A \rule{0.55cm}{0.4pt} & B \rule{0.55cm}{0.4pt} & C \rule{0.55cm}{0.4pt} \\
\midrule
7. Hvem er blevet mere effektiv på sine hjemmearbejdsdage? & A \rule{0.55cm}{0.4pt} & B \rule{0.55cm}{0.4pt} & C \rule{0.55cm}{0.4pt} \\
\end{longtable}
\normalsize

\vspace{0.25em}
\begin{center}
\textbf{\large\color{NCDNavy}Tre personer fortæller om at arbejde hjemmefra}
\end{center}

\begin{emailbox}{A. Freja}
Jeg er sagsbehandler i kommunen og sidder normalt i et stort, åbent kontor på rådhuset. I min afdeling har vi en klar regel om, at man højst må have én hjemmearbejdsdag om ugen, og for mig er det altid onsdag. Da jeg bor alene i en rolig lejlighed, er det den bedste dag på hele ugen for mig. Jeg slipper for den lange bustur i myldretiden og kan sove en halv time længere om morgenen. Sidste år købte jeg en god, ergonomisk stol og et lille skrivebord til stuen, så jeg sidder rigtig behageligt. På kontoret kommer kollegaer tit hen og stiller hurtige spørgsmål, men om onsdagen er der ingen, som afbryder mig. Derfor når jeg faktisk meget mere arbejde om onsdagen, end jeg gør på en normal dag på rådhuset. Midt på dagen laver jeg en varm og sund frokost i mit eget køkken, før jeg arbejder videre.
\end{emailbox}

\begin{emailbox}{B. Omar}
Jeg arbejder som udvikler i et privat softwarefirma i København, hvor jeg må arbejde hjemmefra tre dage om ugen. Min kæreste er sygeplejerske og arbejder på hospitalet hele dagen, så vores lille lejlighed er dejlig stille, når jeg sidder ved computeren. Det sparer mig for meget tid på cyklen, og i mine korte pauser når jeg tit at sætte en vask over eller støvsuge stuen. Desværre har vi ikke et ekstra værelse til et hjemmekontor. Jeg sidder derfor bare ved spisebordet på en almindelig træstol, og efter nogle måneder begyndte jeg at få ondt i ryggen og nakken. En anden ulempe er, at jeg nogle gange føler mig lidt ensom. Selvom vi holder møder på video hver formiddag, mangler jeg de små, tilfældige snakke med kollegaerne ved kaffemaskinen og den gode sociale energi, vi plejer at have sammen.
\end{emailbox}

\begin{emailbox}{C. Mads}
Jeg er ansat i et mellemstort revisionsfirma og kan selv bestemme, om jeg arbejder hjemme to eller tre dage om ugen, alt efter hvor mange kundemøder jeg har. Inde på vores arbejdsplads har vi en fantastisk kantine med varm mad og friske salater, men herhjemme bliver det desværre tit bare til en hurtig rugbrødsmad foran skærmen. Til gengæld elsker jeg at slippe for togturen på halvanden time hver dag. Den sparede tid bruger jeg næsten altid på at løbe en tur i skoven om morgenen, før jeg åbner computeren. Jeg har indrettet et lille hjemmekontor med en god stol, så jeg sidder fint. Min kone arbejder dog på deltid, og vores yngste søn kommer hjem fra skole klokken fjorten. Om eftermiddagen er der derfor ofte lyd fra fjernsynet eller spørgsmål til lektierne, hvilket kan gøre det svært at koncentrere sig.
\end{emailbox}

% =====================================================================
% DEL 2: SKRIVNING (Dual Tasks)
% =====================================================================
\ncdsubsec{Del 2: Skrivning（书面表达｜Opgave 1 投诉信 \& Opgave 2 私人回信）}

\subsubsection*{Opgave 1（半正式投诉信 · 必须覆盖 4 个指定要点）}

\textbf{Situation:} Du har været medlem af fitnesscentret \textit{Fitness Nord} i tre måneder, men du er utilfreds med dårlig rengøring og manglende maskiner eller lukkede hold. Du vil skrive en klage til \textit{Fitness Nord}.\\[2pt]
\textbf{Opgave:} Skriv en klage til fitnesscentret \textit{Fitness Nord}. Du skal fortælle:
\begin{itemize}
  \item hvornår du meldte dig ind
  \item hvilke problemer du har oplevet
  \item hvad du allerede har forsøgt at gøre
  \item hvordan du håber, centret vil løse problemet.
\end{itemize}

\begin{emailbox}{Klage til Fitness Nord}
\vspace{0.3em}
\rule{\linewidth}{0.4pt}\\[0.55em]
\rule{\linewidth}{0.4pt}\\[0.55em]
\rule{\linewidth}{0.4pt}\\[0.55em]
\rule{\linewidth}{0.4pt}\\[0.55em]
\rule{\linewidth}{0.4pt}\\[0.55em]
\rule{\linewidth}{0.4pt}\\[0.55em]
\rule{\linewidth}{0.4pt}\\[0.55em]
\rule{\linewidth}{0.4pt}\\[0.55em]
\rule{\linewidth}{0.4pt}\\[0.55em]
\rule{\linewidth}{0.4pt}
\end{emailbox}

\subsubsection*{Opgave 2（私人邮件回复 · Minimum 90 ord · Tid: 30 minutter）}

\textbf{Situation:} Du har fået e-mailen her nedenfor fra din ven Emil.\\
\textbf{Opgave:} Læs e-mailen, og skriv et svar til Emil. Du skal svare på alle Emils spørgsmål samt indlede og afslutte e-mailen på en passende måde.\\
\textbf{Længde:} Minimum 90 ord. \qquad \textbf{Tid:} 30 minutter.

\begin{emailbox}{Modtaget e-mail fra Emil \hfill \small\normalfont Aarhus d. 14. april 2026}
Kære ....

\vspace{0.35em}
Hvordan går det? Jeg hørte i radioen i morges, at flere danskere cykler og tænker mere på miljøet i hverdagen. Så kom jeg til at tænke på dig! Hvordan kommer du normalt på arbejde eller i skole --- cykler du, eller tager du offentlig transport? Hvad synes du er det bedste eller sværeste ved at transportere sig grønt i Danmark? Hvad gør du ellers derhjemme for at leve mere miljøvenligt i hverdagen? Og er det lettere eller sværere at have grønne vaner her end i dit hjemland?

\vspace{0.35em}
Jeg glæder mig til at høre fra dig, så skriv endelig snart!

\vspace{0.35em}
Mange hilsner\\
Emil
\end{emailbox}

\begin{emailbox}{Dit svar｜Til: Emil \qquad Emne: Grønne vaner og transport i hverdagen}
\vspace{0.3em}
\rule{\linewidth}{0.4pt}\\[0.55em]
\rule{\linewidth}{0.4pt}\\[0.55em]
\rule{\linewidth}{0.4pt}\\[0.55em]
\rule{\linewidth}{0.4pt}\\[0.55em]
\rule{\linewidth}{0.4pt}\\[0.55em]
\rule{\linewidth}{0.4pt}\\[0.55em]
\rule{\linewidth}{0.4pt}\\[0.55em]
\rule{\linewidth}{0.4pt}\\[0.55em]
\rule{\linewidth}{0.4pt}\\[0.55em]
\rule{\linewidth}{0.4pt}
\end{emailbox}

% =====================================================================
% DEL 3: MUNDTLIG KOMMUNIKATION
% =====================================================================
\ncdsubsec{Del 3: Mundtlig kommunikation（口语交际全真题卡｜Opgave 1 \& 2）}

\subsubsection*{Mundtlig Opgave 1 (Emneopgave med Mindmap｜Del 1: 1--2 min. + Del 2: 3--4 min.)}

\begin{tipbox}{DU3 Modul 3 官方口试规则提醒（Opgave 1: Emneopgave）}
\small
\begin{itemize}
  \item 考生须提前准备\textbf{两张}自带主题的 Mindmap（可带入考场），抽签决定其中一张：\textbf{Del 1} 脱稿连贯陈述 \textbf{1--2 分钟}，随后 \textbf{Del 2} 考官就该话题追问 \textbf{3--4 分钟}（要求简明给出原因，如 \textit{fordi / derfor}）。两项合并评估为 \textit{Bestået / Ikke bestået}。
  \item \textbf{框内只能写关键词或短语（stikord），严禁写完整句子！}
  \item \textbf{官方红线提醒}：\textit{Min bolig / Min lejlighed}、\textit{Min familie} 与 \textit{Min hverdag} 因语言复杂度不足，\textbf{严禁在 Modul 3 使用}！务必选择工作、学习、健康或绿色生活等 B1 主题。
\end{itemize}
\end{tipbox}

\begin{grammatikbox}{Mindmap 1A（已填关键词卡）：Mit arbejde som projektkoordinator}
\small
\renewcommand{\arraystretch}{1.34}
\begin{tabularx}{\linewidth}{@{} >{\leavevmode\bfseries\color{NCDNavy}\raggedright\arraybackslash}p{4.6cm} >{\raggedright\arraybackslash}X @{}}
\toprule
\multicolumn{2}{c}{\large\bfseries\color{NCDRed} Emne: Mit arbejde som projektkoordinator \quad\small\color{NCDGray}（仅含 stikord 关键词）} \\
\midrule
1. arbejdssted / hvor længe & grønt logistikfirma i Aarhus N · 1 år og 4 måneder · fast fuldtidsstilling \\
\midrule
2. arbejdsplads \textit{(fx antal ansatte, arbejdssprog)} & ca. 35 ansatte · lyst åbent kontor · dansk og lidt engelsk \\
\midrule
3. arbejdstider / pauser & mandag--fredag kl. 8.30--16.00 · 30 min. frokost i kantinen · fri i weekenden \\
\midrule
4. arbejdsopgaver & planlægge ruter · kundekontakt · ugentlige teammøder · skrive rapporter \\
\midrule
5. kollegaer / chef & 6 tætte kollegaer i teamet · hjælpsom chef (Morten) · fladt hierarki \\
\midrule
6. mening om jobbet \textit{(kan godt lide / ikke lide)} & godt lide: meningsfuldt \& fleksibelt · ikke lide: travle morgener \& mange mails \\
\bottomrule
\end{tabularx}

\vspace{0.45em}
\textbf{\color{NCDRed}Del 2 考官追问预测（Eksaminators opfølgende spørgsmål · 4 题）：}
\begin{enumerate}
  \item \textit{Du siger, at der er fladt hierarki på din arbejdsplads --- hvordan kan du mærke det i hverdagen?}
  \item \textit{Hvilken arbejdsopgave synes du er mest spændende, og hvad kan være svært ved at tale dansk på jobbet?}
  \item \textit{Du nævner, at morgenerne kan være travle. Hvordan håndterer du stress på arbejdet?}
  \item \textit{Kunne du tænke dig at arbejde i det samme firma om fem år, eller har du andre fremtidsplaner?}
\end{enumerate}
\end{grammatikbox}

\begin{grammatikbox}{Mindmap 1B（已填关键词卡）：Mine grønne vaner}
\small
\renewcommand{\arraystretch}{1.34}
\begin{tabularx}{\linewidth}{@{} >{\leavevmode\bfseries\color{NCDNavy}\raggedright\arraybackslash}p{4.6cm} >{\raggedright\arraybackslash}X @{}}
\toprule
\multicolumn{2}{c}{\large\bfseries\color{NCDRed} Emne: Mine grønne vaner \quad\small\color{NCDGray}（仅含 stikord 关键词）} \\
\midrule
1. mad \textit{(fx økologisk, vegetar, undgå madspild)} & mindre kød (3 kødfrie dage) · økologisk mælk \& grønt · madplan \& spise rester \\
\midrule
2. elforbrug \textit{(sparer på …)} & slukke lys \& standbystrøm · vaske ved 30 grader · bruge strøm om natten \\
\midrule
3. vandforbrug \textit{(sparer på …)} & korte bade (5 min.) · lukke vand ved tandbørstning · fylde opvaskemaskinen helt \\
\midrule
4. transport \textit{(fx cykler, offentlig transport, delebil)} & cykel til arbejde (8 km) · tog til andre byer · ingen egen bil \\
\midrule
5. genbrug \textit{(fx brugte ting, affaldssortering)} & sortere i 5 fraktioner (plast, pap, glas, metal, madaffald) · købe brugt tøj \& møbler \\
\midrule
6. andet \textit{(fx produkter uden kemi, deleøkonomi)} & miljømærket sæbe \& vaskemiddel · stofpose til indkøb · låne værktøj af naboer \\
\bottomrule
\end{tabularx}

\vspace{0.45em}
\textbf{\color{NCDRed}Del 2 考官追问预测（Eksaminators opfølgende spørgsmål · 4 题）：}
\begin{enumerate}
  \item \textit{Hvilken grøn vane synes du er nemmest at holde i hverdagen --- og hvilken er sværest?}
  \item \textit{Du nævner, at du køber brugt tøj og møbler. Hvor plejer du at finde de brugte ting?}
  \item \textit{Var du lige så opmærksom på miljø og genbrug i dit hjemland, som du er i Danmark nu?}
  \item \textit{Synes du, det er dyrere eller billigere at leve grønt i hverdagen?}
\end{enumerate}
\end{grammatikbox}

% ---------------------------------------------------------------------
\subsubsection*{Mundtlig Opgave 2 (Paropgave Del 1: ca. 4 min. + Individuel Del 2: ca. 3 min. pr. kursist)}

\begin{grammatikbox}{Opgaveark --- EMNE: Transport i hverdagen（双人四图情景提示卡）}
\small
\begin{tabularx}{\linewidth}{@{} X @{\hspace{0.4cm}} X @{}}
\begin{tcolorbox}[enhanced, colback=white, colframe=NCDTeal, boxrule=0.8pt, arc=3pt, left=7pt, right=7pt, top=6pt, bottom=6pt]
  \textbf{\color{NCDNavy}Billede 1：骑自行车通勤}\\
  {\footnotesize\color{NCDGray}[图片 1：两名佩戴头盔的上班族背着双肩包，迎着晨光骑行在宽阔平坦的城市专用自行车道上，旁边绿树成荫。]}
\end{tcolorbox}
&
\begin{tcolorbox}[enhanced, colback=white, colframe=NCDTeal, boxrule=0.8pt, arc=3pt, left=7pt, right=7pt, top=6pt, bottom=6pt]
  \textbf{\color{NCDNavy}Billede 2：乘坐火车 / 公交}\\
  {\footnotesize\color{NCDGray}[图片 2：通勤乘客坐在整洁的城际火车或城市公交车厢内，有人正戴耳机看书或用笔记本电脑，有人望向窗外。]}
\end{tcolorbox}
\\[0.35em]
\begin{tcolorbox}[enhanced, colback=white, colframe=NCDTeal, boxrule=0.8pt, arc=3pt, left=7pt, right=7pt, top=6pt, bottom=6pt]
  \textbf{\color{NCDNavy}Billede 3：拼车 / 自驾}\\
  {\footnotesize\color{NCDGray}[图片 3：三位同事或朋友合乘一辆家用汽车行驶在公路上，前排驾驶座和副驾正拿着保温咖啡杯轻松交谈。]}
\end{tcolorbox}
&
\begin{tcolorbox}[enhanced, colback=white, colframe=NCDTeal, boxrule=0.8pt, arc=3pt, left=7pt, right=7pt, top=6pt, bottom=6pt]
  \textbf{\color{NCDNavy}Billede 4：步行}\\
  {\footnotesize\color{NCDGray}[图片 4：一位穿着运动鞋的居民拎着布制购物袋，惬意地步行穿过安静的市区公园人行道前往办事。]}
\end{tcolorbox}
\end{tabularx}
\end{grammatikbox}

\begin{emailbox}{Testerark --- 官方考官引导语与 Del 2 个人追问（EMNE: Transport i hverdagen）}
\textbf{\color{NCDRed}Del 1 (Paropgave -- ca. 4 min.):}\\
Opgavearket lægges foran kursisterne, der præsenteres for emnet på følgende måde:
\begin{quote}
\textit{Emnet er ”Transport i hverdagen“. Her er nogle billeder, der viser forskellige måder at komme rundt på i hverdagen. I skal tale med hinanden om, hvad I synes om de forskellige transportformer, og hvorfor. / Kan I lide at cykle / tage toget eller bussen / køre i bil eller samkørsel / gå på arbejde eller i skole? Hvorfor? / Hvorfor ikke? I skal stille spørgsmål til hinanden og svare på hinandens spørgsmål.}
\end{quote}
\small\textcolor{NCDGray}{（若考生未就全部 4 幅图明确表态，考官须提示：\textit{Hvad synes I om at cykle / tage offentlig transport / køre i bil eller køre sammen / gå i hverdagen? Hvorfor? / Hvorfor ikke?}）}

\vspace{0.35em}
\hrule
\vspace{0.35em}
\normalsize
\textbf{\color{NCDRed}Del 2 (Individuel opgave -- ca. 3 min. pr. kursist):}\\
Tester stiller hver kursist nedenstående spørgsmål (uddybende spørgsmål stilles inden for tidsrammen):
\begin{enumerate}
  \item \textit{Hvordan kan du bedst lide at transportere dig i hverdagen? Hvorfor?}
  \item \textit{Hvilke fordele og ulemper er der ved offentlig transport i forhold til bil eller cykel?}
  \item \textit{Kan du fortælle lidt om, hvordan du plejede at komme rundt i dit hjemland sammenlignet med Danmark? (hvordan, hvor langt, nemt/svært)}
\end{enumerate}
\end{emailbox}

% =====================================================================
% DEL 4: RETTEARK, DETALJERET ANALYSE & HØJSCORE-MODELLER
% =====================================================================
\clearpage
\ncdsubsec{Del 4: Retteark, Detaljeret Analyse \& Højscore-Modeller}

\subsubsection*{4.1 官方 \textit{Retteark} 评分卡与考前自查准则}

\small
\renewcommand{\arraystretch}{1.26}
\begin{longtable}{@{\extracolsep{\fill}} >{\leavevmode\bfseries\color{NCDNavy}\raggedright\arraybackslash}p{3.2cm} >{\raggedright\arraybackslash}p{8.4cm} >{\centering\arraybackslash}p{1.6cm} >{\centering\arraybackslash}p{2.3cm} @{}}
\toprule
\textbf{考核题型} & \textbf{Rettenøgle（标准答案）} & \textbf{分值} & \textbf{自测 / 及格线} \\
\midrule
\endfirsthead
\toprule
\textbf{考核题型（续表）} & \textbf{Rettenøgle（标准答案）} & \textbf{分值} & \textbf{自测 / 及格线} \\
\midrule
\endhead
\bottomrule
\endfoot
\textbf{Læsning Opgave 1}\newline（6 题） &
1. \textbf{som} \quad 2. \textbf{fordi} \quad 3. \textbf{desværre} \quad 4. \textbf{om} \quad 5. \textbf{der} \quad 6. \textbf{aldrig}\newline
\textcolor{NCDGray}{\small （多余干扰词：hvis, hvor, selvom, men, mens, kunne｜例词 0 为 godt）} & 6 point & \rule{0.65cm}{0.4pt} / 6 \\
\midrule
\textbf{Læsning Opgave 2A}\newline（3 题） &
1. \textbf{D} \quad 2. \textbf{B} \quad 3. \textbf{E}\newline
\textcolor{NCDGray}{\small （多余干扰项：A, C, F｜例句 0 为 Z）} & 3 point & \rule{0.65cm}{0.4pt} / 3 \\
\midrule
\textbf{Læsning Opgave 2B}\newline（3 题） &
1. \textbf{C} \quad 2. \textbf{F} \quad 3. \textbf{A}\newline
\textcolor{NCDGray}{\small （多余干扰项：B, D, E｜例句 0 为 Z）} & 3 point & \rule{0.65cm}{0.4pt} / 3 \\
\midrule
\textbf{Læsning Opgave 3}\newline（5 题） &
1. \textbf{B} \quad 2. \textbf{A} \quad 3. \textbf{D} \quad 4. \textbf{C} \quad 5. \textbf{A} & 5 point & \rule{0.65cm}{0.4pt} / 5 \\
\midrule
\textbf{Læsning Opgave 4}\newline（7 题） &
1. \textbf{C} \quad 2. \textbf{B} \quad 3. \textbf{A} \quad 4. \textbf{C} \quad 5. \textbf{B} \quad 6. \textbf{C} \quad 7. \textbf{A}\newline
\textcolor{NCDGray}{\small （例题 0 为 A）} & 7 point & \rule{0.65cm}{0.4pt} / 7 \\
\midrule
\textbf{\color{NCDRed}Læsning 总计} & \textbf{Opgave 1--4 合订总评（50 分钟 · 禁用词典）} & \textbf{24 point} & \textbf{\color{NCDRed}15--24 = Bestået} \\
\midrule
\textbf{Skrivning Opgave 1}\newline（投诉信） &
完整覆盖题干 \textbf{4 个要点}（入会时间 / 问题描述 / 已采取行动 / 期望解决方案），称谓（\textit{I / jer / jeres}）一致且 B1 语法清晰。 & 整体评估 & Bestået / Ikke bestået \\
\midrule
\textbf{Skrivning Opgave 2}\newline（私人回信） &
限时 \textbf{30 分钟}，词数严格 \textbf{$\ge$ 90 词}，全面回答 Emil 全部 4 个问题，书信格式得体，主从句语序准确。 & 整体评估 & Bestået / Ikke bestået \\
\midrule
\textbf{Mundtlig Opgave 1}\newline（Emneopgave） &
自带 2 张纯 \textit{stikord} 导图（禁用 \textit{bolig/familie/hverdag}）；Del 1 连贯陈述 1--2 分钟 + Del 2 追问 3--4 分钟并简明给出理由。 & 整体评估 & Bestået / Ikke bestået \\
\midrule
\textbf{Mundtlig Opgave 2}\newline（Paropgave + Ind.） &
Del 1（约 4 分钟）：双人对\textbf{全部 4 张图}明确表达偏好、简要论证并互相提问；Del 2（约 3 分钟）：流利作答个人经验与偏好追问。 & 整体评估 & Bestået / Ikke bestået \\
\end{longtable}
\normalsize

% ---------------------------------------------------------------------
\subsubsection*{4.2 阅读 Opgave 1--4 逐题解析、干扰项排雷与全文中文精译}

\begin{answerbox}
\textbf{\color{NCDNavy}【一、Opgave 1（Klara skifter branche）完形逐题语法解析与全文精译】}
\begin{itemize}
  \item \textbf{(0) godt [Eksempel]}：固定搭配 \textit{kunne godt lide}（喜欢），在本篇中作为示例。
  \item \textbf{(1) som}：关系代词作宾语，修饰前方的机构名词 \textit{et mindre firma}，因从句已有主语 \textit{hun}（\textit{som hun ikke kendte i forvejen}），故只能用 \textbf{som}。\textbf{排除干扰项 hvor}：\textit{kende} 为及物动词（认识某家公司），此处缺宾语而非地点状语，故不能用 \textit{hvor}。
  \item \textbf{(2) fordi}：引导原因状语从句，从句保持“主语 + 动词”语序（\textit{opgaverne lød...}），解释 Klara 没有犹豫的原因。\textbf{排除干扰项 selvom / men}：工作“令人兴奋”与“没有犹豫”是因果而非转折/让步关系。
  \item \textbf{(3) desværre}：主句情态态度副词置于有限动词 \textit{havde} 之后（\textit{Klara havde desværre ingen erfaring...}），引出后文的转折 \textit{men hun forklarede...}。其余词均无法在此处兼顾句法与转折铺垫。
  \item \textbf{(4) om}：动词 \textit{spurgte} 后接间接是非问句，必须用连词 \textbf{om}（是否）。\textbf{核心排雷}：切勿误选 \textbf{hvis}（仅用于条件“如果”），这是 Modul 3 高频必考区分点！
  \item \textbf{(5) der}：句首时间状语 \textit{I begyndelsen} 触发 V2 倒装（\textit{var} 提前），不定复数名词 \textit{mange nye systemer} 需要形式主语 \textbf{der} 引导存在句（\textit{var der mange nye systemer at lære}）。
  \item \textbf{(6) aldrig}：置于情态动词 \textit{vil} 之后修饰省略了动词原形的趋向副词 \textit{tilbage}（\textit{hun vil aldrig tilbage}，她再也不想回去了）。其余副词/连词在语义或句法上均不成立。
\end{itemize}

\vspace{0.3em}
\textbf{\color{NCDTeal}[Opgave 1 全文中文精译]}\\
十年来，克拉拉一直在欧登塞附近的一家大型制药厂担任实验室技术员。她其实挺喜欢她的同事们，但渐渐地，她觉得工作变得过于单调了。每天她都在没有窗户的环境里做着同样的检测，她渴望一些更具体、更环保的工作。\\
一天晚上，她在网上看到一则招聘启事，来自一家她此前并不认识、却专门用电动货车配送有机食材餐盒的小公司。公司正在招聘一名员工来规划配送路线，并确保食物新鲜送达。克拉拉没怎么犹豫，因为这些工作任务对她来说很令人兴奋。\\
就在下一周，她接到了面试通知。主管告诉她，这个职位需要良好的大局观以及与司机和当地农场的顺畅沟通。遗憾的是，克拉拉在物流方面没有任何经验，但她解释说自己习惯了严格遵守截止日期。三天后，主管打来电话，问她是否愿意从下个月一号开始上班。\\
如今克拉拉在新岗位上已经干了半年。一开始有很多新系统需要学习，有些日子早晨很忙碌。现在她处理得快多了，活力也充沛得多。虽然克拉拉挣得比以前略少一些，但她永远也不想回到原来的工作岗位去了。

\vspace{0.45em}
\hrule
\vspace{0.45em}
\textbf{\color{NCDNavy}【二、Opgave 2A \& 2B 对话气泡衔接解析与全文精译】}
\begin{itemize}
  \item \textbf{Opgave 2A（Mikkel \& Sara）}：
    \begin{itemize}
      \item \textbf{1. Svar: D}：Mikkel 询问截止日期提前到周三“你觉得我们能按时完成吗？”只有 \textbf{D}（有点紧张，但如果今天分好工应该来得及）直接回答可行性并引出下一条 Mikkel 提议分工（\textit{God idé}）。
      \item \textbf{2. Svar: B}：Mikkel 问 Sara 想写预算部分还是由他来写。只有 \textbf{B}（我来负责预算吧，表格里数据我都有了）回应分工选择并自然接上 Mikkel 的 \textit{Perfekt!}。
      \item \textbf{3. Svar: E}：Mikkel 提议明天上午 9 点线上碰头。第 5 条 Mikkel 回复“没问题！那就改成 10:30 开会”，反向锁定 \textbf{E}（9 点我看牙医，10:30 回办公室）为唯一逻辑闭合项！
      \item \textbf{干扰项排雷}：\textbf{A} 称上周五已交报告，与第 2 条赶工矛盾；\textbf{C} 称没看过邮件，与例句 Z“刚读完”直接矛盾；\textbf{F} 说 9 点很合适，与第 5 条改期到 10:30 矛盾。
    \end{itemize}
  \item \textbf{Opgave 2B（Nadia \& Henrik）}：
    \begin{itemize}
      \item \textbf{1. Svar: C}：Nadia 问能不能借电钻把书架固定好。只有 \textbf{C}（当然可以！如果缺螺丝也可以顺便借点）正面答应并引发 Nadia 的 \textit{Tusind tak!}。
      \item \textbf{2. Svar: F}：Nadia 问几点敲门最合适，且下一条说“太好了，那我 16:30 过来”，直接锁定 \textbf{F}（16:30 左右过来吧，工具放门厅备好）。
      \item \textbf{3. Svar: A}：Nadia 问今晚还还是明天还，下一条确认“那我明天五点左右送来”，锁定 \textbf{A}（不急，明天下班后再还就行）。
      \item \textbf{干扰项排雷}：\textbf{B} 称明天全天要用且今天不在手头，与后续约定取工具矛盾；\textbf{D} 答非所问只问客人数量；\textbf{E} 要求今晚 18 点前归还，与第 5 条“明天五点还”矛盾。
    \end{itemize}
\end{itemize}

\vspace{0.3em}
\textbf{\color{NCDTeal}[Opgave 2A \& 2B 完整对话中文精译]}\\
\textbf{2A（米克尔与萨拉）}：米：“嗨萨拉！你看到主管发来的绿色项目邮件了吗？”萨（Z）：“看到了，我刚读完。我们要提前两天交报告，对吧？”米：“是的，可惜改到这周三而不是周五了。你觉得我们能赶得完吗？”萨（1-D）：“会有点紧，但如果我们今天把任务分好，应该没问题。”米：“好主意。预算那最后一部分你来写，还是我来？”萨（2-B）：“我来写预算吧，我已经有现成表格里的全部数据了。”米：“太棒了！我们明天上午九点线上碰一下核对全部内容好吗？”萨（3-E）：“九点我有个牙医预约，不过十点半我就能回到办公室。”米：“没问题！那我们改成十点半开会。明天聊。”\\[3pt]
\textbf{2B（娜迪亚与亨里克）}：娜：“嗨亨里克！你今天下午在家吗？我客厅的书架有点松了。”亨（Z）：“嗨娜迪亚。在的，我从四点起在家。你需要帮忙吗？”娜：“谢谢，或者我能借用一下你的电钻吗？今晚来客人前我得把隔板固定好。”亨（1-C）：“当然可以！如果你缺螺丝也可以借一些。”娜：“太感谢了，你真是太好心了！什么时候去敲门最合适？”亨（2-F）：“四点半左右过来就行，我把工具放门厅备好。”娜：“太好了，那我四点半过去。你希望我今晚还还是明天还？”亨（3-A）：“不着急，你明天下班后再送回来就好。”娜：“完美，那我明天五点左右送来。待会儿见！”

\vspace{0.45em}
\hrule
\vspace{0.45em}
\textbf{\color{NCDNavy}【三、Opgave 3（Kaspers nye elbil）段落补句解析与全文精译】}
\begin{itemize}
  \item \textbf{1. Svar: B}：前句说老汽油车油耗高又常修，后句用 \textit{For} 解释“每月修车费几乎快赶上一辆新车”，故中间承上启下的评价句必须是 \textbf{B}（卡斯珀对此真是厌倦透了），其中的 \textit{Det} 精准指代前文麻烦。
  \item \textbf{2. Svar: A}：前句写卡斯珀担心住公寓没自家车道怎么充电，后句用 \textit{Han fortæller nemlig...}（他解释说附近超市刚装了四个新充电桩）呼应，说明 \textit{Han} 指销售员且带来好消息，故唯一衔接句为 \textbf{A}。
  \item \textbf{3. Svar: D}：前句讲电动车买车更贵要掏空积蓄加贷款，后句以 \textit{Strøm er nemlig meget billigere...} 转入算账省钱，故必须选带让步转折词 \textit{Alligevel} 的 \textbf{D}（尽管如此，他的账目显示还是划算的）。
  \item \textbf{4. Svar: C}：前句说高速上光顾着听有声书忘了看电量，后句用 \textit{Heldigvis ligger der en hurtiglader...}（幸好下个出口有快充桩）救场，故中间突发状况只能是 \textbf{C}（突然他发现只剩百分之八的电了）。
  \item \textbf{5. Svar: A}：前句描述趁充电顺便买完一周食材，\textbf{A}（通过这种方式他在忙碌日常中既省时间又省钱）总结成效，完美衔接后文同事也向他咨询买电动车。
\end{itemize}

\vspace{0.3em}
\textbf{\color{NCDTeal}[Opgave 3 全文中文精译]}\\
\textbf{1.} 卡斯珀在瓦埃勒当技术员，却住在二十公里外的一个小镇。多年来他一直开着一辆老旧的汽油车上班，这车越来越费油，而且经常得进修理厂。\textbf{卡斯珀对此真是厌倦透了。}因为每个月花在修车上的钱几乎都快和买一辆新车差不多了。于是他开始研究，买一辆二手电动车对他来说是不是更好的选择。\\
\textbf{2.} 他在网上找到霍森斯一家车行里的一辆三年车龄的电动车，并预约了周六试驾。车开起来既安静又提速快，座椅也非常舒适。不过卡斯珀还是有点担心平时怎么充电，因为他住在公寓里，没有自己的私人车道。\textbf{幸好销售员能用一条好消息让他安心。}销售员告诉他，离卡斯珀那条街两百米远的超市停车场刚刚安装了四个新充电桩。\\
\textbf{3.} 试驾之后，卡斯珀回家仔细算了一笔经济账。电动车买起来比普通汽油车贵，所以他得动用全部积蓄，还要向银行贷一点款。\textbf{尽管如此，他的账目还是显示这笔买卖很划算。}因为电费比油费便宜得多，而且环保车主税也很低。长远来看卡斯珀每月确实能省下钱，同时出行也更环保。\\
\textbf{4.} 周一卡斯珀给霍森斯的车行打了电话并签了合同。十四天后他终于提到了新车。刚到第一个周末，他就想测试一下续航里程并顺便去菲英岛看望姐姐。然而在高速公路上，因为车里安静又在听一部精彩的有声书，他忘了留意电池电量。\textbf{突然他发现只剩下百分之八的电量了。}幸运的是下一个出口就有一个快充站，所以他顺利到达了目的地。\\
\textbf{5.} 如今卡斯珀开电动车已经三个月了，日常生活运转得非常顺畅。他下载了一款能显示晚上最便宜充电时段的应用程序。车在超市充电的功夫，他通常就顺便完成了一周的采购。\textbf{通过这种方式，他在忙碌的日常生活中既省了时间又省了钱。}甚至公司里好几位同事也开始向他请教，因为他们也在考虑把汽油车换成电动车。

\vspace{0.45em}
\hrule
\vspace{0.45em}
\textbf{\color{NCDNavy}【四、Opgave 4（居家办公对比）同义转述定位与全文精译】}
\begin{itemize}
  \item \textbf{(0) A (Freja) [Eksempel]}：\textit{kommunen / på rådhuset}（市政厅公共部门）$\leftrightarrow$ \textit{offentlig institution}。Omar 在私营软件公司，Mads 在会计事务所。
  \item \textbf{1. C (Mads)}：\textit{fantastisk kantine med varm mad og salater}（公司食堂极好）对比 \textit{herhjemme bliver det desværre tit bare til en hurtig rugbrødsmad}（在家只啃黑麦面包）。
  \item \textbf{2. B (Omar)}：\textit{begyndte jeg at få ondt i ryggen og nakken}（背痛和颈痛）同义转述题干 \textit{ondt i kroppen}。Freja 与 Mads 均有良好座椅。
  \item \textbf{3. A (Freja)}：\textit{højst må have én hjemmearbejdsdag om ugen, og for mig er det altid onsdag} $\leftrightarrow$ \textit{én fast dag om ugen}。Omar 为 3 天，Mads 为灵活的 2--3 天。
  \item \textbf{4. C (Mads)}：妻子兼职、小儿子 14 点放学回家后有电视声和作业提问（\textit{lyd fra fjernsynet eller spørgsmål til lektierne}）$\leftrightarrow$ \textit{forstyrret af sin familie}。Freja 独居，Omar 女友全天在医院。
  \item \textbf{5. B (Omar)}：\textit{mangler jeg de små, tilfældige snakke med kollegaerne ved kaffemaskinen} $\leftrightarrow$ \textit{savner den uformelle snak}。
  \item \textbf{6. C (Mads)}：省下每天 1.5 小时火车通勤后“几乎总是在早上到森林里跑步”（\textit{Den sparede tid bruger jeg... på at løbe en tur i skoven}）$\leftrightarrow$ \textit{dyrke motion}。Freja 用来多睡半小时，Omar 用来洗衣服吸尘。
  \item \textbf{7. A (Freja)}：\textit{når jeg faktisk meget mere arbejde om onsdagen, end jeg gør på en normal dag på rådhuset} $\leftrightarrow$ \textit{mere effektiv}。
\end{itemize}

\vspace{0.3em}
\textbf{\color{NCDTeal}[Opgave 4 三人自述中文精译]}\\
\textbf{A. 芙蕾雅}：我是市政厅的一名行政个案专员，平时坐在市政厅一个大型开放式办公室里。我们部门有一项明确规定，每周最多只能居家办公一天，对我来说固定在周三。由于我独自住在一间安静的公寓里，周三是我一周里最喜欢的一天。我不用在高峰期坐长途公交，早上还能多睡半小时。去年我给客厅买了一把好的人体工学椅和一张小书桌，所以坐得很舒服。在办公室同事经常走过来问这问那，但周三没人打扰我。因此我周三实际完成的工作量比平时在市政厅一天要多得多。中午我还在自己厨房做一顿热乎又健康的午餐再继续工作。\\
\textbf{B. 奥马尔}：我在哥本哈根一家私营软件公司当开发工程师，每周可以居家办公三天。我女朋友是护士，白天一整天都在医院，所以我坐在电脑前时公寓非常安静。这省去了我大量骑车通勤的时间，休息间隙我常能顺手洗一锅衣服或给客厅吸吸尘。可惜我们没有多余房间作书房，我只能坐在餐桌旁的普通木椅上，几个月后我的后背和脖子就开始疼了。另一个缺点是有时会觉得有点孤单。虽然每天上午都会开视频会议，但我还是想念茶水间咖啡机旁和同事们随意的闲聊以及平时大家在一起的社交氛围。\\
\textbf{C. 马斯}：我在一家中型会计师事务所工作，每周可以自己决定居家办公两到三天，具体取决于我手头有多少客户会议。在公司里我们有一个超棒的食堂，供应热菜和新鲜沙拉，但在家里很遗憾往往只是对着电脑随便吃片黑麦面包。不过我很喜欢免去每天来回一个半小时的火车通勤。省下来的时间我几乎总会用来早上在森林里跑步，然后再打开电脑。我布置了一个带舒适座椅的小家庭办公室，坐姿很不错。但我妻子做兼职，我们的小儿子下午两点放学回家，所以下午经常会有电视声或问作业的问题，有时会让人很难集中注意力。
\end{answerbox}

% ---------------------------------------------------------------------
\subsubsection*{4.3 核心高频词表（Gloser｜15 个 B1 完整屈折词条）}

\begin{glosertable}
名词 & en branche & \textit{(branchen, brancher, brancherne)} & 行业，领域 \\
名词 & en måltidskasse & \textit{(måltidskassen, -kasser, -kasserne)} & 食材配送箱，餐盒 \\
动词 & at tøve & \textit{(tøver, tøvede, har tøvet)} & 犹豫，迟疑 \\
名词 & en opsparing & \textit{(opsparingen, opsparinger, -erne)} & 储蓄，积蓄 \\
动词 & at berolige & \textit{(beroliger, beroligede, har beroliget)} & 使安心，安抚 \\
动词 & at lade op & \textit{(lader op, ladede op, har ladet op)} & 充电（亦作 \textit{oplade}） \\
名词 & en rækkevidde & \textit{(rækkevidden, rækkevidder, -ne)} & 续航里程，范围 \\
名词 & en ejerafgift & \textit{(ejerafgiften, -afgifter, -afgifterne)} & 车辆持有税（环保税） \\
动词 & at koncentrere sig & \textit{(koncentrerer sig, -ede sig, har -et sig)} & 集中注意力 \\
动词 & at afbryde & \textit{(afbryder, afbrød, har afbrudt)} & 打断，中断（不规则！） \\
名词 & en ulempe & \textit{(ulempen, ulemper, ulemperne)} & 缺点，不利之处 \\
形容词 & meningsfuld & \textit{(meningsfuld, meningsfuldt, meningsfulde)} & 有意义的 \\
形容词 & miljøvenlig & \textit{(miljøvenlig, miljøvenligt, miljøvenlige)} & 环保的，环境友好的 \\
名词 & et medlemskab & \textit{(medlemskabet, medlemskaber, -ne)} & 会员资格 \\
动词 & at opsige & \textit{(opsiger, opsagde, har opsagt)} & 解约，退订；辞职（不规则！） \\
\end{glosertable}

% ---------------------------------------------------------------------
\subsubsection*{4.4 写作高分双档示范（Opgave 1 \& Opgave 2）}

\begin{emailbox}{Opgave 1 范文一：B1 稳过版（107 ord｜四要点齐全 · \textit{I/jer/jeres} 一致）}
Kære Fitness Nord

\vspace{0.25em}
Jeg skriver til \da{jer}, fordi jeg er utilfreds med forholdene i \da{jeres} center. Jeg meldte mig ind i januar, så jeg har været medlem i tre måneder.

\vspace{0.25em}
Jeg har oplevet to store problemer. Omklædningsrummene er ofte beskidte om aftenen, og rengøringen virker \dabold{ikke god} nok. Derudover har tre løbebånd været i stykker i flere uger, og mit yogahold \dabold{blev aflyst} to gange sidste uge.

\vspace{0.25em}
Jeg har allerede talt med en medarbejder i receptionen, som lovede at give beskeden videre, men intet er ændret.

\vspace{0.25em}
Jeg håber, at \da{I} hurtigt vil reparere maskinerne og forbedre rengøringen. Ellers må jeg desværre skifte fitnesscenter.

\vspace{0.25em}
Venlig hilsen\\
Freja Lind

\vspace{0.3em}
\rule{\linewidth}{0.4pt}\\[0.2em]
\small\textbf{\color{NCDTeal}[考官点评 · 107 ord]}：完整覆盖四大要点（1 月入会、卫生与坏机器/取消课程两大问题、已向前台反映但未果、希望尽快维修保洁）。对健身房统一使用复数称谓 \textit{I / jer / jeres}，准确使用 \textit{blev aflyst} 过去被动与从句语序，逻辑清晰稳过。
\end{emailbox}

\begin{emailbox}{Opgave 1 范文二：B1+ 高分冲刺版（139 ord｜\textit{-s} vs \textit{blive} 被动 + 介词 \textit{det at} + 无 \textit{hvis} 倒装）}
Til ledelsen hos Fitness Nord

\vspace{0.25em}
Jeg skriver til \da{jer}, fordi jeg efter tre måneders medlemskab desværre må klage. Jeg meldte mig ind den 12. januar og havde glædet mig til at træne hos \da{jer} ugentligt.

\vspace{0.25em}
I den seneste måned har rengøringen i omklædningsrummene dog været alt for dårlig, og gulvet \dabold{vaskes} tydeligvis ikke dagligt. Samtidig har fire maskiner været i stykker længe, mens flere holdtimer \dabold{er blevet lukket} uden varsel. Jeg er især \da{skuffet over det, at} medlemmerne ikke får besked i forvejen.

\vspace{0.25em}
Jeg skrev til \da{jeres} kundeservice for to uger siden og spurgte, \da{om} fejlene snart \dabold{bliver rettet}, men jeg har stadig ikke fået svar.

\vspace{0.25em}
\dabold{Bliver problemerne ikke løst} inden månedens udgang, bliver jeg desværre nødt til at opsige mit medlemskab. Jeg håber derfor, at \da{I} vil give mig halv pris næste måned.

\vspace{0.25em}
Med venlig hilsen\\
Freja Lind

\vspace{0.3em}
\rule{\linewidth}{0.4pt}\\[0.2em]
\small\textbf{\color{NCDNavy}[高分亮点 · 139 ord]}：\ding{172} 精准区分 \textit{-s} 被动（\textit{gulvet vaskes} 常规保洁制度）与 \textit{blive} 被动（\textit{er blevet lukket / bliver rettet} 具体事件过程）；\ding{173} 运用介词 + 形式宾语 \textit{det} + 从句（\textit{skuffet over det, at...}）及间接疑问 \textit{spurgt, om...}；\ding{174} 末段以省略 \textit{hvis} 的条件倒装句（\textit{Bliver problemerne ikke løst..., bliver jeg...}）收束诉求，语域正式克制，极具 B1+ 冲击力。
\end{emailbox}

\begin{emailbox}{Opgave 2 范文一：B1 稳过版（109 ord｜严格满足 $\ge$ 90 词 · 全面回答 4 问）}
Kære Emil

\vspace{0.25em}
Tak for din søde mail! Det går rigtig godt med mig.

\vspace{0.25em}
Til daglig cykler jeg næsten altid på arbejde, fordi det kun tager tyve minutter. Det bedste ved at cykle i Danmark er de brede cykelstier overalt. Det sværeste er dog vintervejret med blæst og regn, så der tager jeg nogle gange bussen.

\vspace{0.25em}
Derhjemme sorterer jeg alt affald i glas, plast og papir, spiser mindre kød og køber ofte brugt tøj. Jeg synes, det er meget lettere at have grønne vaner i Danmark end i mit hjemland, fordi her er så gode muligheder for at cykle og genbruge.

\vspace{0.25em}
Hvad med dig? Cykler du også?

\vspace{0.25em}
Kærlig hilsen\\
Sara

\vspace{0.3em}
\rule{\linewidth}{0.4pt}\\[0.2em]
\small\textbf{\color{NCDTeal}[考官点评 · 109 ord]}：远超 90 词底线，自然分段应答 Emil 的全部 4 个问题（日常骑车通勤、自行车道优势与冬雨难点、家里垃圾分类与减肉买二手、丹麦与母国对比），人称呼应、V2 倒装与书信收尾完全规范。
\end{emailbox}

\begin{emailbox}{Opgave 2 范文二：B1+ 高分冲刺版（130 ord｜无 \textit{hvis} 倒装 + 被动切换 + \textit{glad for det, at}）}
Kære Emil

\vspace{0.25em}
Tusind tak for din besked! Jeg har det godt og håber, du har haft en god uge.

\vspace{0.25em}
Til daglig cykler jeg året rundt, men \dabold{regner det} kraftigt om morgenen, \dabold{tager jeg} heldigvis letbanen. Det bedste ved grøn transport i Danmark er, at cykelstierne \dabold{bliver ryddet og saltet} om vinteren. Til gengæld er modvinden nogle gange hård på lange ture.

\vspace{0.25em}
Derhjemme køber jeg sjældent nyt tøj, men finder brugte ting online, og så \dabold{sorteres} alt vores affald nøje. Jeg er faktisk rigtig \da{glad for det, at} man kan aflevere flasker i supermarkedet og få pant tilbage. Sammenlignet med mit hjemland er det langt nemmere at leve grønt her, fordi både priser og byplanlægning gør de grønne valg helt naturlige.

\vspace{0.25em}
Skriv endelig, hvordan du selv kommer rundt!

\vspace{0.25em}
Kærlig hilsen\\
Sara

\vspace{0.3em}
\rule{\linewidth}{0.4pt}\\[0.2em]
\small\textbf{\color{NCDNavy}[高分亮点 · 130 ord]}：\ding{172} 省略 \textit{hvis} 条件倒装（\textit{regner det kraftigt..., tager jeg heldigvis letbanen}）；\ding{173} \textit{blive} 被动（\textit{bliver ryddet og saltet} 市政除雪撒盐作业）与 \textit{-s} 被动（\textit{sorteres} 家庭日常分类规则）自然呼应；\ding{174} 介词 + 形式宾语 \textit{det} + 从句（\textit{glad for det, at man kan aflevere flasker...}）地道呈现丹麦押金回收制度（\textit{pant}）。
\end{emailbox}

% ---------------------------------------------------------------------
\subsubsection*{4.5 口语示范逐字稿、核心 IPA 与双人四图互动示范}

\begin{transbox}
\small
\textbf{\color{NCDNavy}[Mindmap 1A 口语陈述示范：Mit arbejde som projektkoordinator｜142 ord · 约 1.5 分钟]}\\
\textit{Jeg vil gerne fortælle om mit arbejde som projektkoordinator \ipa{pʰʁoˈjɛɡdˌkʰoʌdiˌnæːtˢɐ} i et grønt logistikfirma \ipa{loɡiˈsdiɡˌfiɐ̯mɑ} i Aarhus Nord. Jeg har arbejdet der i et år og fire måneder i en fast fuldtidsstilling. Vi er cirka femogtredive ansatte \ipa{ˈanˌsadə} på et lyst, åbent kontor, hvor arbejdssproget hovedsageligt er dansk, men vi taler også lidt engelsk med internationale partnere. Jeg arbejder mandag til fredag fra klokken halv ni til fire, og vi holder en halv times frokostpause sammen i kantinen klokken tolv. Mine vigtigste arbejdsopgaver er at planlægge ruter for vores elvarebiler, have kontakt til kunderne, deltage i ugentlige teammøder og skrive korte rapporter. Jeg har seks rigtig søde kollegaer i mit team og en hjælpsom chef, som hedder Morten. Jeg kan virkelig godt lide, at arbejdet er meningsfuldt og fleksibelt, og at der er fladt hierarki \ipa{hiaˈkʰiːˀ}, selvom morgenerne nogle gange er lidt travle med mange e-mails.}
\end{transbox}

\paragraph*{Mindmap 1A Del 2 考官追问与高分应答示范（4 题）：}
\dialogturn{1}{Eksaminator}{Du siger, at der er fladt hierarki på din arbejdsplads --- hvordan kan du mærke det i hverdagen?}{hierarki \ipa{hiaˈkʰiːˀ} · hverdagen \ipa{ˈvɛːɐˌdæːˀn}}{你说你们工作场所层级很扁平——在日常生活中你从哪里能感受到这一点？}
\dialogturn{2}{Kursist}{Det mærker jeg tydeligt, fordi man altid kan tale direkte med chefen om nye idéer. Hvis jeg er uenig i en plan, lytter Morten faktisk til mine argumenter, og alle spiser frokost ved samme bord.}{tydeligt \ipa{ˈtˢyðəlið} · argumenter \ipa{ɑːɡuˈmɛnˀdɐ}}{我感受非常明显，因为随时都可以直接跟老板讨论新想法。如果我对某个计划有不同意见，Morten 会认真听取我的理由，而且大家都在同一张桌子上吃午饭。}
\dialogturn{3}{Eksaminator}{Hvilken arbejdsopgave synes du er mest spændende, og hvad kan være svært ved at tale dansk på jobbet?}{spændende \ipa{ˈsbɛnənə} · svært \ipa{ˈsvɛːˀɐd}}{哪项工作任务你觉得最有趣？在工作中说丹麦语有什么难点？}
\dialogturn{4}{Kursist}{Jeg synes, det er mest spændende at tale med lokale gårde og planlægge grønne ruter. Det sværeste er nogle gange telefonsamtaler, når folk taler hurtigt i støj, men så spørger jeg bare høfligt, om de vil gentage det.}{telefonsamtaler \ipa{tˢeləˈfoːˀnˌsɑmˌtˢæːˀlɐ} · gentage \ipa{ˈɡɛnˌtˢæːə}}{我觉得最有趣的是和本地农场沟通并规划绿色路线。有时最难的是打电话，当对方在嘈杂环境里说得很快时有点费劲，但我会礼貌地请他们再重复一遍。}
\dialogturn{5}{Eksaminator}{Du nævner, at morgenerne kan være travle. Hvordan håndterer du stress på arbejdet?}{travle \ipa{ˈtˢʁɑu̯lə} · håndterer \ipa{hʌnˈtˢeːˀɐ}}{你提到早晨可能会很忙。你在工作中是如何应对压力的？}
\dialogturn{6}{Kursist}{Når der kommer mange mails på én gang, laver jeg en hurtig liste over de vigtigste opgaver først. Desuden tager jeg en kort pause på fem minutter med en kop kaffe, så jeg kan bevare overblikket.}{vigtigste \ipa{ˈveɡdisdə} · overblikket \ipa{ˈʌu̯ɐˌbleɡəð}}{当大量邮件同时涌来时，我会先列一个最核心任务的清单。此外，我会喝杯咖啡休息五分钟，这样就能保持清晰的条理。}
\dialogturn{7}{Eksaminator}{Kunne du tænke dig at arbejde i det samme firma om fem år, eller har du andre fremtidsplaner?}{fremtidsplaner \ipa{ˈfʁɑmtˢiðsˌpʰlæːˀnɐ} · samme \ipa{ˈsɑmə}}{五年后你还愿意在同一家公司工作吗，还是有其他未来计划？}
\dialogturn{8}{Kursist}{Ja, jeg kunne godt tænke mig at blive, fordi grøn logistik vokser hurtigt, og jeg lærer noget nyt hele tiden. Måske vil jeg på sigt gerne have endnu mere ansvar for større projekter.}{ansvar \ipa{ˈanˌsvɑːˀ} · projekter \ipa{pʰʁoˈjɛɡdɐ}}{是的，我挺想留下来的，因为绿色物流发展很快，我也一直在学新东西。也许长远来看，我希望承担更大项目的责任。}

\begin{transbox}
\small
\textbf{\color{NCDNavy}[Mindmap 1B 口语陈述示范：Mine grønne vaner｜140 ord · 约 1.5 分钟]}\\
\textit{Nu vil jeg gerne fortælle om mine grønne vaner \ipa{ˈɡʁønə ˈvæːnɐ} i hverdagen. Jeg prøver at tænke på miljøet \ipa{milˈjøːˀəð} på flere måder. Når det gælder mad, spiser jeg mindre kød end før og har tre kødfrie dage om ugen. Jeg køber tit økologisk \ipa{økoˈloːˀisɡ} mælk og grøntsager og laver en fast madplan, så jeg undgår madspild \ipa{ˈmæðˌsbilˀ}. For at spare på el og vand slukker jeg altid lyset og fjernsynet helt, vasker tøj ved tredive grader om natten og tager korte bade på fem minutter. I min daglige transport cykler jeg otte kilometer til arbejde hver dag, og når jeg skal til andre byer, tager jeg toget, fordi jeg ikke ejer en bil. Derudover sorterer vi vores affald \ipa{ˈɑu̯ˌfalˀ} i fem fraktioner og køber ofte brugt tøj og brugte møbler. Endelig bruger jeg miljømærket sæbe og har altid en stofpose med, når jeg handler ind.}
\end{transbox}

\paragraph*{Mindmap 1B Del 2 考官追问与高分应答示范（4 题）：}
\dialogturn{1}{Eksaminator}{Hvilken grøn vane synes du er nemmest at holde i hverdagen --- og hvilken er sværest?}{nemmest \ipa{ˈnɛməsd} · sværest \ipa{ˈsvɛːɐsd}}{你觉得日常生活中哪项绿色习惯最容易坚持——哪项最难？}
\dialogturn{2}{Kursist}{Det nemmeste er at sortere affald og cykle, fordi systemet i Danmark fungerer så godt. Det sværeste er nogle gange at undgå madspild, når man har travlt i ugens løb og glemmer rester i køleskabet.}{madspild \ipa{ˈmæðˌsbilˀ} · køleskabet \ipa{ˈkʰøːləˌsgæːˀbəð}}{最容易的是垃圾分类和骑车，因为丹麦的配套体系非常完善。有时最难的是避免食物浪费，平时工作一忙就容易忘记冰箱里的剩菜。}
\dialogturn{3}{Eksaminator}{Du nævner, at du køber brugt tøj og møbler. Hvor plejer du at finde de brugte ting?}{brugt \ipa{ˈbʁuɡd} · plejer \ipa{ˈpʰlɑi̯ɐ}}{你提到你会买二手衣服和家具。你通常去哪里找这些二手物品？}
\dialogturn{4}{Kursist}{Jeg går tit i genbrugsbutikker som Røde Kors i weekenden, og så bruger jeg også apps på nettet til at finde brugte møbler i mit nærområde. Man kan ofte finde rigtig god kvalitet til en billig pris.}{genbrugsbutikker \ipa{ˈɡɛnˌbʁuːsbuˌtˢigɐ} · kvalitet \ipa{kʰvaliˈtˢeːˀd}}{周末我常去像红十字会那样的二手公益店，也会用网上的二手应用在附近淘家具，往往能以很实惠的价格买到质量非常好的东西。}
\dialogturn{5}{Eksaminator}{Var du lige så opmærksom på miljø og genbrug i dit hjemland, som du er i Danmark nu?}{opmærksom \ipa{ʌbˈmɛɐ̯ɡsʌmˀ} · hjemland \ipa{ˈjɛmˌlanˀ}}{你在母国时也像现在在丹麦这样关注环保和回收吗？}
\dialogturn{6}{Kursist}{Ikke helt så meget som nu. I mit hjemland var der ikke så mange cykelstier eller beholdere til sortering overalt. Efter jeg flyttede til Danmark, blev jeg hurtigt inspireret af de gode muligheder her.}{beholdere \ipa{beˈhʌlˀɐ} · inspireret \ipa{ensbiˈʁeːˀɐð}}{以前没有现在这么强。在我母国，不像这里到处都有自行车道和分类垃圾桶。搬到丹麦后，我很快就被这里便利的环保设施带动起来了。}
\dialogturn{7}{Eksaminator}{Synes du, det er dyrere eller billigere at leve grønt i hverdagen?}{dyrere \ipa{ˈdyːɐ} · billigere \ipa{ˈbiliɐ}}{你觉得在日常生活中过绿色生活是更贵还是更省钱？}
\dialogturn{8}{Kursist}{Økologisk mad kan være lidt dyrere, men til gengæld sparer jeg mange penge på at spise mindre kød, købe brugte ting og cykle i stedet for at have en bil. Så samlet set er det faktisk billigere.}{samlet set \ipa{ˈsɑmləð ˈseːˀd}}{有机食品可能会略贵一点，但反过来少吃肉、买二手货以及骑车代替养车帮我省下了很多钱。所以总体而言其实更省钱。}

\vspace{0.35em}
\hrule
\vspace{0.35em}
\paragraph*{Mundtlig Opgave 2 Del 1（Paropgave：Transport i hverdagen · 完整四图双人互动示范）：}

\dialogturn{1}{Kursist A}{Hej! Lad os kigge på billederne om transport i hverdagen. Hvad synes du om billede et, altså at cykle på arbejde?}{transport \ipa{tˢʁɑnˈsbɒːd} · cykle \ipa{ˈsyɡlə}}{你好！我们来看看日常交通这几张图吧。你对第一张图——也就是骑车上班——怎么看？}
\dialogturn{2}{Kursist B}{Jeg kan rigtig godt lide at cykle! Jeg cykler næsten hver dag, fordi det er billigt, sundt og hurtigere end bilen i myldretiden. Men hvad med dig --- cykler du også, når det regner?}{myldretiden \ipa{ˈmylʁəˌtˢiːˀðən}}{我非常喜欢骑车！我几乎每天都骑，因为既便宜健康，高峰期又比汽车快。那你呢——下雨的时候你也骑吗？}
\dialogturn{3}{Kursist A}{Ja, jeg cykler også, hvis jeg har regntøj på, fordi jeg elsker den friske luft om morgenen. Men hvis det stormer meget, tager jeg hellere offentlig transport som på billede to. Hvad synes du om toget og bussen?}{offentlig \ipa{ˈʌfənˀdli} · stormer \ipa{ˈsdɒːmɐ}}{是的，只要穿上雨衣我也会骑，因为我喜欢早晨的新鲜空气。但如果刮大风，我宁愿像图二那样坐公共交通。你觉得坐火车和公交怎么样？}
\dialogturn{4}{Kursist B}{Jeg synes, toget er rigtig rart, især på lange ture, for man kan læse eller slappe af undervejs. Til gengæld kan jeg ikke så godt lide bussen, fordi den tit sidder fast i trafikken eller bliver forsinket. Har du selv bil, eller bruger du samkørsel som på billede tre?}{samkørsel \ipa{ˈsɑmˌkʰøɐ̯səl} · forsinket \ipa{fʌˈseŋˀɡəð}}{我觉得火车很舒适，尤其长途出行时可以在路上看书或放松。不过我不太喜欢公交车，因为它经常堵在路上或者晚点。你自己有车吗，还是会像图三那样拼车？}
\dialogturn{5}{Kursist A}{Jeg har ikke bil, fordi det er alt for dyrt i Danmark. Men jeg synes faktisk, samkørsel med kollegaer er en smart idé, for man sparer benzin og kan snakke sammen. Kunne du tænke dig at køre i bil til arbejde hver dag?}{benzin \ipa{bɛnˈsiːˀn}}{我没有车，因为在丹麦养车太贵了。但我其实觉得跟同事拼车是个聪明的主意，既省油又能聊天。你会想每天开车上班吗？}
\dialogturn{6}{Kursist B}{Nej, helst ikke hver dag, kun hvis jeg bor langt ude på landet. Men hvad siger du til det sidste billede med at gå? Kan du godt lide at gå til dine aftaler?}{langt ude \ipa{ˈlɑŋd ˈuːðə}}{不了，最好不要每天开，除非住在偏远乡下。那最后一张步行的图你怎么看？你喜欢走路去赴约或办事吗？}
\dialogturn{7}{Kursist A}{Ja, jeg elsker at gå, når afstanden er under to kilometer, fordi det giver ro i hovedet, og man ser mere af byen. Men når jeg har travlt eller skal købe tungt ind, tager jeg cyklen i stedet. Hvad foretrækker du?}{afstanden \ipa{ˈɑu̯ˌsdanˀən} · foretrækker \ipa{ˈfɒːəˌtˢʁɛɡɐ}}{是的，只要距离在两公里以内我就很喜欢步行，因为能让大脑放松还能多看看城市。不过赶时间或者采购重物时我就改骑车。你更偏好哪种？}
\dialogturn{8}{Kursist B}{Jeg har det på præcis samme måde! At gå i parken er skønt i weekenden, men til daglig foretrækker jeg helt klart cyklen.}{foretrækker \ipa{ˈfɒːəˌtˢʁɛɡɐ}}{我跟你想的一模一样！周末在公园散步太棒了，但日常通勤我绝对首选骑自行车。}

\paragraph*{Mundtlig Opgave 2 Del 2（个人深度问答 3 题高分示范）：}
\begin{enumerate}
  \item \textbf{Spørgsmål 1:} \textit{Hvordan kan du bedst lide at transportere dig i hverdagen? Hvorfor?}\\
  \textbf{Modelsvar:} \textit{Jeg kan bedst lide at cykle i hverdagen, fordi det giver mig både frihed og daglig motion. Jeg behøver ikke at vente på en bus eller sidde i bilkø, og jeg ved altid præcis, hvornår jeg er fremme. Hvis jeg skal langt væk, tager jeg dog gerne toget, fordi det er behageligt.}
  \item \textbf{Spørgsmål 2:} \textit{Hvilke fordele og ulemper er der ved offentlig transport i forhold til bil eller cykel?}\\
  \textbf{Modelsvar:} \textit{Fordelen ved offentlig transport er, at man kan bruge tiden på at læse, lytte til podcast eller slappe af, og man slipper for at lede efter parkering. Ulempen er, at billetter ofte er dyre i Danmark, og hvis man bor uden for de store byer, kører busserne ikke så tit som en privat bil.}
  \item \textbf{Spørgsmål 3:} \textit{Kan du fortælle lidt om, hvordan du plejede at komme rundt i dit hjemland sammenlignet med Danmark?}\\
  \textbf{Modelsvar:} \textit{I mit hjemland tog jeg næsten altid metro eller bus, fordi afstandene i byen var meget store, og trafikken var alt for farlig til cykler. Da jeg kom til Danmark, blev jeg positivt overrasket over de brede cykelstier og respekten for cyklister. I dag synes jeg, det er meget nemmere og tryggere at cykle her.}
\end{enumerate}

\ncdmodultestmthree{2}{Sundhed, Fritid og Digitalisering}{全真模考卷二：数字排毒、酒店投诉与健康习惯}{一周无社交媒体挑战 / 合买纪念礼物与膝伤休跑聊天 / 社区美食社团 / 三人谈手机与屏幕时间 / 酒店投诉信与夜间提升课程回信 / 健康习惯与数字生活思维导图及运动休闲四图互动}{复合从句连词与词序（fordi/før/at 否定与状语前置）· 跨回合指代与语气衔接 · 因果/转折句段内前瞻后顾 · 被动语态（-s vs blive）· 介词 + 形式宾语 det + 从句 · 书面复合介词（på baggrund af / til trods for）}

% =====================================================================
% DEL 1: LÆSNING
% =====================================================================
\ncdsubsec{Del 1: Læsning（阅读理解｜50 minutter · 24 point）}

\begin{tipbox}{考场规则与配速建议（Læsning｜15--24 point = Bestået）}
\small
本部分共五道模块题，考场将 \textbf{Opgave 1, 2A, 2B, 3 og 4} 一次性合订发放。限时 \textbf{50 分钟}，全程\textbf{禁用任何词典、笔记或电子设备}。每题 1 分，满分 \textbf{24 分}，及格线为 \textbf{15 分}（0--14 分 = \textit{Ikke bestået}）。建议配速：Opgave 1（8 分钟）$\rightarrow$ Opgave 2A \& 2B（10 分钟）$\rightarrow$ Opgave 3（14 分钟）$\rightarrow$ Opgave 4（13 分钟），最后留 5 分钟复查。
\end{tipbox}

% ---------------------------------------------------------------------
\subsubsection*{Opgave 1 (6 point)：En uge uden sociale medier}

\textbf{Instruktion:} Læs teksten \textit{En uge uden sociale medier}. Skriv de ord (1--6), der mangler. Du skal kun bruge hvert ord én gang. Der er seks ord, du ikke skal bruge. Se eksemplet (0).

\vspace{0.35em}
\textbf{\color{NCDNavy}En uge uden sociale medier}

\vspace{0.2em}
Mikkel er niogtyve år og arbejder som grafisk designer i Roskilde. I mange år brugte han næsten tre timer hver dag på sociale medier, og (0)\ \underline{\textbf{det}} begyndte langsomt at irritere ham. Når vækkeuret ringede om morgenen, tjekkede han straks sine beskeder, (1)\ \rule{2.0cm}{0.4pt} han overhovedet stod ud af sengen. På arbejdet mistede han ofte fokus, (2)\ \rule{2.0cm}{0.4pt} skærmen hele tiden lyste med nye notifikationer og videoklip. Derfor besluttede han sidste måned at lave et eksperiment: Han ville slette alle sociale apps fra telefonen i syv dage for at mærke forskellen.

De første to døgn var overraskende hårde. Når Mikkel sad i toget eller ventede på kaffen i kantinen, søgte hånden næsten automatisk ned i lommen efter mobilen. Efter tre dage blev alting (3)\ \rule{2.0cm}{0.4pt} meget roligere. Uden den konstante strøm af billeder sov han bedre om natten og fik pludselig tid (4)\ \rule{2.0cm}{0.4pt} at læse en roman, som havde ligget på natbordet i et halvt år. Han mærkede også tydeligt, (5)\ \rule{2.0cm}{0.4pt} han talte mere med sine kolleger i pauserne i stedet for at stirre ned i skærmen. Da ugen endelig var slut, hentede Mikkel kun to apps tilbage og slog alle notifikationer fra. Nu bruger han højst tyve minutter om dagen på nettet, og det har (6)\ \rule{2.0cm}{0.4pt} ham mere energi i hverdagen.

\vspace{0.35em}
\begin{grammatikbox}{Ordboks（选词框 · 共 13 词：例词 (0) 已用 + 6 答案 + 6 多余干扰词）}
\centering
\large\bfseries\color{NCDNavy}
at \qquad blev \qquad {\color{NCDGray}\textit{det}\textsuperscript{(0)}} \qquad dog \qquad eller \qquad fordi \qquad fra \\[3pt]
før \qquad givet \qquad med \qquad men \qquad som \qquad til
\end{grammatikbox}

% ---------------------------------------------------------------------
\subsubsection*{Opgave 2A (NY · 3 point)：Lukas skriver til sin veninde Emma}

\textbf{Instruktion:}\\
\textbf{Situation:} Lukas skriver til sin veninde Emma.
\begin{itemize}
  \item Læs Lukas' og Emmas chat. Der er 5 beskeder fra Lukas.
  \item Find de svar fra Emma, der passer (A--F).
  \item Skriv bogstavet på det svar, der passer.
\end{itemize}
\hspace*{1.5em}Der er tre svar, du ikke skal bruge.\\
\hspace*{1.5em}Se eksemplet (0).

\vspace{0.25em}
\small
\renewcommand{\arraystretch}{1.35}
\begin{longtable}{@{} >{\raggedright\arraybackslash}p{8.1cm} @{\hspace{0.5cm}} >{\raggedright\arraybackslash}p{8.1cm} @{}}
\toprule
\textbf{\color{NCDNavy}Chat mellem Lukas og Emma} & \textbf{\color{NCDNavy}Svar fra Emma (Z \& A--F)} \\
\midrule
\endfirsthead
\toprule
\textbf{\color{NCDNavy}Chat mellem Lukas og Emma（续表）} & \textbf{\color{NCDNavy}Svar fra Emma（续表）} \\
\midrule
\endhead
\bottomrule
\endfoot
\badge{NCDNavy}{Lukas 1}\ Hej Emma! Signe og Kasper har jo fem års jubilæum næste lørdag. Har du fundet på en god gave fra os alle sammen endnu?\par\vspace{4pt}
\textbf{0. Emma:}\ \framebox[0.85cm]{\textbf{Z}} \quad \textcolor{NCDGray}{\small (Eksempel)} &
\badge{NCDGray}{Z}\ \textcolor{NCDGray}{\textit{Hej! Ikke helt endnu, men jeg tænkte på et gavekort til en god restaurant. Hvad siger du til det?}} \\
\midrule
\badge{NCDNavy}{Lukas 2}\ Ja, det er en super idé! De elsker jo begge to italiensk mad. Hvor meget synes du, at hver person skal betale til gaven?\par\vspace{4pt}
\textbf{1. Emma:}\ \framebox[0.85cm]{\rule{0pt}{0.26cm}} &
\badge{NCDNavy}{A}\ Nej, de aner ingenting! Kaspers bror lukker os stille ind kl. 14.30, mens de er ude at handle.\par\vspace{6pt}
\badge{NCDNavy}{B}\ Hvad med 100 kroner hver? Så bliver det ikke for dyrt for nogen, og vi skal jo også have noget lækkert til kaffen. \\
\midrule
\badge{NCDNavy}{Lukas 3}\ Det lyder helt fair, så bliver det i alt 800 kroner, hvis vi er otte personer. Skal jeg også bestille kagen hos bageren?\par\vspace{4pt}
\textbf{2. Emma:}\ \framebox[0.85cm]{\rule{0pt}{0.26cm}} &
\badge{NCDNavy}{C}\ Nej, vi behøver slet ikke nogen kage, for vi har allerede spist for meget sukker i denne uge.\par\vspace{6pt}
\badge{NCDNavy}{D}\ Det er ikke nødvendigt at give 100 kroner hver, for jeg betaler bare hele gaven og kagen alene. \\
\midrule
\badge{NCDNavy}{Lukas 4}\ Okay, så ringer jeg til bageriet i morgen formiddag og bestiller en chokoladekage uden nødder. Hvad med overraskelsen -- ved de, at vi kommer forbi lørdag eftermiddag?\par\vspace{4pt}
\textbf{3. Emma:}\ \framebox[0.85cm]{\rule{0pt}{0.26cm}} &
\badge{NCDNavy}{E}\ Ja, det må du meget gerne gøre, for du bor lige ved siden af bageriet. Kan du bestille en chokoladekage uden nødder?\par\vspace{6pt}
\badge{NCDNavy}{F}\ Jo, de ved godt, at vi kommer kl. 15, for jeg har allerede sendt dem en invitation. \\
\midrule
\badge{NCDNavy}{Lukas 5}\ Perfekt, så holder jeg det selvfølgelig helt hemmeligt, og jeg skriver lige til de andre gæster om mødetidspunktet. Vi ses på lørdag kl. 14.30! &
\textcolor{NCDGray}{\small （第 5 条为 Lukas 的收尾回复，无需补全）} \\
\end{longtable}
\normalsize

% ---------------------------------------------------------------------
\subsubsection*{Opgave 2B (NY · 3 point)：Sofie skriver til sin løbetræner Peter}

\textbf{Instruktion:}\\
\textbf{Situation:} Sofie skriver til sin løbetræner Peter.
\begin{itemize}
  \item Læs Sofies og Peters chat. Der er 5 beskeder fra Sofie.
  \item Find de svar fra Peter, der passer (A--F).
  \item Skriv bogstavet på det svar, der passer.
\end{itemize}
\hspace*{1.5em}Der er tre svar, du ikke skal bruge.\\
\hspace*{1.5em}Se eksemplet (0).

\vspace{0.25em}
\small
\renewcommand{\arraystretch}{1.35}
\begin{longtable}{@{} >{\raggedright\arraybackslash}p{8.1cm} @{\hspace{0.5cm}} >{\raggedright\arraybackslash}p{8.1cm} @{}}
\toprule
\textbf{\color{NCDNavy}Chat mellem Sofie og Peter} & \textbf{\color{NCDNavy}Svar fra Peter (Z \& A--F)} \\
\midrule
\endfirsthead
\toprule
\textbf{\color{NCDNavy}Chat mellem Sofie og Peter（续表）} & \textbf{\color{NCDNavy}Svar fra Peter（续表）} \\
\midrule
\endhead
\bottomrule
\endfoot
\badge{NCDNavy}{Sofie 1}\ Hej Peter! Jeg har desværre fået lidt ondt i mit højre knæ efter vores lange løbetur i søndags. Tror du, jeg kan være med til intervalløb i aften?\par\vspace{4pt}
\textbf{0. Peter:}\ \framebox[0.85cm]{\textbf{Z}} \quad \textcolor{NCDGray}{\small (Eksempel)} &
\badge{NCDGray}{Z}\ \textcolor{NCDGray}{\textit{Hej Sofie! Det er jeg ked af at høre. Hvis du har smerter, skal du helt sikkert holde en pause fra hårdt løb i aften.}} \\
\midrule
\badge{NCDNavy}{Sofie 2}\ Okay, det har du nok ret i, selvom det er irriterende at gå glip af træningen. Hvor længe skal jeg vente, før jeg må begynde at løbe igen?\par\vspace{4pt}
\textbf{1. Peter:}\ \framebox[0.85cm]{\rule{0pt}{0.26cm}} &
\badge{NCDNavy}{A}\ Ja, tag bare med til intervalløb i aften alligevel, så går smerterne nok væk af sig selv, når du bliver varm.\par\vspace{6pt}
\badge{NCDNavy}{B}\ Du skal helst sidde helt stille på sofaen i syv dage og undgå både vand og strækøvelser. \\
\midrule
\badge{NCDNavy}{Sofie 3}\ En hel uge uden løb lyder lidt langt, men jeg vil jo heller ikke have en alvorlig skade. Må jeg slet ikke bevæge mig i den uge, eller kan jeg træne noget andet i stedet?\par\vspace{4pt}
\textbf{2. Peter:}\ \framebox[0.85cm]{\rule{0pt}{0.26cm}} &
\badge{NCDNavy}{C}\ Start meget langsomt på fladt underlag, og husk at tjekke, om dine løbesko er for gamle og skal skiftes ud.\par\vspace{6pt}
\badge{NCDNavy}{D}\ Jeg vil anbefale dig at tage en uges pause fra løb, indtil knæet slet ikke gør ondt længere. \\
\midrule
\badge{NCDNavy}{Sofie 4}\ Svømning og udstrækning lyder faktisk som en god plan, og der ligger heldigvis en svømmehal tæt på mit arbejde. Er der noget særligt, jeg skal huske, før jeg prøver at løbe en kort tur igen?\par\vspace{4pt}
\textbf{3. Peter:}\ \framebox[0.85cm]{\rule{0pt}{0.26cm}} &
\badge{NCDNavy}{E}\ Jeg foreslår, at du løber en hurtig testtur på ti kilometer i morgen tidlig for at mærke, om knæet holder.\par\vspace{6pt}
\badge{NCDNavy}{F}\ Jo, du må meget gerne holde kroppen i gang. Prøv at svømme et par gange og lave rolige strækøvelser derhjemme, fordi det ikke belaster knæet. \\
\midrule
\badge{NCDNavy}{Sofie 5}\ Tak for de gode råd, Peter! Så køber jeg et par nye sko på fredag og starter stille ud i næste uge. Vi ses til træning snart! &
\textcolor{NCDGray}{\small （第 5 条为 Sofie 的收尾回复，无需补全）} \\
\end{longtable}
\normalsize

% ---------------------------------------------------------------------
\subsubsection*{Opgave 3 (NY · 5 point)：Amira starter i en madklub}

\textbf{Instruktion:}
\begin{itemize}
  \item Læs teksten \textit{Amira starter i en madklub}.
  \item Der er fem afsnit.
  \item I hvert afsnit mangler der en sætning.
  \item Under hvert afsnit står der fire sætninger.
  \item Skriv bogstavet på den sætning, der passer i afsnittet.\\
  Der er tre sætninger, du ikke skal bruge.
\end{itemize}

\vspace{0.25em}
\textbf{\color{NCDNavy}Amira starter i en madklub}

\vspace{0.2em}
\textbf{1.} Amira er for nylig flyttet til Odense for at begynde på et nyt arbejde som bioanalytiker. I starten kendte hun næsten ingen i byen, og efter arbejde sad hun ofte alene i sin lejlighed med en hurtig ret fra supermarkedet. \rule{2.4cm}{0.4pt} En dag så hun et opslag på det lokale bibliotek om en madklub, hvor beboere i kvarteret mødtes hver anden uge for at lave mad sammen. \quad \textbf{Svar:} \rule{1.0cm}{0.4pt}
\begin{enumerate}[label=\textbf{\Alph*.}, leftmargin=2.2em, itemsep=0.12em]
  \item Det var hun rigtig glad for, fordi hun foretrak at spise i ro og mag derhjemme.
  \item Hun savnede både at lave rigtig mad fra bunden og at lære nye mennesker at kende.
  \item Derfor inviterede hun sine mange gamle naboer i Odense på middag hver fredag.
  \item Men hun havde faktisk slet ingen tid til at gå på biblioteket eller lave mad.
\end{enumerate}

\vspace{0.25em}
\textbf{2.} Klubben bestod af otte mennesker i forskellige aldre, som alle boede i det samme kvarter. Første gang Amira kom til mødet, var hun en smule nervøs, fordi hun ikke vidste, om hun ville passe ind i gruppen. \rule{2.4cm}{0.4pt} Allerede efter ti minutter stod hun ved komfuret og grinede sammen med en ældre herre, der hedder Jens, mens de skar grøntsager til en stor suppe. \quad \textbf{Svar:} \rule{1.0cm}{0.4pt}
\begin{enumerate}[label=\textbf{\Alph*.}, leftmargin=2.2em, itemsep=0.12em]
  \item Heldigvis blev hun mødt med stor varme og fik straks en opgave i køkkenet.
  \item Desværre havde de andre medlemmer allerede spist, da hun nåede frem.
  \item Derfor besluttede hun at tage hjem igen med det samme uden at hilse på nogen.
  \item Men ingen af deltagerne havde lyst til at snakke med nye medlemmer.
\end{enumerate}

\vspace{0.25em}
\textbf{3.} I madklubben skiftes medlemmerne til at bestemme menuen og købe ind til aftenen. Da det blev Amiras tur i november, ville hun gerne vise de andre en traditionel ret fra sit hjemland. \rule{2.4cm}{0.4pt} Amira var lidt bange for, at krydderierne ville være for stærke for de danske deltagere, men gryderetten blev hurtigt spist op, og flere bad straks om opskriften. \quad \textbf{Svar:} \rule{1.0cm}{0.4pt}
\begin{enumerate}[label=\textbf{\Alph*.}, leftmargin=2.2em, itemsep=0.12em]
  \item Hun valgte derfor at bestille pizza udefra i stedet for at stå i køkkenet.
  \item De andre sagde dog nej, fordi de kun ønskede at spise klassiske danske retter.
  \item Hun foreslog dog først at lave mad fra et helt andet land, som hun ikke kendte.
  \item Derfor fandt hun en gammel familieopskrift frem og brugte hele eftermiddagen på at forberede krydderierne.
\end{enumerate}

\vspace{0.25em}
\textbf{4.} Ud over selve maden har aftenerne i klubben også haft en anden vigtig betydning for Amira. Når de otte deltagere sidder omkring bordet og taler om arbejde, nyheder og hverdagsliv, får hun trænet sit danske sprog i en afslappet stemning. \rule{2.4cm}{0.4pt} På den måde føler hun sig nu meget mere sikker, når hun skal diskutere eller småsnakke med kollegerne i laboratoriet. \quad \textbf{Svar:} \rule{1.0cm}{0.4pt}
\begin{enumerate}[label=\textbf{\Alph*.}, leftmargin=2.2em, itemsep=0.12em]
  \item Men de øvrige medlemmer vil desværre helst tale engelsk med hende hele aftenen.
  \item Desuden lærer hun mange nye hverdagsudtryk, som ikke står i hendes lærebøger.
  \item Derfor har hun valgt at droppe sit danskkursus og kun læse bøger derhjemme.
  \item Hun synes nemlig, at det er alt for kedeligt at snakke om hverdagen over maden.
\end{enumerate}

\vspace{0.25em}
\textbf{5.} Nu har Amira været medlem af madklubben i næsten et år, og tirsdag aften er blevet hendes yndlingsdag i kalenderen. Selvom hun stadig arbejder meget og nogle gange kommer træt hjem, melder hun sjældent afbud. \rule{2.4cm}{0.4pt} Faktisk er hun allerede i gang med at planlægge næste måneds menu sammen med Jens. \quad \textbf{Svar:} \rule{1.0cm}{0.4pt}
\begin{enumerate}[label=\textbf{\Alph*.}, leftmargin=2.2em, itemsep=0.12em]
  \item Hun regner dog med at melde sig ud næste måned for at få mere tid til sig selv.
  \item Hvis hun har været træt på arbejde, bliver hun næsten altid hjemme på sofaen.
  \item Hun oplever nemlig, at de fælles middage giver hende fornyet energi til resten af ugen.
  \item Jens og de andre har nemlig besluttet at lukke madklubben efter jul.
\end{enumerate}

% ---------------------------------------------------------------------
\subsubsection*{Opgave 4 (NY · 7 point)：Tre personer fortæller om deres forhold til mobil og skærmtid}

\textbf{Instruktion:}
\begin{itemize}
  \item Læs de tre tekster, og læs spørgsmålene (1--7).
  \item Find den person, der passer til hvert af de syv spørgsmål.
  \item Sæt kryds (\textbf{X}) ud for den person (A, B eller C), der passer til spørgsmålet. Se eksemplet (0).
\end{itemize}

\vspace{0.2em}
\small
\renewcommand{\arraystretch}{1.3}
\begin{longtable}{@{} >{\raggedright\arraybackslash}p{11.0cm} @{\hspace{0.3cm}} >{\centering\arraybackslash}p{1.65cm} @{\hspace{0.15cm}} >{\centering\arraybackslash}p{1.65cm} @{\hspace{0.15cm}} >{\centering\arraybackslash}p{1.65cm} @{}}
\toprule
\textbf{Spørgsmål} & \textbf{A. Ida} & \textbf{B. Tobias} & \textbf{C. Lin} \\
\midrule
\endfirsthead
\toprule
\textbf{Spørgsmål（续表）} & \textbf{A. Ida} & \textbf{B. Tobias} & \textbf{C. Lin} \\
\midrule
\endhead
\bottomrule
\endfoot
0. Hvem ser næsten ikke på en skærm i løbet af sin arbejdsdag? \textcolor{NCDGray}{\small (Eksempel)} & A & \textbf{X} & C \\
\midrule
1. Hvem lader ikke sin telefon ligge i soveværelset om natten længere? & A & B & C \\
\midrule
2. Hvem bruger mobilen til at øve sig i et fremmedsprog? & A & B & C \\
\midrule
3. Hvem opdagede sit store skærmforbrug, fordi et familiemedlem sagde noget om det? & A & B & C \\
\midrule
4. Hvem har en aftale med sine venner om at lægge telefonerne væk, når de ses? & A & B & C \\
\midrule
5. Hvem synes, at digitale enheder gør det lettere at bevare kontakten til folk langt væk? & A & B & C \\
\midrule
6. Hvem sætter tid af på én bestemt ugedag til helt at slukke for mobilen? & A & B & C \\
\midrule
7. Hvem sover bedre nu, efter at personen ændrede sine aftenregler? & A & B & C \\
\end{longtable}
\normalsize

\vspace{0.3em}
\begin{center}
\textbf{\large\color{NCDNavy} Tre personer fortæller om deres forhold til mobil og skærmtid}
\end{center}

\begin{emailbox}{A. Ida}
Jeg arbejder som projektkoordinator i et mellemstort firma i Aarhus, hvor jeg får beskeder på både computer og telefon hele dagen lang. Tidligere tog jeg ofte arbejdet med i seng, fordi jeg lige ville tjekke de sidste mails inden sengetid. Det ødelagde min nattesøvn fuldstændigt, for hovedet fortsatte med at tænke på opgaver hele natten. For et halvt år siden lavede jeg derfor faste aftenregler: Efter klokken nitten svarer jeg aldrig på jobmails, og min mobil bliver liggende ude i gangen hele natten, mens et gammelt vækkeur vækker mig om morgenen. Jeg falder meget hurtigere i søvn nu og vågner faktisk frisk og udhvilet. Derudover har mine veninder og jeg indført en fast aftale, når vi mødes på café efter arbejde: Alle telefoner bliver lagt midt på bordet på lydløs, indtil vi går hjem, så vi virkelig er til stede sammen.
\end{emailbox}

\begin{emailbox}{B. Tobias}
Som tømrer i Vejle arbejder jeg udendørs eller på værksted det meste af dagen, så jeg kigger næsten aldrig på en skærm i arbejdstiden. Til gengæld opdagede jeg sidste efterår, at jeg brugte over to timer hver eneste aften hjemme i sofaen på at se sportsvideoer og læse nyheder på nettet. Det var faktisk min tiårige søn, der fik mig til at tænke over det, da han en dag spurgte, hvorfor jeg altid kiggede ned i mobilen, når han ville snakke med mig. Den bemærkning ramte mig ret hårdt. Nu lægger jeg altid telefonen på en høj hylde i køkkenet fra klokken fem til otte om aftenen, så vi kan spise aftensmad og spille brætspil sammen i fred og ro. Desuden har jeg indført en fast mobilfri søndag, hvor min telefon er helt slukket fra morgen til aften, mens hele familien cykler eller fisker.
\end{emailbox}

\begin{emailbox}{C. Lin}
Jeg er medicinstuderende på universitetet i Odense og bruger digitale enheder rigtig mange timer i døgnet. Alle mine lærebøger og forelæsningsnoter ligger nemlig på min tablet, så jeg kan slet ikke undvære en skærm i hverdagen. Samtidig bruger jeg mobilen aktivt til at øve dansk, når jeg tager den lange bustur til campus hver morgen: Jeg lytter til danske nyhedspodcasts og træner nye ord i en sprogapp. Mange siger, at telefoner kun skaber afstand mellem mennesker, men for mig er det helt omvendt. Min nærmeste familie bor i Singapore, og fordi jeg ringer til mine forældre på video hver søndag og sender små billeder fra mit liv i Danmark, føler jeg mig tæt på dem på trods af den store afstand. Jeg tjekker stadig mine beskeder sent om aftenen i soveværelset på grund af tidsforskellen, uden at det skader mit studie næste morgen.
\end{emailbox}

% =====================================================================
% DEL 2: SKRIVNING
% =====================================================================
\ncdsubsec{Del 2: Skrivning（书面表达｜Opgave 1 Klage \& Opgave 2 E-mail）}

\begin{tipbox}{DU3 Modul 3 书面双任务考场提示}
\small
本部分包含两道写作任务：\textbf{Opgave 1}（正式/半正式投诉或推荐信，必须完整覆盖全部 4 个要点；建议写到 105--135 词并用 \textit{I / jer / jeres} 指代机构）；\textbf{Opgave 2}（回复朋友邮件，官方硬性规定 \textbf{Minimum 90 ord · Tid: 30 minutter}，需包含得体开头、结尾并回答全部提问）。
\end{tipbox}

\subsubsection*{Opgave 1（正式投诉信 · Klage til Hotel Kystblik）}

\textbf{Situation:}\\
Du boede i weekenden på \textit{Hotel Kystblik}, men opholdet var skuffende. Der var kraftig støj fra gaden hele natten, vandet i bruseren var koldt, rengøringen på værelset var dårlig, og receptionisten var uhøflig, da du sagde til.

\vspace{0.3em}
\textbf{Opgave:}\\
Skriv en klage til \textit{Hotel Kystblik}. Du skal fortælle:
\begin{itemize}
  \item hvornår du boede der
  \item hvilke problemer der opstod
  \item hvordan personalet reagerede
  \item hvad du forventer som kompensation
\end{itemize}

\begin{emailbox}{Din klage｜Til: Hotel Kystblik \qquad Emne: Klage over ophold}
\vspace{0.3em}
\rule{\linewidth}{0.4pt}\\[0.55em]
\rule{\linewidth}{0.4pt}\\[0.55em]
\rule{\linewidth}{0.4pt}\\[0.55em]
\rule{\linewidth}{0.4pt}\\[0.55em]
\rule{\linewidth}{0.4pt}\\[0.55em]
\rule{\linewidth}{0.4pt}\\[0.55em]
\rule{\linewidth}{0.4pt}
\end{emailbox}

\clearpage
\subsubsection*{Opgave 2（私人回信 · Minimum 90 ord · Tid: 30 minutter）}

\textbf{Situation:}\\
Du har fået e-mailen her nedenfor.

\vspace{0.3em}
\textbf{Opgave:}\\
Læs e-mailen og skriv et svar til Astrid. Du skal svare på alle Astrids spørgsmål. Du skal indlede og afslutte e-mailen på en passende måde.\\
\textbf{Længde:} Minimum 90 ord. \qquad \textbf{Tid:} 30 minutter.

\begin{emailbox}{Modtaget e-mail fra Astrid}
Hej!

Jeg hørte fra fælles venner, at du er startet på et spændende aftenkursus eller en ny fritidsinteresse efter arbejde! Det lyder virkelig sejt, men også lidt travlt. Hvilket kursus eller hvilken aktivitet går du til, og hvorfor valgte du det? Hvordan når du at balancere dit arbejde, din danskundervisning og din fritid i hverdagen? Og hvad gør du for at slappe af i weekenden, når ugen har været hård? Jeg glæder mig til at høre fra dig!

Kærlig hilsen\\
Astrid
\end{emailbox}

\begin{emailbox}{Dit svar｜Til: Astrid \qquad Emne: Sv: Aftenkursus og hverdag}
\vspace{0.3em}
\rule{\linewidth}{0.4pt}\\[0.55em]
\rule{\linewidth}{0.4pt}\\[0.55em]
\rule{\linewidth}{0.4pt}\\[0.55em]
\rule{\linewidth}{0.4pt}\\[0.55em]
\rule{\linewidth}{0.4pt}\\[0.55em]
\rule{\linewidth}{0.4pt}\\[0.55em]
\rule{\linewidth}{0.4pt}\\[0.55em]
\rule{\linewidth}{0.4pt}\\[0.55em]
\rule{\linewidth}{0.4pt}
\end{emailbox}

% =====================================================================
% DEL 3: MUNDTLIG KOMMUNIKATION
% =====================================================================
\ncdsubsec{Del 3: Mundtlig kommunikation（口语交际全真题卡｜Opgave 1 \& 2）}

\subsubsection*{Mundtlig Opgave 1 (Emneopgave med Mindmap｜Del 1: 1--2 min. + Del 2: 3--4 min.)}

\begin{tipbox}{考场规则提醒（DU3 Modul 3 Opgave 1 · 难度与关键词红线）}
\small
考生需自带\textbf{两张}提前填写好的思维导图（至少涉及工作/学习/医疗健康/环保/数字化/家务分工等达到 B1 复杂度的主题）。导图框内\textbf{只能写关键词或短语（stikord），严禁写完整句子}，以防考场照本宣科！抽签选定 1 张后，先进行 \textbf{1--2 分钟}连贯陈述（Del 1），随后回答考官 \textbf{3--4 分钟}追问（Del 2）并简要说明原因。\\
\textbf{官方禁用红线}：\textit{Min bolig/lejlighed}、\textit{Min familie} 与 \textit{Min hverdag} 复杂度停留在 A2，\textbf{不可作为 Modul 3 Opgave 1 话题}！
\end{tipbox}

\begin{grammatikbox}{Mindmap 2A（已填关键词卡）：Mine sunde og usunde vaner}
\small
\renewcommand{\arraystretch}{1.35}
\begin{tabularx}{\linewidth}{@{} >{\leavevmode\bfseries\color{NCDNavy}\raggedright\arraybackslash}p{3.1cm} @{\hspace{0.2cm}} >{\raggedright\arraybackslash}X @{\hspace{0.35cm}} >{\leavevmode\bfseries\color{NCDNavy}\raggedright\arraybackslash}p{3.1cm} @{\hspace{0.2cm}} >{\raggedright\arraybackslash}X @{}}
\toprule
\multicolumn{4}{c}{\large\bfseries\color{NCDRed} Emne: Mine sunde og usunde vaner \quad\small\color{NCDGray}（仅含 stikord 关键词，不含整句）} \\
\midrule
1. sund mad & havregryn · frisk frugt og grøntsager · fuldkornsbrød · fisk 2 gange ugentligt · meget vand &
2. usund mad & chokolade om aftenen · fredagsslik · fastfood ved travlhed · op til 4 kopper kaffe \\
\midrule
3. motion & cykle 8 km til arbejde · svømmehal onsdag aften · lang gåtur i skoven om søndagen &
4. rygning / alkohol & aldrig røget · ingen cigaretter · 1 glas vin i weekenden · alkoholfrit øl til hverdag \\
\midrule
5. søvn & ca. 7 timers søvn · fast sengetid kl. 23 · ingen mobil efter kl. 22 · træt mandag morgen &
6. mening om din livsstil (kan godt lide / ikke lide) & tilfreds med motion og mad · god balance · vil gerne skære ned på sukker · ikke for streng \\
\bottomrule
\end{tabularx}
\end{grammatikbox}

\textbf{Del 2 考官追问模拟（Mine sunde og usunde vaner · 4 题）：}
\begin{enumerate}
  \item \textit{Du siger, at du spiser lidt for meget chokolade om aftenen -- hvornår har du mest lyst til usunde ting, og hvorfor?}
  \item \textit{Hvorfor kan du godt lide at svømme og cykle til arbejde, og hvordan holder du motivationen om vinteren?}
  \item \textit{Du nævner, at du lægger telefonen væk klokken ti -- har det gjort en forskel for din søvn?}
  \item \textit{Hvad synes du generelt er det vigtigste for at leve et sundt liv i en travl hverdag i Danmark?}
\end{enumerate}

\begin{grammatikbox}{Mindmap 2B（已填关键词卡）：Mine digitale vaner}
\small
\renewcommand{\arraystretch}{1.35}
\begin{tabularx}{\linewidth}{@{} >{\leavevmode\bfseries\color{NCDNavy}\raggedright\arraybackslash}p{3.1cm} @{\hspace{0.2cm}} >{\raggedright\arraybackslash}X @{\hspace{0.35cm}} >{\leavevmode\bfseries\color{NCDNavy}\raggedright\arraybackslash}p{3.1cm} @{\hspace{0.2cm}} >{\raggedright\arraybackslash}X @{}}
\toprule
\multicolumn{4}{c}{\large\bfseries\color{NCDRed} Emne: Mine digitale vaner \quad\small\color{NCDGray}（仅含 stikord 关键词，不含整句）} \\
\midrule
1. skærmtid / telefon & ca. 2,5 timer dagligt · tidsgrænse på apps · telefon i tasken under aftensmad &
2. sociale medier & LinkedIn og lidt Instagram · slået notifikationer fra · sjældent opslag selv \\
\midrule
3. nyheder / podcasts & DR Nyheder om morgenen · danske podcasts på cyklen · øve dansk udtale &
4. nethandel / MitID & MitID til netbank, e-Boks og borger.dk · undgå impulskøb · lokale butikker \\
\midrule
5. arbejde / studie online & onlinemøder i Teams · digitale lærebøger og ordbøger · huske korte skærmpauser &
6. mening om digitale vaner (kan godt lide / ikke lide) & praktisk og tidsbesparende · nem kontakt · bekymret for uro · ønsker flere skærmfrie aftener \\
\bottomrule
\end{tabularx}
\end{grammatikbox}

\textbf{Del 2 考官追问模拟（Mine digitale vaner · 4 题）：}
\begin{enumerate}
  \item \textit{Du fortæller, at du har slået notifikationer fra -- hvordan har det påvirket din hverdag på arbejde?}
  \item \textit{Hvad synes du om, at næsten alting i Danmark foregår digitalt med MitID og e-Boks?}
  \item \textit{Vil du fortælle lidt mere om, hvordan du bruger danske podcasts eller apps til at lære dansk?}
  \item \textit{Tror du, vi bruger telefonen for meget i dag, eller er der mest fordele ved digitale løsninger?}
\end{enumerate}

% ---------------------------------------------------------------------
\subsubsection*{Mundtlig Opgave 2 (Paropgave Del 1 ca. 4 min. + Individuel Del 2 ca. 3 min. pr. kursist)}

\begin{center}
\large\bfseries\color{NCDTeal} OPGAVEARK \qquad|\qquad EMNE: Fritid og motion
\end{center}

\small
\renewcommand{\arraystretch}{1.35}
\begin{tabularx}{\linewidth}{@{} >{\raggedright\arraybackslash}X @{\hspace{0.4cm}} >{\raggedright\arraybackslash}X @{}}
\toprule
\badge{NCDNavy}{Billede 1}\ \textbf{\color{NCDNavy}[图片 1: 健身房器械训练]}\par\vspace{3pt}
\textit{Træning i et moderne fitnesscenter med vægte, løbebånd og maskiner indendørs.}
&
\badge{NCDNavy}{Billede 2}\ \textbf{\color{NCDNavy}[图片 2: 森林户外跑步]}\par\vspace{3pt}
\textit{Løbetur alene eller med en ven på en grøn skovsti i frisk luft.}
\\
\midrule
\badge{NCDNavy}{Billede 3}\ \textbf{\color{NCDNavy}[图片 3: 羽毛球/排球等团体球类]}\par\vspace{3pt}
\textit{Holdsport i en lokal sportshal, hvor man spiller badminton eller volleyball sammen.}
&
\badge{NCDNavy}{Billede 4}\ \textbf{\color{NCDNavy}[图片 4: 居家瑜伽与拉伸]}\par\vspace{3pt}
\textit{Rolig yoga og udstrækning på en træningsmåtte i stuen derhjemme.}
\\
\bottomrule
\end{tabularx}
\normalsize

\vspace{0.35em}
\begin{tipbox}{TESTERARK｜考官官方引导语与 Del 1 / Del 2 考核流程}
\small
\textbf{Del 1 (Paropgave -- ca. 4 min.)}\\
Opgavearket lægges foran kursisterne, der præsenteres for emnet på følgende måde:\\[3pt]
\textit{Emnet er 'Fritid og motion'. Her er nogle billeder, der viser forskellige måder at dyrke motion på i sin fritid. I skal tale med hinanden om, hvad I synes om de forskellige måder at dyrke motion på, og hvorfor. I skal stille spørgsmål til hinanden og svare på hinandens spørgsmål.}\\[4pt]
Hvis kursisterne ikke udtrykker deres præferencer i forhold til alle fire billeder, må tester sørge for at dette sker ved at spørge:\\[2pt]
\textit{Kan I lide at træne i et fitnesscenter / løbe ude i naturen / dyrke holdsport i en hal / lave yoga og udstrækning derhjemme? Hvorfor? / Hvorfor ikke?}\\[6pt]
\textbf{Del 2 (Individuel opgave -- ca. 3 min. pr. kursist)}\\
Tester stiller hver kursist nedenstående spørgsmål:
\begin{enumerate}[leftmargin=1.6em, itemsep=0.15em]
  \item \textit{Hvordan kan du bedst lide at slappe af eller dyrke motion i din fritid? Hvorfor?}
  \item \textit{Hvornår og hvor ofte har du tid til fritidsinteresser i en almindelig uge?}
  \item \textit{Kan du fortælle lidt om en fritidsaktivitet eller sport, du har prøvet eller gerne vil prøve? (hvor, med hvem, hvorfor, nemt/svært)}
\end{enumerate}
Hvis kursisten har besvaret et spørgsmål under samtalen i del 1, stilles uddybende spørgsmål, fx: \textit{Du siger, at \dots{} vil du fortælle lidt mere om \dots{}.} Der stilles evt. opfølgende og uddybende spørgsmål inden for tidsrammen.
\end{tipbox}

% =====================================================================
% DEL 4: RETTEARK, DETALJERET ANALYSE & HØJSCORE-MODELLER
% =====================================================================
\ncdsubsec{Del 4: Retteark, Detaljeret Analyse \& Højscore-Modeller}

\subsubsection*{4.1 官方 Retteark 评分卡（Læsning 24 point）}

\small
\renewcommand{\arraystretch}{1.28}
\begin{longtable}{@{} >{\leavevmode\bfseries\color{NCDNavy}\raggedright\arraybackslash}p{2.9cm} @{\hspace{0.25cm}} >{\raggedright\arraybackslash}p{9.6cm} @{\hspace{0.25cm}} >{\centering\arraybackslash}p{1.75cm} @{\hspace{0.25cm}} >{\centering\arraybackslash}p{1.65cm} @{}}
\toprule
\textbf{模块 / 题号} & \textbf{Rettenøgle（标准答案）} & \textbf{分值} & \textbf{自测得分} \\
\midrule
\endfirsthead
\toprule
\textbf{模块 / 题号（续表）} & \textbf{Rettenøgle（标准答案）} & \textbf{分值} & \textbf{自测得分} \\
\midrule
\endhead
\bottomrule
\endfoot
\textbf{Læsning Opgave 1}\newline（6 题） &
1. \textbf{før} \quad 2. \textbf{fordi} \quad 3. \textbf{dog} \quad 4. \textbf{til} \quad 5. \textbf{at} \quad 6. \textbf{givet}\newline
\textcolor{NCDGray}{\small （多余干扰词：blev, eller, fra, med, men, som｜例词 0 为 det）} & 6 point & \rule{0.75cm}{0.4pt} / 6 \\
\midrule
\textbf{Læsning Opgave 2A}\newline（3 题 · 聊天补全） &
1. \textbf{B} \quad 2. \textbf{E} \quad 3. \textbf{A}\newline
\textcolor{NCDGray}{\small （多余干扰项：C, D, F｜例句 0 为 Z）} & 3 point & \rule{0.75cm}{0.4pt} / 3 \\
\midrule
\textbf{Læsning Opgave 2B}\newline（3 题 · 聊天补全） &
1. \textbf{D} \quad 2. \textbf{F} \quad 3. \textbf{C}\newline
\textcolor{NCDGray}{\small （多余干扰项：A, B, E｜例句 0 为 Z）} & 3 point & \rule{0.75cm}{0.4pt} / 3 \\
\midrule
\textbf{Læsning Opgave 3}\newline（5 题 · 段内插句） &
1. \textbf{B} \quad 2. \textbf{A} \quad 3. \textbf{D} \quad 4. \textbf{B} \quad 5. \textbf{C} & 5 point & \rule{0.75cm}{0.4pt} / 5 \\
\midrule
\textbf{Læsning Opgave 4}\newline（7 题 · 三人匹配） &
1. \textbf{A} \quad 2. \textbf{C} \quad 3. \textbf{B} \quad 4. \textbf{A} \quad 5. \textbf{C} \quad 6. \textbf{B} \quad 7. \textbf{A}\newline
\textcolor{NCDGray}{\small （例题 0 为 B. Tobias）} & 7 point & \rule{0.75cm}{0.4pt} / 7 \\
\midrule
\textbf{Læsning 总计} & \textbf{Bestået：15--24 point} \quad|\quad \textbf{Ikke bestået：0--14 point} & \textbf{24 point} & \rule{0.75cm}{0.4pt} / 24 \\
\end{longtable}
\normalsize

\begin{tipbox}{Skrivning \& Mundtlig 自检过关准则（Helhedsvurdering）}
\small
\begin{itemize}
  \item \textbf{Skrivning（整体评估 Bestået）}：\ding{172} Opgave 1 投诉信必须不漏任何一条指定要点，且全程用 \textbf{\textit{I / jer / jeres}} 称呼酒店（切忌中途跳成 \textit{du / dig}）；\ding{173} Opgave 2 私人邮件严格满 \textbf{90 词以上}（建议写到 105--135 词），回答 Astrid 全部 4 问，并含得体称呼与落款；\ding{174} 主动运用 NCD 3 核心语法（\textit{-s/blive} 被动、\textit{til trods for at / på baggrund af}、从句中否定副词前置等）。
  \item \textbf{Mundtlig Opgave 1 \& 2（整体评估 Bestået）}：Opgave 1 陈述紧扣 6 个 \textit{stikord} 并简要说明原因；Opgave 2 双人对话必须对\textbf{全部 4 张运动图}明确表达喜欢/不喜欢并给出理由，同时积极向搭档发问。
\end{itemize}
\end{tipbox}

% ---------------------------------------------------------------------
\subsubsection*{4.2 阅读 Opgave 1--4 逐题解析、全文精译与核心词表}

\paragraph*{一、Opgave 1（一周无社交媒体实验）词汇语法定位与干扰项排除}
\begin{itemize}
  \item \textbf{(0) det [Eksempel]}：回指前分句“每天花近三小时刷社交媒体”这件事，作 \textit{begyndte} 的形式/回指主语。
  \item \textbf{1. før}：时间从句连词“在……之前”，与 \textit{overhovedet stod ud af sengen} 搭配（“甚至还没下床就立刻查看消息”）。\textit{men / eller / som} 无法引导该时间状语从句，\textit{fra} 与 \textit{overhovedet} 语义不合。
  \item \textbf{2. fordi}：引导原因状语从句；注意从句内部状语 \textit{hele tiden} 位于变位动词 \textit{lyste} \textbf{之前}（典型从句语序），由此直接排除并列连词 \textit{men / eller}（并列句不改变常规主句语序）及介词/关系代词 \textit{fra / med / som}。
  \item \textbf{3. dog}：句中让步对比副词“然而/不过”，承接上文“前两天出奇难熬”与后文“三天后一切\textbf{却}平静得多”。该位置位于主句动词 \textit{blev} 与表语 \textit{meget roligere} 之间，其余候选词均不符合词性或句法位置。
  \item \textbf{4. til}：高频固定搭配 \textit{tid til at + 不定式}（“有时间去做某事”），不可替换为 \textit{med / fra}。
  \item \textbf{5. at}：引导 \textit{mærkede også tydeligt} 的宾语从句（“明显感觉到……”），只有连接词 \textit{at} 合法。
  \item \textbf{6. givet}：跟在助动词 \textit{har} 后构成现在完成时，且后面紧接双宾语 \textit{ham mere energi}（“给了他更多精力”）；其余干扰词 \textit{blev} 不能接在 \textit{har} 后，也无法带双宾语。
\end{itemize}

\begin{transbox}
\small
\textbf{\color{NCDTeal}[Opgave 1 全文中文精译]}\\
米克尔今年二十九岁，在罗斯基勒担任平面设计师。多年来，他每天几乎花三个小时在社交媒体上，\textbf{这件事}慢慢开始让他感到烦躁。每天早上闹钟一响，他甚至\textbf{在}还没下床\textbf{之前}就立刻查看消息。在工作时他也常常难以集中注意力，\textbf{因为}屏幕总是不断亮起新通知和短视频。于是上个月他决定做一个实验：把手机上所有社交软件删除七天，感受一下变化。最初两天出奇地难熬。每当米克尔坐在火车上或在食堂等咖啡时，手几乎会条件反射般地伸进口袋掏手机。\textbf{不过}三天之后，一切都变得平静多了。少了源源不断的图片信息流，他晚上睡得更香，还突然有了时间\textbf{去}读一本在床头柜躺了半年的小说。他也明显感觉到，\textbf{（就是）}自己在休息时间与同事聊得更多了，而不是只低头盯着屏幕。等这周终于结束时，米克尔只重新下载了两个应用，并且关掉了所有推送通知。现在他每天最多只花二十分钟上网，而这\textbf{给}他在日常生活中带来了更多精力。
\end{transbox}

\paragraph*{二、Opgave 2A \& 2B（聊天对话补全）双向衔接解析与精译}
\begin{itemize}
  \item \textbf{Opgave 2A（Lukas 与 Emma 合买周年礼物）}：
    \begin{itemize}
      \item \textbf{1. $\rightarrow$ B}：Lukas 2 询问每人分摊多少（\textit{Hvor meget... hver person}），B 提出“每人 100 克朗”且顺带引出“还需要买配咖啡的点心（\textit{noget lækkert til kaffen}）”，正好触发 Lukas 3 的“8 人共 800 克朗”及“要不要我顺便去面包店订蛋糕”。干扰项 D 虽出现 \textit{100 kroner hver}，却否定收钱，与 Lukas 3 的 \textit{i alt 800 kroner} 直接矛盾。
      \item \textbf{2. $\rightarrow$ E}：Lukas 3 主动问能否去订蛋糕，E 欣然同意并请他预订“无坚果巧克力蛋糕（\textit{chokoladekage uden nødder}）”，紧接 Lukas 4 的“好的，那我明天上午打电话订无坚果巧克力蛋糕”。干扰项 C 说“根本不需要蛋糕”，与 Lukas 4 矛盾。
      \item \textbf{3. $\rightarrow$ A}：Lukas 4 问两人是否知道周六聚会，A 回答“他们毫不知情，哥哥 14:30 悄悄放我们进去”，完美引出 Lukas 5 的“太好了，那我当然保密，周六 14:30 见”。干扰项 F 说“已发邀请、15 点见”，与保密及时间双双冲突。
    \end{itemize}
  \item \textbf{Opgave 2B（Sofie 与跑步教练 Peter 沟通膝伤）}：
    \begin{itemize}
      \item \textbf{1. $\rightarrow$ D}：Sofie 2 问需要停跑多久，D 建议“停跑一周直到膝盖完全不痛”，正对 Sofie 3 的“一整周不跑听起来有点久（\textit{En hel uge uden løb}）”。干扰项 A 劝她今晚照常间歇跑，与例句 Z 矛盾。
      \item \textbf{2. $\rightarrow$ F}：Sofie 3 问这周能不能改练别的项目，F 建议“游泳和居家轻柔拉伸（\textit{svømme ... rolige strækøvelser}）”，正对 Sofie 4 开头的“游泳和拉伸听起来确实是好计划”。干扰项 B 说“七天一动不动并避开水和拉伸”，与 Sofie 4 冲突。
      \item \textbf{3. $\rightarrow$ C}：Sofie 4 问复跑前有什么注意事项，C 叮嘱“在平坦路面慢跑起步，并检查跑鞋是否太旧该换了”，正好引出 Sofie 5 的“那我周五就去买双新鞋，下周慢慢开始”。干扰项 E 建议“明早快跑 10 公里测试”，与休息一周和慢慢重启双双矛盾。
    \end{itemize}
\end{itemize}

\begin{transbox}
\small
\textbf{\color{NCDTeal}[Opgave 2A \& 2B 完整对话中文精译]}\\
\textbf{2A（Lukas \& Emma）}：L1：嗨艾玛！西格妮和卡斯珀下周六就是五周年纪念日了，你想好大家合送什么礼物了吗？ $\rightarrow$ \textbf{Z（例）}：嗨！还没完全定，不过我在想可以送一家好餐厅的礼品卡，你觉得呢？ $\rightarrow$ L2：好主意！他们俩都爱吃意餐。你觉得每人出多少钱合适？ $\rightarrow$ \textbf{1. B}：每人 100 克朗怎么样？谁都不会有压力，而且配咖啡我们也得准备点甜点。 $\rightarrow$ L3：很公平，那 8 个人总共就是 800 克朗。蛋糕也要我去面包店订吗？ $\rightarrow$ \textbf{2. E}：对，麻烦你了，因为你就住在面包店旁边。你能订一个不含坚果的巧克力蛋糕吗？ $\rightarrow$ L4：没问题，我明天上午就打电话订无坚果巧克力蛋糕。惊喜安排呢——他们知道我们周六下午过去吗？ $\rightarrow$ \textbf{3. A}：不知道，他们毫不知情！卡斯珀的弟弟会在 14:30 趁他们出门买东西时悄悄放我们进去。 $\rightarrow$ L5：太棒了，我肯定严格保密，这就给其他客人发集合时间。周六 14:30 见！\\[4pt]
\textbf{2B（Sofie \& Peter）}：S1：嗨彼得！周日长跑后我右膝盖有点疼，你觉得我今晚能参加间歇跑吗？ $\rightarrow$ \textbf{Z（例）}：嗨索菲！真遗憾听到这个消息。如果有疼痛，今晚绝对要暂停高强度跑步。 $\rightarrow$ S2：好吧，你说得对，虽然错过训练有点郁闷。我要等多久才能重新开始跑？ $\rightarrow$ \textbf{1. D}：我建议你暂停跑步一周，直到膝盖彻底不疼为止。 $\rightarrow$ S3：一周不跑听起来有点久，但我也不想落下重伤。这周我是一点都不能动，还是可以练点别的？ $\rightarrow$ \textbf{2. F}：可以保持身体活动。试着去游几次泳并在家做轻柔拉伸，因为这不会给膝盖施压。 $\rightarrow$ S4：游泳和拉伸听起来确实不错，而且我公司旁边刚好有游泳馆。试着短跑前有什么要特别记住的吗？ $\rightarrow$ \textbf{3. C}：一定在平坦路面上极慢地起步，而且记得检查跑鞋是不是太旧该换新了。 $\rightarrow$ S5：谢谢你的好建议，彼得！那我周五就去买双新鞋，下周轻轻恢复。训练场很快见！
\end{transbox}

\paragraph*{三、Opgave 3（Amira 加入美食俱乐部）段内插句精析与全文精译}
\begin{itemize}
  \item \textbf{第 1 段 $\rightarrow$ B}：前句写她刚搬来人生地不熟、晚饭总靠超市快餐将就；后句写她在图书馆看到“邻里做饭俱乐部”海报。选项 B“她既怀念亲手从头做热饭，也渴望认识新朋友”承前启后，直接解释了她的动机；A（更喜欢独自在家吃）、C（老邻居众多）、D（没时间去图书馆）均与上下文自相矛盾。
  \item \textbf{第 2 段 $\rightarrow$ A}：前句写她第一次赴约有点紧张（\textit{en smule nervøs}），担心融不进去；后句却画面一转“短短十分钟后她就和长者延斯在炉灶边有说有笑地切菜”。A“幸好她受到热情欢迎并立刻分到厨房任务”精准完成转折衔接；B（别人早吃完了）、C（立刻转身回家）、D（没人愿意和新成员说话）皆违背下文事实。
  \item \textbf{第 3 段 $\rightarrow$ D}：前句说轮到阿米拉定菜单，她想展示家乡传统菜；后句说她有点担心香料对丹麦人太辣，但炖菜很快被抢光。D“因此她翻出一份古老的家族食谱，花了一整个下午准备香料”将“传统菜”和后文的 \textit{krydderierne / gryderetten} 紧密扣合；A（点披萨）、B（大家只吃丹麦菜）、C（做别国菜）均断链。
  \item \textbf{第 4 段 $\rightarrow$ B}：前句指大家围坐聊天帮她练习了丹麦语；后句写“通过这种方式，她现在在实验室与同事闲聊讨论更有自信”。B“此外，她学到许多课本里找不到的日常地道表达”恰当延伸口语收获并顺接 \textit{På den måde}；其余选项（只说英语、退掉语言课、嫌聊天无聊）均错误。
  \item \textbf{第 5 段 $\rightarrow$ C}：前句说尽管工作累她也很少请假缺席（\textit{melder sjældent afbud}）；后句说她已经在和延斯策划下个月菜单。C“她感觉到共同晚餐能给剩下的一周注入新能量”以 \textit{nemlig} 解释原因并自然过渡；A（想退社）、B（累了就宅在沙发）、D（俱乐部圣诞后关闭）均与前后句冲突。
\end{itemize}

\begin{transbox}
\small
\textbf{\color{NCDTeal}[Opgave 3 全文中文精译]}\\
\textbf{1.} 阿米拉最近搬到欧登塞，开始担任生物分析员的新工作。起初她在城里几乎谁也不认识，下班后经常独自坐在公寓里对付一份超市快餐。\textbf{[B] 她既怀念亲手从头做一顿正餐，也渴望结识新朋友。}有一天，她在当地图书馆看到一张社区美食俱乐部的海报，邻里们每两周聚在一起做一次饭。\textbf{2.} 俱乐部由住在同一街区、年龄各异的八个人组成。阿米拉第一次去聚会时有点紧张，因为不知道自己能否融入这个集体。\textbf{[A] 幸好大家非常热情地欢迎了她，并立刻给她安排了厨房里的活儿。}短短十分钟后，她就已经站在炉灶旁，和一位名叫延斯的长者一边切汤菜一边开心大笑了。\textbf{3.} 在美食俱乐部里，成员们轮流决定菜单并采购食材。十一月轮到阿米拉做东时，她很想向大家展示一道家乡传统菜。\textbf{[D] 因此她找出一份古老的家族食谱，花了整个下午准备香料。}阿米拉原本有点担心香料对丹麦成员来说太辣，但那锅炖菜很快一扫而空，好几位成员立刻找她要食谱。\textbf{4.} 除了美食本身，俱乐部的夜晚对阿米拉还有另一重重要意义。当八位成员围坐在餐桌旁聊工作、新闻与日常生活时，她在轻松氛围中锻炼了丹麦语。\textbf{[B] 此外，她还学到了许多教材里没有的日常表达。}通过这种方式，她现在在实验室里与同事讨论或闲聊时感觉自信多了。\textbf{5.} 如今阿米拉加入俱乐部快一年了，周二晚上成了她日历上最喜欢的一天。尽管工作依然很忙、有时下班很累，她却很少请假缺席。\textbf{[C] 因为她感受到，一起共进晚餐能给她剩下的一周重新注入活力。}事实上，她已经在和延斯一起策划下个月的菜单了。
\end{transbox}

\paragraph*{四、Opgave 4（三人谈手机与屏幕时间）同义转述定位与精译}
\begin{itemize}
  \item \textbf{(0) B. Tobias [Eksempel]}：\textit{ser næsten ikke på en skærm i løbet af sin arbejdsdag} $\leftrightarrow$ \textit{kigger næsten aldrig på en skærm i arbejdstiden}（木匠工作几乎不看屏幕；Ida 全天用电脑手机回消息，Lin 白天用平板读书）。
  \item \textbf{1. A. Ida}：\textit{lader ikke sin telefon ligge i soveværelset om natten længere} $\leftrightarrow$ \textit{min mobil bliver liggende ude i gangen hele natten}（过去曾带手机上床，如今手机整晚留在走廊并换用老式闹钟；Lin 睡前仍在卧室刷消息）。
  \item \textbf{2. C. Lin}：\textit{øve sig i et fremmedsprog} $\leftrightarrow$ \textit{øve dansk ... lytter til danske nyhedspodcasts og træner nye ord i en sprogapp}。
  \item \textbf{3. B. Tobias}：\textit{fordi et familiemedlem sagde noget om det} $\leftrightarrow$ \textit{min tiårige søn, der fik mig til at tænke over det, da han en dag spurgte, hvorfor jeg altid kiggede ned i mobilen}（注意不要被 Ida 自己因失眠立规矩干扰）。
  \item \textbf{4. A. Ida}：\textit{aftale med sine venner om at lægge telefonerne væk} $\leftrightarrow$ \textit{mine veninder og jeg indført en fast aftale ... Alle telefoner bliver lagt midt på bordet på lydløs}。
  \item \textbf{5. C. Lin}：\textit{bevare kontakten til folk langt væk} $\leftrightarrow$ \textit{ringer til mine forældre på video hver søndag ... føler jeg mig tæt på dem på trods af den store afstand}。
  \item \textbf{6. B. Tobias}：\textit{sætter tid af på én bestemt ugedag til helt at slukke for mobilen} $\leftrightarrow$ \textit{indført en fast mobilfri søndag, hvor min telefon er helt slukket fra morgen til aften}（Ida 只是每晚不看工作邮件并放走廊，Lin 每周日反而要跟父母视频）。
  \item \textbf{7. A. Ida}：\textit{sover bedre nu, efter at personen ændrede sine aftenregler} $\leftrightarrow$ \textit{falder meget hurtigere i søvn nu og vågner faktisk frisk og udhvilet}（Tobias 仅谈及陪儿子玩桌游，Lin 从未调整晚间习惯）。
\end{itemize}

\begin{transbox}
\small
\textbf{\color{NCDTeal}[Opgave 4 三人自述中文精译]}\\
\textbf{A. 艾达}：我在奥胡斯一家中型公司做项目协调员，整天都要在电脑和手机上接收各种消息。以前我经常把工作带到床上，因为睡前总想再看一眼最后的邮件。这彻底毁了我的夜间睡眠，因为大脑整晚都在继续琢磨工作任务。因此半年前我制定了固定的晚间规定：晚上七点后绝不回复工作邮件，手机整晚放在门厅走廊里，早上靠一个老式闹钟叫醒我。现在我入睡快得多，醒来也真正感到神清气爽。此外，我和闺蜜们下班后约咖啡时也立下规矩：所有手机静音放在桌子正中央直到散场回家，这样大家才算真正专心相伴。\\
\textbf{B. 托比亚斯}：作为瓦埃勒的木匠，我白天大部分时间都在户外或工坊干活，所以工作时间几乎从不看屏幕。可另一方面，去年秋天我却发现自己每晚窝在沙发上刷体育视频、看网上新闻超过两个小时。其实是我十岁的儿子点醒了我——有一天他问我，为什么每次他想和我说话时，我总是在低头看手机。那句话深深戳中了我。现在每天傍晚五点到八点，我都会把手机放在厨房高处的架子上，好让我们安安静静地一起吃晚饭和玩桌游。此外我还推行了“无手机星期日”，周日从早到晚彻底关机，全家人一起骑车或钓鱼。\\
\textbf{C. 林}：我是欧登塞大学的医学生，一昼夜里使用电子设备的时间非常长。我的所有教材和讲义笔记都在平板电脑上，所以在日常生活中我根本离不开屏幕。与此同时，每天早上坐长途公交去校园时，我也会积极用手机练习丹麦语：我会听丹麦语新闻播客，并在语言软件里背新词。很多人说手机只会拉远人与人的距离，但对我来说恰恰相反。我的直系亲属住在新加坡，正因为我每周日都能跟父母视频通话，日常也能发些丹麦生活照，尽管相隔万里，我依然觉得与他们很亲近。由于时差，我睡前仍会在卧室看一眼消息，而这并不会影响我第二天早上的学业。
\end{transbox}

\paragraph*{五、核心高频词表（Gloser｜15 个必背屈折变化词条）}

\begin{glosertable}
名词 & en skærmtid & \textit{(skærmtiden, ingen flertal)} & 屏幕使用时间 \\
名词 & et eksperiment & \textit{(eksperimentet, eksperimenter, eksperimenterne)} & 实验，尝试 \\
动词 & at irritere & \textit{(irriterer, irriterede, har irriteret)} & 使烦躁，惹恼 \\
动词 & at stirre & \textit{(stirrer, stirrede, har stirret)} & 盯着看，凝视 \\
名词 & et jubilæum & \textit{(jubilæet, jubilæer, jubilæerne)} & 周年纪念日 \\
动词 & at anbefale & \textit{(anbefaler, anbefalede, har anbefalet)} & 推荐，建议 \\
动词 & at belaste & \textit{(belaster, belastede, har belastet)} & 给……施加负荷/压力 \\
短语 & at melde afbud & \textit{(melder afbud, meldte afbud, har meldt afbud)} & 取消预约，请假缺席 \\
动词 & at undvære & \textit{(undværer, undværede, har undværet)} & 缺少，离不开（\textit{ikke kunne undvære}） \\
副词 & til gengæld & （副词短语，句首触发 V2 倒装） & 相反，另一方面 \\
名词 & et ophold & \textit{(opholdet, ophold, opholdene)} & 逗留，住宿期间 \\
形容词 & uhøflig & \textit{(uhøflig, uhøfligt, uhøflige)} & 无礼的，态度恶劣的 \\
动词 & at refundere & \textit{(refunderer, refunderede, har refunderet)} & 退款，退还费用 \\
名词 & en balance & \textit{(balancen, balancer, balancerne)} & 平衡，均衡 \\
短语 & at koble af & \textit{(kobler af, koblede af, har koblet af)} & 彻底放松，摆脱忙碌 \\
\end{glosertable}

% =====================================================================
% 4.3 写作双档高分示范与点评
% =====================================================================
\subsubsection*{4.3 写作双档高分示范与点评（Skrivning Modeller）}

\begin{emailbox}{Opgave 1 范文一：B1 稳过版（115 词｜全覆 4 问 · 称谓一致 · 因果清晰）}
Kære Hotel Kystblik

Jeg skriver for at klage over mit ophold på jeres hotel fra den 14. til den 16. november, hvor jeg boede på værelse 204.

Desværre opstod der flere problemer. For det første var der kraftig støj fra gaden hele natten, så jeg næsten ikke sov. Derudover var vandet i bruseren helt koldt lørdag morgen, og værelset var ikke gjort ordentligt rent.

Da jeg gik ned i receptionen for at fortælle om dette, var receptionisten meget uhøflig. Han sagde bare kort, at han ikke kunne hjælpe, og gav ingen undskyldning.

Derfor forventer jeg, at I betaler halvdelen af prisen tilbage som kompensation. Jeg hører gerne fra jer hurtigst muligt.

Venlig hilsen\\
Katrine Holm

\vspace{0.3em}
\rule{\linewidth}{0.4pt}\\[0.2em]
\small\textbf{\color{NCDTeal}[考官点评 · 115 ord]}：稳健达标并覆盖全部 4 项要点（入住日期与房号、三大问题、前台反应、明确半价赔偿要求）。全文自始至终保持 \textit{jeres / I / jer} 第二人称复数指代酒店，结构连接词（\textit{For det første / Derudover / Da / Derfor}）运用地道，从句否定语序正确（\textit{så jeg næsten ikke sov}、\textit{at han ikke kunne hjælpe}）。
\end{emailbox}

\begin{emailbox}{Opgave 1 范文二：B1+ 高分示范版（136 词｜NCD 3 被动 + 复合介词 + 整句关系从句）}
Kære Hotel Kystblik

Jeg skriver denne klage \da{på baggrund af} et meget skuffende weekendophold på jeres hotel fra den 14. til den 16. november, hvor jeg havde bestilt værelse 204.

\da{Til trods for at} hotellet \dabold{beskrives} som roligt, opstod der flere alvorlige problemer under opholdet. For det første holdt høj musik fra gaden mig vågen begge nætter, fordi de gamle vinduer ikke lukkede tæt. Desuden var vandet i bruseren iskoldt lørdag morgen, og rengøringen \dabold{var ikke blevet udført} ordentligt.

Da jeg henvendte mig i receptionen, reagerede medarbejderen desværre uhøfligt. I stedet for at beklage svarede han surt, at støj var normalt i weekenden, \da{hvilket gjorde mig} meget skuffet.

Da jeg ikke fik den service, der \dabold{blev lovet} ved bestillingen, forventer jeg, at I refunderer halvtreds procent af beløbet inden fjorten dage.

Med venlig hilsen\\
Katrine Holm

\vspace{0.3em}
\rule{\linewidth}{0.4pt}\\[0.2em]
\small\textbf{\color{NCDNavy}[高分亮点 · 136 ord]}：\ding{172} 精准调用 NCD 3 书面复合介词 \textit{på baggrund af} 与让步结构 \textit{Til trods for at}；\ding{173} 灵活用出两种被动表达（现状描述用 \textit{-s} 被动 \textit{beskrives}，过去完成事件与动作分别用 \textit{var ikke blevet udført} 和 \textit{blev lovet}）；\ding{174} 末尾用整句关系从句 \textit{..., hvilket gjorde mig meget skuffet} 点评态度，措辞礼貌而坚定。
\end{emailbox}

\begin{emailbox}{Opgave 2 范文一：B1 稳过版（111 词｜超过 90 词底线 · 四问全答 · 自然亲切）}
Hej Astrid

Tak for din søde mail! Det er dejligt at høre fra dig igen.

Ja, jeg er lige begyndt på et aftenkursus i projektledelse hver tirsdag aften. Jeg valgte det, fordi jeg gerne vil lære nye ting og få mere ansvar på mit arbejde næste år.

Det kan godt være travlt at nå både fuldtidsarbejde, danskundervisning og fritid i hverdagen. Derfor laver jeg altid en ugeplan om søndagen og forbereder aftensmad til flere dage ad gangen.

Når det bliver weekend, vil jeg helst slappe af. Så slukker jeg arbejdsmailen, går en lang tur i skoven og mødes med venner til kaffe. Hvad laver du selv for tiden?

Mange hilsner\\
Rasmus

\vspace{0.3em}
\rule{\linewidth}{0.4pt}\\[0.2em]
\small\textbf{\color{NCDTeal}[考官点评 · 111 ord]}：稳超官方 90 词最低门槛，行文流畅自然。依次应答了 Astrid 全部问题（参加周二项目管理夜校、明年职场晋升动机、周日列周计划与提前备餐的时间管理法、周末关邮箱逛森林与喝咖啡放松），并在结尾反问互动，语气完全契合挚友通信语域。
\end{emailbox}

\begin{emailbox}{Opgave 2 范文二：B1+ 高分示范版（137 词｜条件倒装/方式状语 + 形式宾语 + 从句切换）}
Kære Astrid

Hvor er det dejligt at høre fra dig! Det passer, at jeg for nylig startede på et spændende aftenkursus i digital kommunikation. Jeg valgte holdet, fordi jeg ønsker at styrke mine faglige kompetencer på det danske arbejdsmarked, \da{selvom mine hverdage allerede er} ret fyldte.

Til tider kræver det god planlægning at få fuldtidsjob, sprogskole to aftener om ugen og fritid til at hænge sammen. Men \da{ved at skrive} alle aftaler i kalenderen og lægge mobilen væk efter klokken ni \dabold{når jeg} stadig både lektier, madlavning og søvn, selvom det \dabold{ikke altid er} nemt.

I weekenden lægger jeg derfor vægt \da{på, at jeg får koblet} helt af fra arbejdet. Jeg bager frisk brød lørdag morgen, laver rolig yoga derhjemme og bruger søndagen i naturen. Skal vi ikke snart mødes til en kop kaffe?

Kærlig hilsen\\
Rasmus

\vspace{0.3em}
\rule{\linewidth}{0.4pt}\\[0.2em]
\small\textbf{\color{NCDNavy}[高分亮点 · 137 ord]}：\ding{172} 方式状语短语 \textit{ved at skrive ... og lægge ...} 置句首直接触发主句 V2 倒装 \textit{når jeg}；\ding{173} 准确驾驭“介词 + 从句”结构（\textit{lægger vægt på, at...}）及让步从句中副词前置（\textit{selvom det ikke altid er nemt}）；\ding{174} 词汇丰富地道（\textit{faglige kompetencer, hænge sammen, koble af}），自然冲击 B1 上限高分。
\end{emailbox}

% =====================================================================
% 4.4 口语示范逐字稿、核心 IPA 与双人四图互动全真逐字稿
% =====================================================================
\subsubsection*{4.4 口语示范逐字稿、核心 IPA 与 Opgave 2 双人四图互动逐字稿}

\begin{transbox}
\small
\textbf{\color{NCDNavy}[Mindmap 2A 口语陈述范文：Mine sunde og usunde vaner｜150 词 · 约 1.5 分钟]}\\
\textit{Jeg vil gerne fortælle om mine sunde og usunde vaner \ipa{ˈvæːnɐ}, som spiller en stor rolle i min hverdag. På den sunde side starter jeg næsten altid dagen med havregryn \ipa{ˈhɑu̯ʁɐˌɡʁyːˀn} og frisk frugt, og til aftensmad spiser vi mange grøntsager \ipa{ˈɡʁœnˌsæːɐ}, fuldkornsbrød og fisk to gange om ugen. Jeg drikker også masser af vand i stedet for sodavand. Alligevel har jeg nogle usunde vaner, for når jeg er træt om aftenen, spiser jeg tit lidt for meget chokolade \ipa{ɕoɡoˈlæːðə} eller fredagsslik, og på kontoret drikker jeg nogle gange fire kopper kaffe. Med hensyn til motion \ipa{moˈɕoːˀn} cykler jeg otte kilometer til arbejde hver dag, svømmer \ipa{ˈsvømɐ} om onsdagen og går tur i skoven om søndagen. Jeg har aldrig røget, og jeg drikker kun et glas vin i weekenden. Til gengæld burde jeg forbedre min søvn \ipa{ˈsœu̯ˀn} ved at lægge mobilen væk klokken ti. Alt i alt er jeg tilfreds med min livsstil \ipa{ˈliusˌsdiːˀl}, fordi balancen fungerer godt.}
\end{transbox}

\paragraph*{Mindmap 2A Del 2 考官追问实战示范对话：}
\dialogturn{1}{Eksaminator}{Du siger, at du spiser lidt for meget chokolade om aftenen -- hvornår har du mest lyst til usunde ting, og hvorfor?}{chokolade \ipa{ɕoɡoˈlæːðə} · usunde \ipa{ˈuˌsʌnə}}{你说你晚上会吃有点多的巧克力——你什么时候最想吃不健康的东西，为什么？}
\dialogturn{2}{Kursist}{Det sker næsten altid efter klokken ni om aftenen, når jeg endelig sidder stille i sofaen efter en lang arbejdsdag. Jeg tror, det handler om, at jeg er træt i hovedet og bare vil belønne mig selv lidt.}{arbejdsdag \ipa{ˈɑːbɑi̯ˀdsˌdæːˀ} · belønne \ipa{beˈlœnˀə}}{几乎总是在晚上九点以后，经历漫长工作终于安静坐在沙发上的时候。我觉得这是因为脑子累了，想稍微犒劳一下自己。}
\dialogturn{3}{Eksaminator}{Hvorfor kan du godt lide at svømme og cykle til arbejde, og hvordan holder du motivationen om vinteren?}{motivationen \ipa{motˢivaˈɕoːˀnən}}{为什么你喜欢游泳和骑车上班，冬天又是怎么保持动力的？}
\dialogturn{4}{Kursist}{Når jeg cykler, vågner jeg godt op om morgenen, og i svømmehallen slapper hele kroppen af. Om vinteren har jeg varmt regntøj og lygter på cyklen, og fordi jeg har en fast aftale med en ven i svømmehallen hver onsdag, bliver jeg aldrig hjemme.}{svømmehallen \ipa{ˈsvøməˌhalˀən} · lygter \ipa{ˈløɡdɐ}}{骑车能让我早上彻底清醒，在游泳馆全身都能放松。冬天我备好了保暖雨衣和车灯，加上每周三和朋友约好了去游泳，所以我从不缺席。}
\dialogturn{5}{Eksaminator}{Du nævner, at du lægger telefonen væk klokken ti -- har det gjort en forskel for din søvn?}{forskel \ipa{ˈfʌːˌsɡɛlˀ} · søvn \ipa{ˈsœu̯ˀn}}{你提到你十点把手机放一边——这对你的睡眠带来了改变吗？}
\dialogturn{6}{Kursist}{Ja, en ret stor forskel faktisk. Før lå jeg og læste nyheder til midnat, så det var svært at falde i søvn. Nu læser jeg tyve minutter i en almindelig bog i stedet, og derfor sover jeg meget hurtigere.}{nyheder \ipa{ˈnyˌheðˀɐ} · almindelig \ipa{alˈmenˀəli}}{是的，其实改变还挺大的。以前我会躺着刷新闻到半夜，很难入睡。现在我改看二十分钟纸质书，因此入睡快得多。}
\dialogturn{7}{Eksaminator}{Hvad synes du generelt er det vigtigste for at leve et sundt liv i en travl hverdag i Danmark?}{generelt \ipa{ɡenəˈʁɑlˀd}}{总体而言，你认为在丹麦忙碌的日常生活中过上健康生活最重要的是什么？}
\dialogturn{8}{Kursist}{Jeg synes, det vigtigste er at have en god balance og nogle enkle rutiner, for eksempel at cykle til arbejde. Hvis reglerne er alt for strenge, giver man hurtigt op igen.}{balance \ipa{baˈlɑŋsə} · strenge \ipa{ˈsdʁɛŋə}}{我觉得最重要的是保持良好平衡并养成简单习惯，比如骑车上班。如果规矩订得太苛刻，人很快就会放弃。}

\vspace{0.3em}
\begin{transbox}
\small
\textbf{\color{NCDNavy}[Mindmap 2B 口语陈述范文：Mine digitale vaner｜146 词 · 约 1.5 分钟]}\\
\textit{Nu vil jeg gerne fortælle lidt om mine digitale vaner. I gennemsnit bruger jeg cirka to og en halv time på min telefon om dagen, men jeg har sat en skærmtidsgrænse \ipa{ˈsɡɛɐ̯mˌtˢiðsˌɡʁɛnsə} på flere apps, og ved aftensmaden ligger mobilen i tasken. På sociale medier \ipa{soˈɕæːlə ˈmeːðiɐ} bruger jeg mest LinkedIn til arbejde og kigger lidt på Instagram i bussen, men jeg har slået alle notifikationer \ipa{notˢifikaˈɕoːˀnɐ} fra for ikke at blive forstyrret. Jeg bruger også telefonen aktivt til at følge med i samfundet: Om morgenen læser jeg DR Nyheder, og når jeg cykler på arbejde, lytter jeg til danske podcasts for at øve udtalen \ipa{ˈuðˌtˢæːlən}. Derudover bruger jeg MitID næsten dagligt til netbank, e-Boks og borger.dk, mens jeg undgår impulskøb \ipa{emˈpʰulˀsˌkʰøːˀb} online. På arbejdet har vi ofte møder i Teams, hvilket sparer transporttid, men kræver skærmpauser. Generelt er jeg glad for digitale løsninger, selvom jeg gerne vil have flere skærmfrie aftener.}
\end{transbox}

\paragraph*{Mindmap 2B Del 2 考官追问实战示范对话：}
\dialogturn{1}{Eksaminator}{Du fortæller, at du har slået notifikationer fra -- hvordan har det påvirket din hverdag på arbejde?}{påvirket \ipa{ˈpʰʌˌviɐ̯ɡəð}}{你说你关掉了消息通知——这对你的日常工作有什么影响？}
\dialogturn{2}{Kursist}{Det har hjulpet rigtig meget på min koncentration. Før kiggede jeg hele tiden på skærmen, hver gang telefonen vibrerede, så det tog lang tid at blive færdig med en opgave. Nu tjekker jeg kun beskeder i min frokostpause.}{koncentration \ipa{kʰʌnsəntˢʁaˈɕoːˀn}}{这对我集中注意力帮助极大。以前手机一震动我就看屏幕，一项任务拖很久才做完；现在我只在午休时看消息。}
\dialogturn{3}{Eksaminator}{Hvad synes du om, at næsten alting i Danmark foregår digitalt med MitID og e-Boks?}{foregår \ipa{ˈfɒːˌɡʌˀ}}{几乎所有事在丹麦都要用 MitID 和 e-Boks 线上办理，你对此怎么看？}
\dialogturn{4}{Kursist}{For mig er det super praktisk, fordi jeg kan ordne lægetid, bank og breve fra kommunen hjemmefra på fem minutter. Men til trods for at det er smart, synes jeg også, det kan være svært for ældre mennesker eller nyankomne udlændinge.}{praktisk \ipa{ˈpʰʁɑɡtˢisɡ} · udlændinge \ipa{ˈuðˌlɛneŋə}}{对我来说非常方便，在家五分钟就能搞定预约医生、银行和市政信件。不过尽管很智能，我也觉得这对老年人或新来的外国人可能有点难。}
\dialogturn{5}{Eksaminator}{Vil du fortælle lidt mere om, hvordan du bruger danske podcasts eller apps til at lære dansk?}{udtalen \ipa{ˈuðˌtˢæːlən}}{能再多讲讲你是怎么用丹麦语播客或软件来学丹麦语的吗？}
\dialogturn{6}{Kursist}{Ja, selvfølgelig! Jeg lytter til en kort nyhedspodcast på femten minutter hver morgen, hvor de taler tydeligt. Hvis jeg hører et nyt ord, skriver jeg det ind i min ordbogsapp bagefter, så jeg kan øve det igen om aftenen.}{nyhedspodcast \ipa{ˈnyheðsˌpʰʌdˌkʰɑsd}}{当然！我每天早上听十五分钟语速清晰的短新闻播客。如果听到新词，之后就会记进词典应用里晚上再复习。}
\dialogturn{7}{Eksaminator}{Tror du, vi bruger telefonen for meget i dag, eller er der mest fordele ved digitale løsninger?}{fordele \ipa{ˈfɒːˌdeːlə}}{你觉得如今人们手机用得太多了吗，还是说数字化方案利大于弊？}
\dialogturn{8}{Kursist}{Jeg tror faktisk, mange af os bruger den lidt for meget om aftenen, fordi man let glemmer tiden med sociale medier. Så selvom teknologien har store fordele på arbejde og i samfundet, er det sundt selv at sætte klare grænser.}{grænser \ipa{ˈɡʁɛnsɐ}}{我确实觉得很多人晚上用得过多，因为刷社交媒体容易忘了时间。所以尽管科技在工作和社会中优势巨大，自己设好明确界限才健康。}

\paragraph*{Mundtlig Opgave 2（EMNE: Fritid og motion）Del 1 双人四图互动全真逐字稿：}
\dialogturn{1}{Kursist A}{Nå, vores emne er 'Fritid og motion'. Hvis vi kigger på de fire billeder her, hvad synes du så bedst om? Kan du godt lide at træne i et fitnesscenter, som på det første billede?}{fitnesscenter \ipa{ˈfidnəsˌsɛnˀdɐ}}{好，我们的话题是“休闲与运动”。看看这四张图，你最喜欢哪一种？你喜欢像第一张图那样在健身房训练吗？}
\dialogturn{2}{Kursist B}{Altså, jeg går faktisk i fitnesscenter to gange om ugen, fordi det er meget fleksibelt. Jeg kan træne efter arbejde, også når det regner eller er mørkt udenfor. Men nogle gange bliver det lidt kedeligt at stå på et løbebånd. Hvad med dig -- foretrækker du fitnesscenter eller at løbe ude i skoven som på billede to?}{løbebånd \ipa{ˈløːbəˌbʌnˀ} · skoven \ipa{ˈsɡʌu̯ˀən}}{其实我每周去两次健身房，因为时间非常灵活。哪怕外面下雨或天黑，我下班后也能练。不过有时在跑步机上跑也有点无聊。你呢——你更喜欢健身房，还是像图二那样在森林里户外跑？}
\dialogturn{3}{Kursist A}{Jeg kan helt klart bedst lide at løbe ude i skoven! Selvom man selvfølgelig bliver lidt våd i dårligt vejr, giver den friske luft og de grønne træer mig meget mere energi end et fitnesscenter med høj musik. Og så er det helt gratis. Hvad tænker du om holdsport i en hal, som vi ser på det tredje billede?}{friske luft \ipa{ˈfʁɛsɡə lufd} · holdsport \ipa{ˈhʌlˀˌsbʌːd}}{我明显更喜欢在森林里跑步！虽然坏天气难免淋湿，但新鲜空气和绿树比放着响亮音乐的健身房更能给我充电，而且还免费。你怎么看图三那种在体育馆里进行的团体球类运动？}
\dialogturn{4}{Kursist B}{Jeg elsker holdsport! Jeg spiller badminton hver torsdag i en lokal klub. Det bedste ved holdsport er det sociale fællesskab, for man snakker, griner og glemmer næsten, at man dyrker hård motion. Ulempen er bare, at man skal passe faste tider hver uge. Kunne du selv tænke dig at spille badminton eller volleyball?}{fællesskab \ipa{ˈfɛləsˌsɡæːˀb} · ulempen \ipa{ˈuˌlɛmbən}}{我超爱团体运动！我每周四都在当地俱乐部打羽毛球。团体运动最棒的是社交氛围，大家边说边笑，几乎忘了自己正在剧烈运动。唯一缺点是要迁就每周固定时间。你自己愿意试试打羽毛球或排球吗？}
\dialogturn{5}{Kursist A}{Ja, det lyder faktisk rigtig hyggeligt, især hvis man kan melde sig på et begynderhold sammen med en ven. Men når jeg har haft en meget stressende uge, har jeg nogle gange bare brug for ro. Derfor kan jeg også rigtig godt lide billede fire med yoga og udstrækning derhjemme. Laver du nogensinde rolige øvelser i stuen?}{stressende \ipa{ˈsdʁɛsənə} · udstrækning \ipa{ˈuðˌsdʁɛɡneŋ}}{是的，听起来真的挺惬意，尤其是和朋友一起报个初级班。不过经历高压的一周后，有时我只想静一静，因此我也很喜欢图四那样在家做瑜伽和拉伸。你在客厅做过舒缓训练吗？}
\dialogturn{6}{Kursist B}{Ja, engang imellem laver jeg strækøvelser efter arbejde, når min ryg er stiv. Det er rart, at man ikke behøver at tage nogen steder hen og kun bruger ti minutter på en måtte. Til gengæld er det lidt svært at holde koncentrationen alene derhjemme, så jeg foretrækker stadig holdsport eller fitness.}{måtte \ipa{ˈmʌdə} · stiv \ipa{ˈsdiːu̯ˀ}}{会的，下班后如果背部僵硬，我偶尔会做拉伸。不用出门、垫子上十分钟就能练完很方便。不过一个人在家不太容易保持专注，所以我还是更偏爱团体运动或健身房。}

\paragraph*{Mundtlig Opgave 2 Del 2（个人追问 3 题高分示范）：}
\dialogturn{Q1}{Tester}{Hvordan kan du bedst lide at slappe af eller dyrke motion i din fritid? Hvorfor?}{slappe af \ipa{ˈslɑbə ɑˀ}}{你空闲时最喜欢如何放松或锻炼？为什么？}
\dialogturn{A1}{Kursist}{Jeg kan allerbedst lide at løbe en lang tur i skoven alene efter arbejde, fordi det hjælper mig med at tømme hovedet for arbejdsstress. Når man hører fuglene og ser årstiderne skifte, føles det slet ikke som hård træning, men mere som en pause for hjernen.}{årstiderne \ipa{ˈɒːsˌtˢiðˀɐnə}}{我最喜欢下班后独自去森林跑长距离，这能帮我清空工作压力。听着鸟叫、看着四季变化时，完全不像苦练，更像给大脑放了个假。}
\dialogturn{Q2}{Tester}{Hvornår og hvor ofte har du tid til fritidsinteresser i en almindelig uge?}{almindelig \ipa{alˈmenˀəli}}{在平常的一周里，你一般什么时候、多久有时间投入业余爱好？}
\dialogturn{A2}{Kursist}{I en normal uge har jeg tid til motion omkring tre gange. Jeg prøver at løbe en kort tur tirsdag og torsdag aften, når aftensmaden er spist, og så går jeg en længere tur eller laver yoga søndag formiddag, hvor der ikke er så mange faste aftaler.}{formiddag \ipa{ˈfʌːmiˌdæːˀ}}{正常一周里我大概有三次锻炼时间。周二和周四吃完晚饭后会试着短跑一次，周日上午没什么固定安排，就会走远一点或做做瑜伽。}
\dialogturn{Q3}{Tester}{Kan du fortælle lidt om en fritidsaktivitet eller sport, du har prøvet eller gerne vil prøve?}{aktivitet \ipa{ɑɡtˢiviˈtˢeːˀd}}{你能讲讲你尝试过或想尝试的一项休闲活动或运动吗？}
\dialogturn{A3}{Kursist}{Jeg meldte mig til et begynderhold i badminton for et halvt år siden med min kollega Morten. I starten var det svært at ramme bolden i det høje tempo, men vores træner var meget tålmodig, og i dag elsker jeg det, fordi man både får sved på panden og nye venner.}{tålmodig \ipa{tˢʌlˈmoːˀði}}{半年前我和同事莫滕一起报了羽毛球零基础班。一开始节奏很快、很难击中球，但教练很有耐心，现在我非常喜欢这项运动，既能出汗又能交到新朋友。}

\ncdmodultestmthree{3}{Uddannelse, Fællesskab og Kultur}{全真模考卷三：音乐节义工、好评推荐与社团文化}{音乐节志愿者故事 / 小组报告与邻居代收快递聊天 / 开设自行车咖啡馆 / 三人谈社团生活 / 车行好评推荐与求职面试回信 / 大学学业与志愿工作陈述及消费购物四图双人互动}{\textit{-s} 被动 vs \textit{blive} 被动 · 介词 + 形式宾语 \textit{det} + 从句 · 现在/过去分词定语 · 省略 \textit{hvis} 条件倒装 · 篇章因果与评价转述}

% =====================================================================
% DEL 1: LÆSNING
% =====================================================================
\ncdsubsec{Del 1: Læsning（阅读理解｜50 minutter · 24 point）}

\begin{tipbox}{考场规则与配速建议（Læsning｜15--24 point = Bestået）}
\small
\textit{Læsning opgave 1, 2A, 2B, 3 og 4 udleveres samtidig. Samlet tid: 50 minutter. Kursisten må ikke bruge hjælpemidler.}（阅读五部分合订发放，\textbf{总计时 50 分钟}，全程禁用词典或手机。满分 \textbf{24 分}，\textbf{15--24 分及格（Bestået）}；0--14 分不及格（Ikke bestået）。建议配速：Opgave 1（8 min）$\rightarrow$ Opgave 2A \& 2B（10 min）$\rightarrow$ Opgave 3（14 min）$\rightarrow$ Opgave 4（13 min），最后留 5 分钟复查。）
\end{tipbox}

% ---------------------------------------------------------------------
\subsubsection*{Opgave 1 (6 point)：Frivillig på sommerens musikfestival}

\textbf{Instruktion:}
\begin{itemize}
  \item Læs teksten \textit{Frivillig på sommerens musikfestival}.
  \item Skriv de ord (1--6), der mangler. Du skal kun bruge hvert ord én gang. Der er seks ord, du ikke skal bruge. Se eksemplet (0).
\end{itemize}

\begin{emailbox}{Frivillig på sommerens musikfestival}
Freja elsker livemusik og har altid haft lyst til at opleve en stor dansk musikfestival indefra. I år har hun \textbf{(0)} \underline{\textbf{besluttet}} sig for at arbejde som frivillig på en festival i Roskilde. Det kræver en stor indsats at stå op i mange timer. Freja glæder sig meget til opgaven, \textbf{(1)}\ \rule{1.9cm}{0.4pt} hun ved, at det også bliver hårdt.

\vspace{0.35em}
Den første dag møder Freja ind om morgenen og får sin T-shirt udleveret. Hun skal stå i et stort telt ved indgangen, \textbf{(2)}\ \rule{1.9cm}{0.4pt} alle gæsterne henter deres armbånd. Hun skal også hjælpe folk med at finde deres \textbf{(3)}\ \rule{1.9cm}{0.4pt} på campingområdet, så der ikke opstår lange køer. Det er travlt om eftermiddagen, når tusindvis af gæster ankommer med telte og rygsække. Heldigvis arbejder Freja \textbf{(4)}\ \rule{1.9cm}{0.4pt} med fire andre frivillige, og de hjælper hele tiden hinanden.

\vspace{0.35em}
På den tredje dag begynder det desværre at regne kraftigt, og hendes sko bliver helt våde. Freja overvejer et øjeblik at tage tidligt hjem, \textbf{(5)}\ \rule{1.9cm}{0.4pt} hun fryser og er træt i benene. Men da hendes yndlingsband går på scenen om aftenen, glemmer hun hurtigt regnen og kulden. Til sidst er hun glad for, at hun blev. Vejret om sommeren er ikke altid perfekt, men hun vil \textbf{(6)}\ \rule{1.9cm}{0.4pt} melde sig som frivillig igen næste år.
\end{emailbox}

\begin{grammatikbox}{Ordboks（选词框 · 共 13 词：1 例词 + 6 答案 + 6 多余干扰词）}
\centering
\large\bfseries\color{NCDNavy}
\tikz[baseline=(w.base)]{\node[inner sep=1pt](w){\color{NCDGray}besluttet}; \draw[NCDRed,line width=0.9pt](w.west)--(w.east);} \qquad
selvom \qquad alene \qquad hvor \qquad derfor \qquad plads \qquad heller\\[3pt]
sammen \qquad hvornår \qquad fordi \qquad pause \qquad alligevel \qquad som
\end{grammatikbox}

% ---------------------------------------------------------------------
\subsubsection*{Opgave 2A (NY · 3 point)：Jonas skriver til sin studiegruppeven Maria}

\textbf{Instruktion:}
\begin{itemize}
  \item \textbf{Situation:} Jonas skriver til sin studiegruppeven Maria.
  \item Læs Jonas’ og Marias chat. Der er 5 beskeder fra Jonas.
  \item Find de svar fra Maria, der passer (A--F).
  \item Skriv bogstavet på det svar, der passer. Der er tre svar, du ikke skal bruge. Se eksemplet (0).
\end{itemize}

\small
\renewcommand{\arraystretch}{1.3}
\begin{longtable}{@{} >{\raggedright\arraybackslash}p{7.6cm} @{\hspace{0.4cm}} >{\raggedright\arraybackslash}p{8.8cm} @{}}
\toprule
\textbf{Chatbeskeder mellem Jonas og Maria} & \textbf{Marias mulige svar（Z 为例答，A--F 选 3 项）} \\
\midrule
\endfirsthead
\toprule
\textbf{Chatbeskeder（续表）} & \textbf{Marias mulige svar（续表）} \\
\midrule
\endhead
\bottomrule
\endfoot
\badge{NCDNavy}{Jonas}\ Hej Maria! Har du tid til at arbejde videre på vores oplæg til på torsdag?\newline
\textbf{0.}\ \badge{NCDTeal}{Maria}\ \fbox{\parbox[c][1.1em][c]{1.1em}{\centering\textbf{Z}}} &
\badge{NCDGray}{Z · Eksempel}\ \textcolor{NCDGray}{\textit{Hej Jonas! Ja, det har jeg. Skal vi ikke mødes på biblioteket i morgen formiddag?}} \\
\midrule
\badge{NCDNavy}{Jonas}\ Jo, det er en god idé. Skal vi sige kl. 10 i stueetagen?\newline
\textbf{1.}\ \badge{NCDTeal}{Maria}\ \fbox{\parbox[c][1.1em][c]{1.1em}{\centering\phantom{Z}}} &
\badge{NCDNavy}{A}\ Super, kl. 10 passer mig perfekt. Men jeg har allerede lavet hele oplægget, så vi behøver slet ikke mødes længe.\newline
\badge{NCDNavy}{B}\ Det passer mig rigtig godt med afslutningen. Men måske skal vi booke et grupperum, så der ikke er så meget støj? \\
\midrule
\badge{NCDNavy}{Jonas}\ Fint, så ses vi kl. 10.30. Jeg kan godt tage indledningen og tallene, hvis du vil lave afslutningen?\newline
\textbf{2.}\ \badge{NCDTeal}{Maria}\ \fbox{\parbox[c][1.1em][c]{1.1em}{\centering\phantom{Z}}} &
\badge{NCDNavy}{C}\ Nej, lad den endelig blive hjemme, for vi skal bare læse mine noter igennem på papir.\newline
\badge{NCDNavy}{D}\ Jeg har desværre lægetid om morgenen, så jeg kan først kl. 10.30. Har du tænkt på, hvordan vi skal fordele delene? \\
\midrule
\badge{NCDNavy}{Jonas}\ Jeg skal nok booke et grupperum i dag, så vi kan sidde i fred og ro. Skal jeg også tage min computer med?\newline
\textbf{3.}\ \badge{NCDTeal}{Maria}\ \fbox{\parbox[c][1.1em][c]{1.1em}{\centering\phantom{Z}}} &
\badge{NCDNavy}{E}\ Ja, meget gerne, for min egen er desværre ret langsom for tiden. Så kan vi lettere samle vores slides.\newline
\badge{NCDNavy}{F}\ Jeg vil hellere have, at vi aflyser mødet på biblioteket og bare sender vores tekst til hinanden i aften. \\
\midrule
\badge{NCDNavy}{Jonas}\ Aftale! Så tager jeg den med, og vi ses i morgen kl. 10.30. &
\textcolor{NCDGray}{\small （提示：从 A--F 中选出 3 项填入左侧 1--3 方框，另有 3 项多余。）} \\
\end{longtable}
\normalsize

% ---------------------------------------------------------------------
\subsubsection*{Opgave 2B (NY · 3 point)：Hanne skriver til sin nabo Ali}

\textbf{Instruktion:}
\begin{itemize}
  \item \textbf{Situation:} Hanne skriver til sin nabo Ali.
  \item Læs Hannes og Alis chat. Der er 5 beskeder fra Hanne.
  \item Find de svar fra Ali, der passer (A--F).
  \item Skriv bogstavet på det svar, der passer. Der er tre svar, du ikke skal bruge. Se eksemplet (0).
\end{itemize}

\small
\renewcommand{\arraystretch}{1.3}
\begin{longtable}{@{} >{\raggedright\arraybackslash}p{7.6cm} @{\hspace{0.4cm}} >{\raggedright\arraybackslash}p{8.8cm} @{}}
\toprule
\textbf{Chatbeskeder mellem Hanne og Ali} & \textbf{Alis mulige svar（Z 为例答，A--F 选 3 项）} \\
\midrule
\endfirsthead
\toprule
\textbf{Chatbeskeder（续表）} & \textbf{Alis mulige svar（续表）} \\
\midrule
\endhead
\bottomrule
\endfoot
\badge{NCDNavy}{Hanne}\ Hej Ali! Jeg skal på konference med arbejdet fra onsdag til søndag. Har du mulighed for at hjælpe mig med en lille ting?\newline
\textbf{0.}\ \badge{NCDTeal}{Ali}\ \fbox{\parbox[c][1.1em][c]{1.1em}{\centering\textbf{Z}}} &
\badge{NCDGray}{Z · Eksempel}\ \textcolor{NCDGray}{\textit{Hej Hanne! Ja, selvfølgelig. Hvad kan jeg hjælpe med?}} \\
\midrule
\badge{NCDNavy}{Hanne}\ Tusind tak! Der kommer nemlig en vigtig pakke til pakkeshoppen i Brugsen torsdag, og jeg har også nogle planter på altanen.\newline
\textbf{1.}\ \badge{NCDTeal}{Ali}\ \fbox{\parbox[c][1.1em][c]{1.1em}{\centering\phantom{Z}}} &
\badge{NCDNavy}{A}\ Det gør jeg. Hvornår kommer du egentlig hjem igen på søndag, og hvornår skal jeg aflevere nøglen tilbage?\newline
\badge{NCDNavy}{B}\ Jeg skal desværre selv rejse væk fra torsdag til søndag, så jeg kan ikke passe dine ting alligevel. \\
\midrule
\badge{NCDNavy}{Hanne}\ Du får koden til pakken på sms, og min ekstranøgle ligger i din postkasse i morgen aften. Planterne skal bare have vand fredag.\newline
\textbf{2.}\ \badge{NCDTeal}{Ali}\ \fbox{\parbox[c][1.1em][c]{1.1em}{\centering\phantom{Z}}} &
\badge{NCDNavy}{C}\ Intet problem! Jeg henter gerne pakken og vander dine planter. Hvad skal jeg bruge i pakkeshoppen, og hvordan kommer jeg ind?\newline
\badge{NCDNavy}{D}\ Super, så henter jeg pakken. Du behøver ikke lægge nogen nøgle til mig, for min kone har allerede din ekstranøgle. \\
\midrule
\badge{NCDNavy}{Hanne}\ Den er ret lille og vejer næsten ingenting, så du kan bare lægge den på bordet i min entré.\newline
\textbf{3.}\ \badge{NCDTeal}{Ali}\ \fbox{\parbox[c][1.1em][c]{1.1em}{\centering\phantom{Z}}} &
\badge{NCDNavy}{E}\ Fint, jeg stiller bare pakken ind i min egen stue, indtil du kommer hjem om et par uger.\newline
\badge{NCDNavy}{F}\ Det lyder helt fint. Er pakken stor og tung, og hvor skal jeg stille den, når jeg har hentet den? \\
\midrule
\badge{NCDNavy}{Hanne}\ Jeg lander først søndag aften ved 21-tiden, så jeg henter bare nøglen hos dig mandag aften. Endnu en gang tak! &
\textcolor{NCDGray}{\small （提示：从 A--F 中选出 3 项填入左侧 1--3 方框，另有 3 项多余。）} \\
\end{longtable}
\normalsize

% ---------------------------------------------------------------------
\subsubsection*{Opgave 3 (NY · 5 point)：Søren åbner en lille cykelcafé}

\textbf{Instruktion:}
\begin{itemize}
  \item Læs teksten \textit{Søren åbner en lille cykelcafé}. Der er fem afsnit.
  \item I hvert afsnit mangler der en sætning. Under hvert afsnit står der fire sætninger.
  \item Skriv bogstavet på den sætning, der passer i afsnittet. Der er tre sætninger i hvert afsnit, du ikke skal bruge.
\end{itemize}

\begin{tcolorbox}[enhanced, breakable, colback=NCDBgGray, colframe=NCDNavy!60, boxrule=0.7pt, arc=2pt, left=9pt, right=9pt, top=6pt, bottom=6pt]
\textbf{\color{NCDNavy}1.}\ Søren har arbejdet som cykelmekaniker i Odense i næsten ti år, og han er rigtig dygtig til sit fag. Alligevel har han længe drømt om at prøve noget nyt, hvor han både kan bruge sine hænder og snakke endnu mere med folk i kvarteret. En dag læser han en artikel om cykelcaféer i udlandet, hvor man kan få repareret sin cykel, mens man drikker en god kop kaffe. \rule{2.3cm}{0.4pt} For han har allerede en lille opsparing i banken og kender mange cyklister i byen.\\[4pt]
\textbf{A.} Men Søren synes, at kaffe og cykler slet ikke passer sammen.\\
\textbf{B.} Derfor vælger han at blive ansat på et endnu større værksted.\\
\textbf{C.} Han beslutter sig straks for selv at åbne sådan et sted i Odense.\\
\textbf{D.} Desværre har han aldrig før prøvet at reparere en cykel.
\end{tcolorbox}

\begin{tcolorbox}[enhanced, breakable, colback=NCDBgGray, colframe=NCDNavy!60, boxrule=0.7pt, arc=2pt, left=9pt, right=9pt, top=6pt, bottom=6pt]
\textbf{\color{NCDNavy}2.}\ Det tager dog flere måneder at finde de rigtige lokaler til den nye café. Søren kigger på tre forskellige butikker i centrum, men huslejen er alt for høj til en nystartet forretning. Til sidst finder han et gammelt hjørnelokale på en stille sidegade tæt på en populær cykelsti. \rule{2.3cm}{0.4pt} Lokalet trænger nemlig til ny maling og et bedre gulv, før han kan sætte kaffemaskinen op.\\[4pt]
\textbf{A.} Prisen er heldigvis rimelig, selvom der skal laves en del arbejde indvendigt.\\
\textbf{B.} Derfor åbner han caféen allerede dagen efter uden at ændre noget.\\
\textbf{C.} Men beliggenheden er dårlig, fordi der aldrig kommer cyklister forbi.\\
\textbf{D.} Alligevel er huslejen alt for dyr, så han må droppe drømmen helt.
\end{tcolorbox}

\begin{tcolorbox}[enhanced, breakable, colback=NCDBgGray, colframe=NCDNavy!60, boxrule=0.7pt, arc=2pt, left=9pt, right=9pt, top=6pt, bottom=6pt]
\textbf{\color{NCDNavy}3.}\ Før åbningen får Søren brug for ekstra hjælp til cafédelen, fordi han ikke selv ved så meget om bønner og mælkeskum. Han sætter derfor et opslag op på sociale medier, og allerede næste dag skriver en ung kvinde ved navn Maja, som har arbejdet som barista i fire år. De mødes til en samtale og bliver hurtigt enige om et samarbejde. \rule{2.3cm}{0.4pt} På den måde kan Søren fokusere på dæk, bremser og gear, mens Maja sørger for kaffe og nybagte boller.\\[4pt]
\textbf{A.} Men Maja vil kun arbejde der, hvis hun også må reparere cyklerne.\\
\textbf{B.} Derfor beslutter Søren sig for slet ikke at servere kaffe i butikken.\\
\textbf{C.} Hun har desværre ingen erfaring fra cafébranchen overhovedet.\\
\textbf{D.} Maja skal arbejde deltid og stå for disken, især om morgenen og i weekenden.
\end{tcolorbox}

\begin{tcolorbox}[enhanced, breakable, colback=NCDBgGray, colframe=NCDNavy!60, boxrule=0.7pt, arc=2pt, left=9pt, right=9pt, top=6pt, bottom=6pt]
\textbf{\color{NCDNavy}4.}\ Lørdag den 3. maj slår Søren og Maja endelig dørene op til \textit{Kæde \& Kaffe}. Vejret er solrigt og lunt, og allerede fra klokken ni står de første nysgerrige naboer klar udenfor. Nogle kommer kun for at drikke en latte i solen, mens andre triller ind med en punkteret cykel. \rule{2.3cm}{0.4pt} Det betyder, at stemningen i lokalet er utrolig hyggelig, selvom der er travlt hele dagen.\\[4pt]
\textbf{A.} Folk bliver dog sure over, at de skal vente hele tre minutter på kaffen.\\
\textbf{B.} Mange falder hurtigt i snak med hinanden ved de lange træborde.\\
\textbf{C.} Desværre er gaden helt tom, fordi ingen har set reklamen.\\
\textbf{D.} Derfor lukker de caféen tidligt allerede klokken elleve om formiddagen.
\end{tcolorbox}

\begin{tcolorbox}[enhanced, breakable, colback=NCDBgGray, colframe=NCDNavy!60, boxrule=0.7pt, arc=2pt, left=9pt, right=9pt, top=6pt, bottom=6pt]
\textbf{\color{NCDNavy}5.}\ Her et halvt år efter åbningen går forretningen stadig over al forventning. Hver morgen stopper mange pendlere på vej til arbejde for at få pumpet dækkene og tage en kaffe med på farten. En gang om måneden holder Søren også et lille aftenkursus, hvor gæsterne kan lære at lappe deres egne slanger. \rule{2.3cm}{0.4pt} Søren har derfor aldrig fortrudt, at han sagde sit gamle job op og fulgte sin idé.\\[4pt]
\textbf{A.} Men pendlere vil hellere tage bilen end at cykle i byen.\\
\textbf{B.} Alligevel overvejer han at sælge værkstedet, fordi det er for kedeligt.\\
\textbf{C.} De aftener er næsten altid udsolgte, og flere deltagere kommer igen næste måned.\\
\textbf{D.} Kurserne bliver dog aflyst hver gang, fordi ingen tilmelder sig.
\end{tcolorbox}

% ---------------------------------------------------------------------
\subsubsection*{Opgave 4 (NY · 7 point)：Tre personer fortæller om at være med i en forening}

\textbf{Instruktion:}
\begin{itemize}
  \item Læs de tre tekster (A, B og C), og læs spørgsmålene (1--7).
  \item Find den person, der passer til hvert af de syv spørgsmål.
  \item Sæt kryds (\textbf{X}) ud for den person (A, B eller C), der passer til spørgsmålet. Se eksemplet (0).
\end{itemize}

\small
\renewcommand{\arraystretch}{1.28}
\begin{longtable}{@{} >{\raggedright\arraybackslash}p{10.5cm} @{\hspace{0.25cm}} >{\centering\arraybackslash}p{1.9cm} @{\hspace{0.15cm}} >{\centering\arraybackslash}p{1.9cm} @{\hspace{0.15cm}} >{\centering\arraybackslash}p{1.9cm} @{}}
\toprule
\textbf{Spørgsmål} & \textbf{A. Katrine} & \textbf{B. David} & \textbf{C. Yuna} \\
\midrule
\endfirsthead
\toprule
\textbf{Spørgsmål（续表）} & \textbf{A. Katrine} & \textbf{B. David} & \textbf{C. Yuna} \\
\midrule
\endhead
\bottomrule
\endfoot
0. Hvem har været medlem af sin forening i længst tid? \textcolor{NCDGray}{\small (Eksempel)} & A\ \textcolor{NCDRed}{\textbf{X}} & B \ \underline{\hspace{0.35cm}} & C \ \underline{\hspace{0.35cm}} \\
\midrule
1. Hvem bruger også foreningen til at blive bedre til det danske sprog? & A \ \underline{\hspace{0.35cm}} & B \ \underline{\hspace{0.35cm}} & C \ \underline{\hspace{0.35cm}} \\
\midrule
2. Hvem begyndte i foreningen, fordi personen havde smerter i kroppen? & A \ \underline{\hspace{0.35cm}} & B \ \underline{\hspace{0.35cm}} & C \ \underline{\hspace{0.35cm}} \\
\midrule
3. Hvem var nødt til at lære bestemte regler og teknikker, før personen måtte deltage fuldt ud? & A \ \underline{\hspace{0.35cm}} & B \ \underline{\hspace{0.35cm}} & C \ \underline{\hspace{0.35cm}} \\
\midrule
4. Hvem oplever stor aldersforskel blandt medlemmerne i sin forening? & A \ \underline{\hspace{0.35cm}} & B \ \underline{\hspace{0.35cm}} & C \ \underline{\hspace{0.35cm}} \\
\midrule
5. Hvem optræder af og til for et publikum? & A \ \underline{\hspace{0.35cm}} & B \ \underline{\hspace{0.35cm}} & C \ \underline{\hspace{0.35cm}} \\
\midrule
6. Hvem mødes med sin forening tidligt om morgenen på hverdage? & A \ \underline{\hspace{0.35cm}} & B \ \underline{\hspace{0.35cm}} & C \ \underline{\hspace{0.35cm}} \\
\midrule
7. Hvem fortæller, at det er ret billigt at være med, fordi folk selv hjælper til? & A \ \underline{\hspace{0.35cm}} & B \ \underline{\hspace{0.35cm}} & C \ \underline{\hspace{0.35cm}} \\
\end{longtable}
\normalsize

\begin{center}
\textbf{\large\color{NCDNavy} Tre personer fortæller om at være med i en forening}
\end{center}

\begin{emailbox}{A. Katrine}
Jeg meldte mig ind i vores lokale naturforening for tre år siden, kort tid efter at jeg var flyttet til Silkeborg. Jeg kendte næsten ingen i byen, og i stedet for at sidde alene i min lejlighed ville jeg gerne ud i skoven sammen med andre. Hver anden søndag mødes vi en gruppe på femten til tyve personer og går en tur på otte til ti kilometer. Nogle gange har vi en biolog med, som fortæller om planter og fugle undervejs, og bagefter laver vi varm suppe eller kaffe over bål. Medlemmerne er mellem tyve og halvfjerds år, så man snakker med mange forskellige mennesker. Det koster kun tre hundrede kroner om året, fordi alle skiftes til at planlægge turene frivilligt. Det har været den bedste måde for mig at lære både naturen og området at kende på.
\end{emailbox}

\begin{emailbox}{B. David}
For halvandet år siden begyndte jeg i en roklub ved fjorden, fordi jeg havde fået ondt i ryggen af at sidde foran en skærm hele dagen. I starten var jeg nødt til at tage et obligatorisk begynderkursus over fire aftener, hvor vi lærte teknik og sikkerhed på vandet, før vi måtte ro ud alene. Det var faktisk ret hårdt for armene i begyndelsen! Nu ror jeg sammen med et fast hold på fire mand hver tirsdag og torsdag morgen allerede klokken halv syv, før jeg møder på arbejde. Vandet er helt stille på det tidspunkt, og det giver mig masser af energi til resten af dagen. Kontingentet er lidt højt, omkring to tusinde kroner om året, men til gengæld er alle både og redningsveste med i prisen. Om sommeren tager klubben også på en weekendtur til en anden ø.
\end{emailbox}

\begin{emailbox}{C. Yuna}
Da jeg kom til Aarhus fra Japan for to år siden, savnede jeg at lave noget kreativt efter arbejde og samtidig øve mit dansk i en afslappet stemning. En kollega tog mig med til en øveaften i et lokalt kor, og jeg blev bidt af det med det samme. Man behøver ikke at kunne læse noder eller gå til optagelsesprøve for at være med, så alle kan starte. Vi øver hver onsdag aften i et kulturhus, hvor vi synger både danske sange og engelske popnumre. På mine første øveaftener var det lidt svært at følge med i de hurtige danske tekster, men de andre medlemmer hjalp mig tålmodigt med udtalen. To gange om året giver vi små koncerter på plejehjem eller i en kirke. For mig giver koret et stærkt fællesskab, og flere af sangerne er i dag mine nære venner.
\end{emailbox}

% =====================================================================
% DEL 2: SKRIVNING
% =====================================================================
\ncdsubsec{Del 2: Skrivning（书面表达｜Opgave 1 \& Opgave 2）}

\subsubsection*{Opgave 1（半正式好评推荐信 · Anbefaling）}

\textbf{Situation:} Du har fået repareret din cykel hos den lokale cykelhandler \textit{”Byens Cykler“}, og du er meget tilfreds med servicen. Du vil skrive en anbefaling på internettet.\\
\textbf{Opgave:} Skriv en anbefaling af \textit{”Byens Cykler“}. Du skal fortælle:
\begin{itemize}
  \item hvor værkstedet ligger
  \item hvad du fik lavet
  \item hvorfor du synes, servicen var god
  \item hvem du vil anbefale værkstedet til
\end{itemize}

\begin{emailbox}{Anbefaling af ”Byens Cykler“}
\vspace{0.25em}
\rule{\linewidth}{0.4pt}\\[0.52em]
\rule{\linewidth}{0.4pt}\\[0.52em]
\rule{\linewidth}{0.4pt}\\[0.52em]
\rule{\linewidth}{0.4pt}\\[0.52em]
\rule{\linewidth}{0.4pt}\\[0.52em]
\rule{\linewidth}{0.4pt}\\[0.52em]
\rule{\linewidth}{0.4pt}\\[0.52em]
\rule{\linewidth}{0.4pt}
\end{emailbox}

\subsubsection*{Opgave 2（私人回信 · Minimum 90 ord · Tid: 30 minutter）}

\textbf{Situation:} Du har fået e-mailen her nedenfor fra din ven, Henrik.\\
\textbf{Opgave:} Læs e-mailen, og skriv et svar til Henrik. Du skal svare på alle Henriks spørgsmål, og du skal indlede og afslutte e-mailen på en passende måde.\\
\textbf{Længde:} Minimum 90 ord. \qquad \textbf{Tid:} 30 minutter.

\begin{emailbox}{Modtaget e-mail fra Henrik}
\hfill Aarhus d. 14. oktober 2025

Kære ....

Hvordan går det? Jeg hørte fra din søster, at du har været til en spændende jobsamtale for nylig! Tillykke med at du kom til samtale! Jeg er jo selv ved at søge arbejde, så jeg vil rigtig gerne høre om det.

Hvordan gik det til samtalen? Hvad handler jobbet helt præcist om? Hvordan var selve samtalen og personerne, du talte med? Og hvordan forberedte du dig på forhånd på dansk?

Jeg glæder mig meget til at høre fra dig!

Mange hilsner\\
Henrik
\end{emailbox}

\begin{emailbox}{Dit svar｜Til: Henrik \qquad Emne: Jobsamtalen}
\vspace{0.25em}
\rule{\linewidth}{0.4pt}\\[0.52em]
\rule{\linewidth}{0.4pt}\\[0.52em]
\rule{\linewidth}{0.4pt}\\[0.52em]
\rule{\linewidth}{0.4pt}\\[0.52em]
\rule{\linewidth}{0.4pt}\\[0.52em]
\rule{\linewidth}{0.4pt}\\[0.52em]
\rule{\linewidth}{0.4pt}\\[0.52em]
\rule{\linewidth}{0.4pt}
\end{emailbox}

% =====================================================================
% DEL 3: MUNDTLIG KOMMUNIKATION
% =====================================================================
\ncdsubsec{Del 3: Mundtlig kommunikation（口语交际全真题卡｜Opgave 1 \& 2）}

\subsubsection*{Mundtlig Opgave 1 (Emneopgave med Mindmap｜Del 1: 1--2 min. + Del 2: 3--4 min.)}

\begin{tipbox}{考场规则提醒（Opgave 1 · 官方红线与陈述要求）}
\small
考生需自带\textbf{两张}提前准备好的个人主题思维导图（至少 1 张须属于工作 / 实习 / 教育 / 丹麦语学习领域）。导图格内\textbf{只能填写关键词或短语（stikord），严禁写完整句子}！官方明确禁止使用复杂度过低的题目（如 \textit{Min bolig}、\textit{Min familie}、\textit{Min hverdag}）。抽签决定 1 张后，先进行 \textbf{1--2 分钟}连贯陈述（Del 1），再接受考官 \textbf{3--4 分钟}追问并简要说明理由（Del 2）。
\end{tipbox}

\begin{grammatikbox}{Mindmap 3A（已填关键词卡）：Min skole / uddannelse}
\small
\renewcommand{\arraystretch}{1.35}
\begin{tabularx}{\linewidth}{@{} >{\leavevmode\bfseries\color{NCDNavy}\raggedright\arraybackslash}p{3.4cm} @{\hspace{0.2cm}} >{\raggedright\arraybackslash}X @{\hspace{0.4cm}} >{\leavevmode\bfseries\color{NCDNavy}\raggedright\arraybackslash}p{3.4cm} @{\hspace{0.2cm}} >{\raggedright\arraybackslash}X @{}}
\toprule
\multicolumn{4}{c}{\large\bfseries\color{NCDRed} Emne: Min skole / uddannelse \quad\small\color{NCDGray}（仅含 stikord 关键词，不含整句）} \\
\midrule
1. sted / tid & København · erhvervsakademi · 2. år · man.--fre. kl. 9--15 &
2. skolen / uddannelsen & offentlig · ca. 200 studerende på årgangen · multimediedesign \\
\midrule
3. transport & cykel 20 min. langs kanalen · metro når det regner &
4. undervisningen & grafisk design · kodning · gruppeprojekter · bedst: praksis / sværest: danske rapporter \\
\midrule
5. lærere / medstuderende & uformelle lærere · danskere og internationale · fast studiegruppe &
6. mening om uddannelsen & rigtig glad · jobrelevant · godt fællesskab · travlt før eksamen \\
\bottomrule
\end{tabularx}
\end{grammatikbox}

\textbf{Del 2 考官追问模拟（Min skole / uddannelse · 4 题）：}
\begin{enumerate}
  \item \textit{Du siger, at I laver mange gruppeprojekter --- kan du bedst lide at arbejde alene eller i en gruppe? Hvorfor?}
  \item \textit{Du nævner, at skriftlige rapporter på dansk er det sværeste --- hvad gør du for at blive bedre til at skrive på dansk?}
  \item \textit{Hvordan er forholdet mellem lærere og studerende på din uddannelse i forhold til dit hjemland?}
  \item \textit{Hvilken slags job drømmer du om at få, når du er færdig med uddannelsen?}
\end{enumerate}

\begin{grammatikbox}{Mindmap 3B（已填关键词卡）：Mit frivillige arbejde}
\small
\renewcommand{\arraystretch}{1.35}
\begin{tabularx}{\linewidth}{@{} >{\leavevmode\bfseries\color{NCDNavy}\raggedright\arraybackslash}p{3.4cm} @{\hspace{0.2cm}} >{\raggedright\arraybackslash}X @{\hspace{0.4cm}} >{\leavevmode\bfseries\color{NCDNavy}\raggedright\arraybackslash}p{3.4cm} @{\hspace{0.2cm}} >{\raggedright\arraybackslash}X @{}}
\toprule
\multicolumn{4}{c}{\large\bfseries\color{NCDRed} Emne: Mit frivillige arbejde \quad\small\color{NCDGray}（仅含 stikord 关键词，不含整句）} \\
\midrule
1. sted / organisation & kulturhus i midtbyen · folkekøkken ”Naboens Bord“ &
2. tidsforbrug & hver anden torsdag aften · kl. 16.30--20.30 · ca. 8 timer/md. \\
\midrule
3. arbejdsopgaver & skære grøntsager · dække langborde · tage imod gæster · rydde op &
4. mennesker / hold & 10 faste frivillige · alle aldre · naboer og studerende · god holdånd \\
\midrule
5. udbytte / dansk & øver hverdagssnak · nye danske venner · lærer om madkultur &
6. mening om frivilligt arbejde & meningsfuldt · giver energi · lidt hårdt for benene · stolt af indsatsen \\
\bottomrule
\end{tabularx}
\end{grammatikbox}

\textbf{Del 2 考官追问模拟（Mit frivillige arbejde · 4 题）：}
\begin{enumerate}
  \item \textit{Hvordan fandt du frem til lige præcis dette folkekøkken, og hvorfor valgte du at blive frivillig?}
  \item \textit{Hvad kan du bedst lide af dine opgaver i køkkenet eller ved bordene, og er der noget, der er mindre sjovt?}
  \item \textit{Hvad har det frivillige arbejde betydet for dit danske sprog og dit netværk i byen?}
  \item \textit{Synes du, at det er let eller svært at passe frivilligt arbejde ved siden af hverdag og arbejde?}
\end{enumerate}

% ---------------------------------------------------------------------
\subsubsection*{Mundtlig Opgave 2 (Paropgave Del 1: ca. 4 min. + Individuel Del 2: ca. 3 min. pr. kursist)}

\begin{grammatikbox}{Kursistark（Opgaveark · EMNE: Indkøb og forbrug）}
\small
\textbf{\color{NCDRed}Opgave:} \textit{Testeren introducerer dig til opgaven.} \qquad \textbf{\color{NCDNavy}Emne:} \textit{Indkøb og forbrug.}\\[4pt]
\renewcommand{\arraystretch}{1.35}
\begin{tabularx}{\linewidth}{@{} >{\raggedright\arraybackslash}X @{\hspace{0.4cm}} >{\raggedright\arraybackslash}X @{}}
\toprule
\badge{NCDNavy}{[图片 1: 本地超市常规采购]}\newline
\textit{En kunde med en indkøbskurv blandt hylder og gule tilbudsskilte i et almindeligt dansk supermarked.} &
\badge{NCDTeal}{[图片 2: 网上购物与送货]}\newline
\textit{En person bestiller dagligvarer og tøj på sin bærbare computer, mens et par pakker allerede står ved hoveddøren.} \\
\midrule
\badge{NCDGold}{[图片 3: 跳蚤市场与二手店]}\newline
\textit{Besøgende kigger på brugte møbler, bøger, lamper og tøj i boderne på et udendørs loppemarked eller i en genbrugsbutik.} &
\badge{NCDRed}{[图片 4: 农场直购/本地集市]}\newline
\textit{Et lokalt torvemarked eller en gårdbutik med friske økologiske grøntsager, æg og brød direkte fra producenten.} \\
\bottomrule
\end{tabularx}
\end{grammatikbox}

\begin{emailbox}{Testerark（考官引导语与 Del 1 / Del 2 考核指令）}
\small
\textbf{\color{NCDNavy}Del 1 (Paropgave -- ca. 4 min.)}\\
Opgavearket lægges foran kursisterne, der præsenteres for emnet på følgende måde:\\[2pt]
\textit{”Emnet er ’Indkøb og forbrug’. Her er nogle billeder, der viser forskellige måder at købe ind på. I skal tale med hinanden om, hvad I synes om de forskellige måder at handle på, og hvorfor. I skal stille spørgsmål til hinanden og svare på hinandens spørgsmål.“}\\[4pt]
Hvis kursisterne ikke udtrykker deres præferencer i forhold til alle fire billeder, må tester sørge for, at dette sker ved at spørge:\\[2pt]
\textit{”Hvad synes I om at handle i et almindeligt supermarked / købe ind på nettet / handle på loppemarked eller i genbrugsbutik / købe varer i en gårdbutik eller på et lokalt marked? Hvorfor? / Hvorfor ikke?“}\\[6pt]
\textbf{\color{NCDNavy}Del 2 (Individuel opgave -- ca. 3 min. pr. kursist)}\\
Tester stiller hver kursist nedenstående spørgsmål:
\begin{enumerate}
  \item \textit{Hvordan kan du bedst lide at købe ind --- og hvorfor?}
  \item \textit{Hvad er vigtigst for dig, når du handler i hverdagen (fx pris, kvalitet, tid, klima eller økologi)?}
  \item \textit{Kan du fortælle lidt om et godt eller dårligt køb, du har lavet for nylig? (hvad, hvor, hvornår, hvorfor godt/dårligt)}
\end{enumerate}
\end{emailbox}

% =====================================================================
% DEL 4: RETTEARK, DETALJERET ANALYSE & HØJSCORE-MODELLER
% =====================================================================
\ncdsubsec{Del 4: Retteark, Detaljeret Analyse \& Højscore-Modeller}

\subsubsection*{4.1 官方 Retteark 评分卡（Læsning 24 point \& Skrivning / Mundtlig 评估基准）}

\small
\renewcommand{\arraystretch}{1.28}
\begin{longtable}{@{} >{\leavevmode\bfseries\color{NCDNavy}\raggedright\arraybackslash}p{3.3cm} @{\hspace{0.25cm}} >{\raggedright\arraybackslash}p{9.45cm} @{\hspace{0.25cm}} >{\centering\arraybackslash}p{1.8cm} @{\hspace{0.25cm}} >{\centering\arraybackslash}p{1.7cm} @{}}
\toprule
\textbf{模块 / 题号} & \textbf{Rettenøgle（标准答案）} & \textbf{分值} & \textbf{自测得分} \\
\midrule
\endfirsthead
\toprule
\textbf{模块 / 题号（续表）} & \textbf{Rettenøgle（标准答案）} & \textbf{分值} & \textbf{自测得分} \\
\midrule
\endhead
\bottomrule
\endfoot
\textbf{Læsning Opgave 1}\newline（6 题） &
1. \textbf{selvom} \quad 2. \textbf{hvor} \quad 3. \textbf{plads} \quad 4. \textbf{sammen} \quad 5. \textbf{fordi} \quad 6. \textbf{alligevel}\newline
\textcolor{NCDGray}{\small （例词 0：besluttet｜多余干扰词：alene, derfor, heller, hvornår, pause, som）} & 6 point & \rule{0.75cm}{0.4pt} / 6 \\
\midrule
\textbf{Læsning Opgave 2A}\newline（3 题） &
1. \textbf{D} \quad 2. \textbf{B} \quad 3. \textbf{E}\newline
\textcolor{NCDGray}{\small （例答 0：Z｜多余干扰项：A, C, F）} & 3 point & \rule{0.75cm}{0.4pt} / 3 \\
\midrule
\textbf{Læsning Opgave 2B}\newline（3 题） &
1. \textbf{C} \quad 2. \textbf{F} \quad 3. \textbf{A}\newline
\textcolor{NCDGray}{\small （例答 0：Z｜多余干扰项：B, D, E）} & 3 point & \rule{0.75cm}{0.4pt} / 3 \\
\midrule
\textbf{Læsning Opgave 3}\newline（5 题） &
1. \textbf{C} \quad 2. \textbf{A} \quad 3. \textbf{D} \quad 4. \textbf{B} \quad 5. \textbf{C} & 5 point & \rule{0.75cm}{0.4pt} / 5 \\
\midrule
\textbf{Læsning Opgave 4}\newline（7 题） &
1. \textbf{C} \quad 2. \textbf{B} \quad 3. \textbf{B} \quad 4. \textbf{A} \quad 5. \textbf{C} \quad 6. \textbf{B} \quad 7. \textbf{A}\newline
\textcolor{NCDGray}{\small （例题 0 为 A）} & 7 point & \rule{0.75cm}{0.4pt} / 7 \\
\midrule
\textbf{Læsning 总计} & \textbf{Bestået：15--24 point（及格）} \quad|\quad \textbf{Ikke bestået：0--14 point} & \textbf{24 point} & \rule{0.75cm}{0.4pt} / 24 \\
\end{longtable}
\normalsize

\begin{tipbox}{Skrivning \& Mundtlig Opgave 1--2 整体评估自检准绳（Bestået 底线）}
\small
\begin{itemize}
  \item \textbf{Skrivning Opgave 1（Anbefaling）}：\ding{172} 必须完整涵盖指定 4 个要点（地点、维修项目、为何满意、推荐给谁）；\ding{173} 语气礼貌得体，灵活运用 \textit{som / fordi / selvom} 与过去时叙事 + 结尾推荐句型（如 \textit{Jeg vil varmt anbefale...}）。
  \item \textbf{Skrivning Opgave 2（Privat e-mail）}：\ding{172} 字数硬性要求 \textbf{minimum 90 ord}（建议写到 105--135 词）；\ding{173} 回答 Henrik 的全部 4 个问题；\ding{174} 具备地道开头（\textit{Kære Henrik}）与落款（\textit{Kærlig/Mange hilsner}）；\ding{175} 保持正确的 B1 从句语序与时态呼应。
  \item \textbf{Mundtlig Opgave 1 \& 2}：Opgave 1 仅凭 \textit{stikord} 连贯陈述 1--2 分钟，并在 Del 2 给出带简短理由（\textit{fordi / derfor}）的回应；Opgave 2 Del 1 两人必须交互问答，对\textbf{全部 4 张图}都明确表达偏好与原因，Del 2 能围绕个人消费习惯展开 3 分钟问答。
\end{itemize}
\end{tipbox}

% ---------------------------------------------------------------------
\subsubsection*{4.2 阅读 Opgave 1--4 逐题深度解析、全文精译与核心词表}

\paragraph*{一、Opgave 1（音乐节志愿者）完形语法逻辑解析与全文精译}
\begin{itemize}
  \item \textbf{(1) selvom}：前后形成“虽然知道辛苦，但依然非常期待”的让步关系（\textit{glæder sig meget..., selvom hun ved...}）。干扰词 \textit{fordi} 逻辑反向（不会“因为工作辛苦而期待”）；\textit{derfor} 属于副词不能直接作从句连词，且后接主句需倒装。
  \item \textbf{(2) hvor}：修饰地点名词 \textit{et stort telt ved indgangen}，引导定语从句说明“在那个帐篷里所有客人都来领取手环”。干扰词 \textit{som} 不能作地点状语（从句中 \textit{alle gæsterne} 和 \textit{deres armbånd} 已完整）；\textit{hvornår} 为疑问副词，不可作定语从句关系词。
  \item \textbf{(3) plads}：放在物主代词 \textit{deres} 之后充当单数共性名词（\textit{finde deres plads på campingområdet}，“找到他们在露营区的营位”）。唯一另一个名词干扰词 \textit{pause}（“找到他们的休息时间”）语义不通。
  \item \textbf{(4) sammen}：构成固定搭配 \textit{arbejde sammen med...}（“与……一起工作”）。干扰词 \textit{alene}（“独自”）与后面的 \textit{med fire andre} 及 \textit{hjælper hinanden} 直接矛盾。
  \item \textbf{(5) fordi}：解释前文“考虑提前回家”的原因（“因为她冷得发抖而且双腿很累”），从句语序 \textit{hun fryser} 正确。干扰词 \textit{derfor} 因果倒置且需触发 V2 倒装（*\textit{derfor fryser hun}）。
  \item \textbf{(6) alligevel}：放在情态动词 \textit{vil} 与实义动词 \textit{melde} 之间，呼应前半句转折连词 \textit{men}（“夏天天气未必完美，但她\textbf{还是}想明年再来当志愿者”）。干扰词 \textit{heller} 只能用于否定句（需配合 \textit{ikke}）。
\end{itemize}

\begin{transbox}
\small
\textbf{\color{NCDTeal}[Opgave 1 全文中文精译]}\\
弗蕾娅热爱现场音乐，一直很想亲身体验一次丹麦大型音乐节的幕后生活。今年她\textbf{（0）决定}去罗斯基勒的一场音乐节做志愿者。这需要付出很大精力、长时间站立。不过弗蕾娅非常期待这项任务，\textbf{（1）尽管}她知道这也会很辛苦。第一天，弗蕾娅一早到岗并领到了自己的 T 恤。她要守在入口处的一个大帐篷里，\textbf{（2）在那里}所有观众会领取他们的手环。她还得帮大家在露营区找到各自的\textbf{（3）营位}，以免排起长队。下午当成千上万名观众背着帐篷和双肩包同时涌入时，工作非常忙碌。幸好弗蕾娅是和另外四名志愿者\textbf{（4）一起}工作，他们一直互相照应。到了第三天，不幸下起了大雨，她的鞋子全湿透了。弗蕾娅有那么一瞬间考虑提前回家，\textbf{（5）因为}她既冷又觉得腿酸。但当她最喜欢的乐队在晚上登台时，她很快便忘记了雨水和寒冷。到最后，她很庆幸自己留了下来。夏天的天气并不总是完美，但她明年\textbf{（6）依然}想再报名当志愿者。
\end{transbox}

\paragraph*{二、Opgave 2A \& 2B（双对话气泡补全）跨回合衔接解析与全文精译}
\begin{itemize}
  \item \textbf{Opgave 2A（Jonas \& Maria）}：
    \begin{itemize}
      \item \textbf{1. D}：Jonas 提议十点在一楼见面；下一句 Jonas 回应 \textit{Fint, så ses vi kl. 10.30. Jeg kan godt tage indledningen... hvis du vil lave afslutningen?}，说明 Maria 必先将时间改到 10:30 并询问“如何分工”（\textit{fordele delene}）。\textbf{排除 A}（A 说 10 点合适且自己写完了，与 10:30 及分工矛盾）。
      \item \textbf{2. B}：承接 Jonas 询问“你做结尾部分可以吗？”，Maria 在 B 中先接受分工（\textit{Det passer mig rigtig godt med afslutningen}），再提议预订小组讨论室以免太吵；Jonas 紧接着在下一条回复 \textit{Jeg skal nok booke et grupperum i dag}，精准扣合。
      \item \textbf{3. E}：回应 Jonas 问“我也要带电脑吗？”，Maria 在 E 中表示“非常欢迎，因为我自己的电脑最近很慢”，且便于合并幻灯片；最后 Jonas 确认 \textit{Så tager jeg den med}。\textbf{排除 C}（拒绝带电脑，与结尾句直接冲突）与 \textbf{F}（提议取消图书馆会面，与订讨论室及明天见矛盾）。
    \end{itemize}
  \item \textbf{Opgave 2B（Hanne \& Ali）}：
    \begin{itemize}
      \item \textbf{1. C}：Hanne 说明有两件事（周四去 Brugsen 包裹点取件 + 给阳台植物浇水）；下一条 Hanne 回复了“短信发取件码、备用钥匙放信箱、周五浇水”，可见 Ali 在 C 中不仅一口答应，还问了“包裹点需要什么”和“怎么进你家”。\textbf{排除 B}（自己要出去旅行）与 \textbf{D}（自称已经有钥匙）。
      \item \textbf{2. F}：承接 Hanne 交代取件码与钥匙后，Ali 在 F 中询问“包裹又大又重吗？取完放哪儿？”，紧接着 Hanne 回复 \textit{Den er ret lille og vejer næsten ingenting, så du kan bare lægge den på bordet i min entré}，问与答严丝合缝。\textbf{排除 E}（说放自己客厅等几周，与放进 Hanne 玄关且周日回来矛盾）。
      \item \textbf{3. A}：Ali 先答应照办（\textit{Det gør jeg}），再问“你周日几点回来、我什么时候还钥匙？”，正好引出 Hanne 最后一条交代周日晚九点落地及周一取钥匙的回复。
    \end{itemize}
\end{itemize}

\begin{transbox}
\small
\textbf{\color{NCDTeal}[Opgave 2A \& 2B 对话中文精译]}\\
\textbf{2A（乔纳斯与玛丽亚）}：乔纳斯：嗨玛丽亚！你有空继续准备周四的报告吗？$\rightarrow$ 玛丽亚（Z 例）：嗨乔纳斯！有的。我们明天上午在图书馆见怎么样？$\rightarrow$ 乔纳斯：好主意。约十点在一楼好吗？$\rightarrow$ 玛丽亚（1. D）：抱歉我早上有医生预约，十点半才能到。你想过怎么分配各部分了吗？$\rightarrow$ 乔纳斯：行，那十点半见。我可以负责开头和数据，你来做结尾好吗？$\rightarrow$ 玛丽亚（2. B）：我做结尾很合适。不过也许我们该订个小组讨论室，这样不会那么吵？$\rightarrow$ 乔纳斯：我今天去订讨论室，好让我们安静坐着。我还要带电脑吗？$\rightarrow$ 玛丽亚（3. E）：太好了，我自己的电脑最近特别慢。这样我们也更好合并幻灯片。$\rightarrow$ 乔纳斯：说定了！我把电脑带上，明天十点半见。\\[3pt]
\textbf{2B（汉娜与阿里）}：汉娜：嗨阿里！我周三到周日要出差开会，能帮我个小忙吗？$\rightarrow$ 阿里（Z 例）：嗨汉娜！当然可以，有什么能帮忙的？$\rightarrow$ 汉娜：太感谢了！周四有个重要包裹会到 Brugsen 的取件点，我阳台上还有一些花。$\rightarrow$ 阿里（1. C）：没问题！我乐意帮忙取件和浇花。去取件需要什么信息？我怎么进你公寓？$\rightarrow$ 汉娜：取件码会发你短信，我的备用钥匙明晚放你信箱。植物周五浇水就行。$\rightarrow$ 阿里（2. F）：听起来很容易。包裹又大又重吗？取回来放哪儿？$\rightarrow$ 汉娜：它挺小的而且几乎不沉，放我玄关桌子上就行。$\rightarrow$ 阿里（3. A）：好的。你周日到底几点回来？我什么时候还钥匙？$\rightarrow$ 汉娜：我周日晚上九点左右落地，周一晚上我去你那儿拿钥匙就行。再次感谢！
\end{transbox}

\paragraph*{三、Opgave 3（开设自行车咖啡馆）段落逻辑补句解析与全文精译}
\begin{itemize}
  \item \textbf{第 1 段 $\rightarrow$ C}：前句写到 Søren 读到国外的自行车咖啡馆报道；后句以 \textit{For} 引出原因“因为他银行里已经有一笔小额存款，而且认识城里很多骑车人”。C 项“他立刻决定自己在欧登塞也开这样一家店”完美构成因果链条（\textbf{排除 A/B/D}：均否定开店或否定十年修车经验）。
  \item \textbf{第 2 段 $\rightarrow$ A}：前句说明他终于在安静小街且靠近热门自行车道旁找到一间旧街角店面；后句用 \textit{nemlig} 解释“店铺确实需要重新粉刷并修整地板才能摆上咖啡机”。A 项“幸好租金合理，尽管内部还得做不少装修”承上启下，\textit{arbejde indvendigt} 直接对应后句的刷墙铺地。
  \item \textbf{第 3 段 $\rightarrow$ D}：前句写咖啡师 Maja 来面试并达成合作意向；后句以 \textit{På den måde} 总结“这样一来，Søren 可以专注轮胎、刹车和变速，而 Maja 负责咖啡和现烤小面包”。D 项“Maja 将做兼职并负责吧台，尤其是早间和周末”准确补出具体的分工方案。
  \item \textbf{第 4 段 $\rightarrow$ B}：前句描述开业当天阳光明媚，有的客人只喝拿铁，有的推着爆胎的自行车进来；后句总结 \textit{Det betyder, at stemningen i lokalet er utrolig hyggelig}（“这意味着店里气氛格外温馨舒适”）。B 项“许多人在长木桌旁很快就聊了起来”正好解释了为何气氛温馨。
  \item \textbf{第 5 段 $\rightarrow$ C}：空前一句提到 Søren 每月还举办一次教客人自己补胎的晚间小课程；后句总结“因此 Søren 从未后悔辞掉旧工作”。C 项中的回指短语 \textit{De aftener}（“那些夜晚”）精准指代前文的 \textit{et lille aftenkursus}，并交代场场爆满的盛况。
\end{itemize}

\begin{transbox}
\small
\textbf{\color{NCDTeal}[Opgave 3 全文中文精译]}\\
\textbf{1.} 索伦在欧登塞当了近十年自行车修理工，手艺非常精湛。尽管如此，他一直梦想尝试点新事物，既能动手干活，又能和街坊邻里多聊聊天。一天，他读到一篇介绍国外“自行车咖啡馆”的文章——在那里人们可以一边喝杯好咖啡，一边修好自行车。\textbf{[C. 他立刻决定自己在欧登塞也开这样一家店。]} 因为他银行里已经有一笔小额存款，而且认识城里许多骑车人。\textbf{2.} 不过找到合适的咖啡馆铺面花了好几个月。索伦看了市中心的三家店铺，但租金对新创小店来说太高了。最后，他在一条靠近热门自行车道的安静支路上找到一间老街角铺面。\textbf{[A. 幸好价格合理，尽管内部还得做不少工程。]} 因为店里在安装咖啡机前确实需要重新粉刷和铺设更好的地板。\textbf{3.} 开业前，索伦在咖啡区需要额外帮手，因为他自己对咖啡豆和打奶泡了解不多。于是他在社交媒体发了招聘帖，第二天就有一位做了四年咖啡师的年轻姑娘玛雅联系他。两人面谈后很快达成合作。\textbf{[D. 玛雅将兼职负责吧台，尤其是早上和周末。]} 这样一来，索伦可以专心处理轮胎、刹车和变速器，而玛雅负责咖啡和现烤面包。\textbf{4.} 五月三日周六，索伦和玛雅终于推开了“链条与咖啡”的大门。天气晴朗温暖，早上九点不到第一批好奇的邻居就等在门外。有人只是来晒太阳喝拿铁，也有人推着扎破胎的车进来。\textbf{[B. 许多人在长木桌旁很快聊了起来。]} 这使得店里的氛围虽然全天忙碌却无比温馨。\textbf{5.} 开业半年后，生意依然远超预期。每天早上都有不少上班族在通勤路上停下来打气并顺手买杯外带咖啡。索伦每月还办一次教客人自己补内胎的晚间小课堂。\textbf{[C. 那些夜晚几乎场场报满，还有好几位学员下个月会再来。]} 因此索伦从未后悔过辞去旧工作、追随自己的想法。
\end{transbox}

\paragraph*{四、Opgave 4（三人谈社团生活）同义定位解析与三人自述精译}
\begin{itemize}
  \item \textbf{(0) A [Eksempel]}：Katrine 三年前入会（\textit{for tre år siden}），David 一年半前（\textit{halvandet år siden}），Yuna 两年前（\textit{to år siden}），故 Katrine 入会时间最长。
  \item \textbf{1. C (Yuna)}：明确提到加入合唱团是为了同时在轻松氛围里练习丹麦语（\textit{samtidig øve mit dansk}），且其他成员耐心帮她纠正发音。
  \item \textbf{2. B (David)}：因为整天坐在电脑屏幕前导致背痛（\textit{havde fået ondt i ryggen}），才加入赛艇俱乐部。对应题干 \textit{havde smerter i kroppen}。
  \item \textbf{3. B (David)}：必须先完成四晚的必修入门课程学习划桨技巧与水上安全，才能单独划出去（\textit{obligatorisk begynderkursus... lærte teknik og sikkerhed... før vi måtte ro ud alene}）。Yuna 明确相反（无需识谱或试音）。
  \item \textbf{4. A (Katrine)}：成员年龄跨度从二十岁到七十岁（\textit{mellem tyve og halvfjerds år}），对应题干 \textit{stor aldersforskel}。
  \item \textbf{5. C (Yuna)}：每年两次在养老院或教堂举办小型音乐会（\textit{giver vi små koncerter}），对应题干 \textit{optræder ... for et publikum}。
  \item \textbf{6. B (David)}：每周二、周四早上六点半上班前一起划船（\textit{tirsdag og torsdag morgen allerede klokken halv syv, før jeg møder på arbejde}）。Katrine 在周日，Yuna 在周三晚。
  \item \textbf{7. A (Katrine)}：年费只需三百克朗，因为大家志愿轮流策划徒步活动（\textit{koster kun tre hundrede kroner om året, fordi alle skiftes til at planlægge turene frivilligt}）。David 反而说会费偏高（约两千克朗）。
\end{itemize}

\begin{transbox}
\small
\textbf{\color{NCDTeal}[Opgave 4 三人自述中文精译]}\\
\textbf{A. 卡特琳娜}：三年前刚搬到锡尔克堡不久，我就加入了当地的自然协会。那时我在城里几乎谁也不认识，不想一个人待在公寓里，很想和大家一起走进森林。每隔一个周日，我们十五到二十个人聚在一起徒步八到十公里。有时会有生物学家同行，沿途讲解植物和鸟类，之后我们还会架起篝火煮热汤或咖啡。会员年龄在二十到七十岁之间，因此能和各种各样的人交流。每年会费只需三百克朗，因为大家会自愿轮流策划路线。这对我来说是认识当地自然风光与周边地区最好的方式。\\
\textbf{B. 大卫}：一年半前我加入了峡湾边的赛艇俱乐部，因为整天坐在屏幕前让我腰背疼痛。一开始我必须参加四晚的必修入门课，学习水上技术与安全规则后才允许独自出海划艇。刚开始手臂真的挺吃力的！现在我每周二、周四早上六点半上班前就和固定的四人小组一起划艇。那个时段水面非常平静，能为我一整天注入充沛精力。会费略高，一年大约两千克朗，不过所有赛艇和救生衣都包含在内。夏天俱乐部还会组织去另一座岛的周末短途旅行。\\
\textbf{C. 优娜}：两年前我从日本来到奥胡斯时，很想在下班后做些有创意的事，同时在轻松的氛围里练习丹麦语。一位同事带我去了一家本地合唱团的排练夜，我一下子就迷上了。参加这里不需要会看乐谱也不用参加入学试音，所以人人都能加入。我们每周三晚上在文化馆排练，既唱丹麦歌曲也唱英文流行曲。最初几次排练时，我有点跟不上语速较快的丹麦歌词，但其他团员都耐心地帮我纠正发音。每年两次我们会在养老院或教堂举办小型演出。合唱团给了我强烈的归属感，如今好几位团员都成了我的挚友。
\end{transbox}

\paragraph*{五、核心高频词表（Gloser｜15 个必背 B1 屈折词条）}

\begin{glosertable}
名词 & en vagt & \textit{(vagten, vagter, vagterne)} & 值班，班次（如志愿者轮班） \\
名词 & et armbånd & \textit{(armbåndet, armbånd, armbåndene)} & 手环，入场腕带（单复数同形） \\
动词 & at overveje & \textit{(overvejer, overvejede, har overvejet)} & 考虑，斟酌 \\
副词 & alligevel & （句首触发 V2；句中置于中心状语位） & 仍然，尽管如此 \\
动词 & at fordele & \textit{(fordeler, fordelte, har fordelt)} & 分配，分摊 \\
名词 & en støj & \textit{(støjen, ingen flertal)} & 噪音，嘈杂声 \\
名词 & et lokale & \textit{(lokalet, lokaler, lokalerne)} & 店铺场地，房间，空间 \\
动词 & at opstå & \textit{(opstår, opstod, er opstået)} & 出现，产生（完成时用 \textit{er}！） \\
动词 & at fortryde & \textit{(fortryder, fortrød, har fortrudt)} & 后悔（不规则强变化！） \\
名词 & en opsparing & \textit{(opsparingen, opsparinger, -erne)} & 储蓄，积蓄 \\
名词 & en forening & \textit{(foreningen, foreninger, foreningerne)} & 协会，社团，俱乐部 \\
名词 & et kontingent & \textit{(kontingentet, kontingenter, -erne)} & 会费，会员费 \\
名词 & et fællesskab & \textit{(fællesskabet, fællesskaber, -ne)} & 共同体，社群归属感 \\
形容词 & pålidelig & \textit{(pålidelig, pålideligt, pålidelige)} & 可靠的，值得信赖的 \\
动词 & at anbefale & \textit{(anbefaler, anbefalede, har anbefalet)} & 推荐，建议 \\
\end{glosertable}

% =====================================================================
% 4.3 写作范文
% =====================================================================
\subsubsection*{4.3 写作双任务双档高分示范与点评（Skrivning Modeller）}

\begin{emailbox}{Opgave 1 范文一：B1 稳过版（109 词｜清晰覆盖 4 要点 · 句式自然）}
\textbf{Anbefaling af Byens Cykler}

Jeg vil gerne anbefale cykelhandleren \textit{Byens Cykler}, som ligger på Søndergade 42 i Aarhus lige ved biblioteket.

I sidste uge punkterede mit bagdæk på vej hjem fra arbejde, og mine bremser virkede heller ikke ordentligt. Jeg kørte derfor cyklen ind til værkstedet om morgenen uden at have bestilt tid.

Servicen var helt fantastisk. Mekanikeren tog imod mig med et smil og forklarede tydeligt, hvad der skulle laves, og hvad det ville koste. Allerede samme eftermiddag var cyklen færdig, og de havde endda smurt kæden gratis.

Jeg vil klart anbefale \textit{Byens Cykler} til alle i byen, der ønsker hurtig, ærlig og venlig hjælp til deres cykel.

\vspace{0.3em}
\rule{\linewidth}{0.4pt}\\[0.2em]
\small\textbf{\color{NCDTeal}[考官点评 · 109 ord]}：完整覆盖全部 4 项要求（地址、维修项目、服务优点、推荐对象）。语法亮点包括：\ding{172} 关系从句主语标配 \textit{hvad \textbf{der} skulle laves}；\ding{173} 句首状语倒装 \textit{Allerede samme eftermiddag \textbf{var} cyklen færdig}；\ding{174} 过去完成时 \textit{havde smurt} 补充背景；\ding{175} 不定式介词短语 \textit{uden at have bestilt}。
\end{emailbox}

\begin{emailbox}{Opgave 1 范文二：B1+ 高分示范版（141 词｜被动语态 + 介词接从句 + 高级衔接）}
\textbf{Den bedste cykelservice i byen}

Hvis du leder efter et pålideligt cykelværksted, kan jeg på det varmeste anbefale \textit{Byens Cykler}, som ligger centralt på Søndergade 42 i Odense, fem minutter fra banegården.

Forleden opdagede jeg, at gearet på min cykel var gået i stykker, og samtidig trængte både kæden og forlygten til at blive skiftet. \da{Selvom jeg ikke havde aftalt} nogen tid på forhånd, \dabold{blev cyklen taget ind} med det samme.

\da{Grunden til, at} jeg er så tilfreds med servicen, er den ærlige rådgivning. Inden reparationen gik i gang, fik jeg en klar pris, så jeg ikke \dabold{blev overrasket} ved kassen. Desuden sendte værkstedet en sms efter fire timer, og cyklen kørte som ny.

Jeg vil især anbefale stedet til pendlere og studerende, som er afhængige af deres cykel i hverdagen og sætter pris på god kvalitet til en fair pris.

\vspace{0.3em}
\rule{\linewidth}{0.4pt}\\[0.2em]
\small\textbf{\color{NCDNavy}[高分亮点 · 141 ord]}：\ding{172} 熟练运用 \textit{blive} 被动描述具体维修过程（\textit{blive skiftet / blev taget ind / blev overrasket}）；\ding{173} 让步从句内否定前置 + 主句倒装（\textit{Selvom jeg \textbf{ikke} havde aftalt..., \textbf{blev} cyklen...}）；\ding{174} 精准使用 B1 评价结构 \textit{Grunden til, at... er...} 与 \textit{sætter pris på}。
\end{emailbox}

\begin{emailbox}{Opgave 2 范文一：B1 稳过版（112 词｜严格 $\ge$ 90 词 · 回答全部 4 问）}
Kære Henrik

Tak for din søde mail!

Ja, det gik rigtig godt til jobsamtalen i torsdags, og jeg blev tilbudt jobbet allerede i går! Jeg er selvfølgelig utrolig glad.

Stillingen er som logistikmedarbejder på et mindre lager. Arbejdet handler om at pakke ordrer, tale med chauffører og holde styr på varerne i systemet.

Samtalen varede fyrre minutter og var meget afslappet. Jeg talte med chefen og en medarbejder. De var meget venlige og spurgte til min erfaring og mit dansk.

Før samtalen havde jeg forberedt mig grundigt på dansk. Jeg læste om firmaet, skrev små noter og øvede svarene højt med min nabo. Held og lykke med din jobsøgning!

Kærlig hilsen\\
Klara

\vspace{0.3em}
\rule{\linewidth}{0.4pt}\\[0.2em]
\small\textbf{\color{NCDTeal}[考官点评 · 112 ord]}：远超 90 词底线，称呼落款得体。顺畅使用过去被动 \textit{blev tilbudt} 与过去完成时 \textit{havde forberedt mig} 区分“面试前准备”与“面试当天经过”，完整回答 Henrik 的全部四个问题，信息密度充实。
\end{emailbox}

\begin{emailbox}{Opgave 2 范文二：B1+ 高分示范版（138 词｜介词 + om 从句 + 对比连接）}
Kære Henrik

Tak for din besked -- det er dejligt at høre, at du også søger arbejde!

Det gik over al forventning til min jobsamtale i sidste uge, for i går ringede de og tilbød mig stillingen.

Jobbet er som kontorassistent i en grøn virksomhed, hvor mine vigtigste opgaver bliver at svare på kundehenvendelser, planlægge møder og hjælpe med mindre regnskaber.

\da{Selvom jeg i starten var lidt spændt på, om} mit danske sprog var godt nok, \dabold{blev samtalen} hurtigt tryg og uformel. Vi drak kaffe, og de to ledere lagde vægt på min motivation frem for perfekte endelser.

For at blive klar havde jeg brugt weekenden på at forberede mig. \da{Ud over at undersøge} hjemmesiden trænede jeg spørgsmål med en dansk veninde, så jeg kendte de vigtigste fagord. Skriv endelig, hvis du vil låne mine noter!

Mange hilsner\\
Klara

\vspace{0.3em}
\rule{\linewidth}{0.4pt}\\[0.2em]
\small\textbf{\color{NCDNavy}[高分亮点 · 138 ord]}：\ding{172} 精准示范介词接间接疑问从句（\textit{spændt på, om mit danske sprog var godt nok}）；\ding{173} 高级短语 \textit{over al forventning}、\textit{lagde vægt på... frem for...} 与 \textit{Ud over at...} 极大提升书面表达成熟度，且字数控制在考场 30 分钟黄金区间。
\end{emailbox}

% =====================================================================
% 4.4 口语示范稿与四图双人对话
% =====================================================================
\subsubsection*{4.4 口语示范逐字稿、核心 IPA 与四图双人互动全真对话}

\begin{transbox}
\small
\textbf{\color{NCDNavy}[Mindmap 3A 口语陈述范文：Min skole / uddannelse｜152 词 · 约 1.5 分钟]}\\
\textit{Jeg vil gerne fortælle lidt om min uddannelse \ipa{ˈuðˌdanˀəlsə}. Jeg studerer multimediedesign på et erhvervsakademi \ipa{ɛɐˈvɛːˀɐsˌɑkʰadəˌmiˀ} i København, hvor jeg går på andet år. Det er en offentlig uddannelse, og der går cirka to hundrede studerende \ipa{sd̥uˈdeːˀɐnə} på min årgang. Jeg møder fra mandag til fredag mellem klokken ni og femten. Hver morgen cykler jeg langs kanalen i tyve minutter, men hvis det regner kraftigt, tager jeg metroen i stedet for. I undervisningen \ipa{ˈʌnɐˌviːˀsneŋ} arbejder vi både med grafisk design, kodning og store gruppeprojekter \ipa{ˈɡʁubəˌpʰʁoˌjɛɡdɐ}. Jeg kan bedst lide de praktiske opgaver, mens det sværeste for mig stadig er at skrive lange rapporter på dansk. Vores lærere er meget uformelle og hjælpsomme, og mit hold består af både danskere og internationale studerende. Jeg mødes fast med en lille studiegruppe \ipa{ˈsd̥uˀdjəˌɡʁubə}, hvor vi støtter hinanden. Alt i alt er jeg rigtig glad for uddannelsen, fordi den er jobrelevant, selvom det bliver ret travlt før eksamen \ipa{ɛɡˈsæːmən}.}
\end{transbox}

\paragraph*{Mindmap 3A Del 2 考官追问实战示范对话（4 题全解）：}
\dialogturn{1}{Eksaminator}{Du siger, at I laver mange gruppeprojekter --- kan du bedst lide at arbejde alene eller i en gruppe? Hvorfor?}{gruppeprojekter \ipa{ˈɡʁubəˌpʰʁoˌjɛɡdɐ} · alene \ipa{aˈleːnə}}{你说你们做很多小组项目——你更喜欢独自工作还是小组合作？为什么？}
\dialogturn{2}{Kursist}{Jeg kan faktisk bedst lide at arbejde i en gruppe, fordi man får flere gode idéer sammen. Når man er usikker på et dansk fagord, kan ens studiekammerater også hjælpe med det samme.}{idéer \ipa{iˈdeːˀɐ} · studiekammerater \ipa{ˈsd̥uˀdjəˌkʰamɑˌʁɑːˀdɐ}}{其实我更喜欢小组合作，因为大家一起能碰撞出更多好点子。当对某个丹麦语专业词拿不准时，同学也能立刻帮忙。}
\dialogturn{3}{Eksaminator}{Du nævner, at skriftlige rapporter på dansk er det sværeste --- hvad gør du for at blive bedre til det?}{skriftlige \ipa{ˈsɡʁefdliə} · rapporter \ipa{ʁaˈpʰʌdɐ}}{你提到用丹麦语写书面报告是最难的——你做些什么来提升呢？}
\dialogturn{4}{Kursist}{Jeg læser mange danske artikler om design, og når vi skriver i gruppen, beder jeg ofte mine danske kammerater om at rette mine sætninger. På den måde lærer jeg hurtigt af mine fejl.}{artikler \ipa{ɑˈtˢiɡlɐ} · fejl \ipa{ˈfɑi̯ˀl}}{我读很多关于设计的丹麦语文章，而且小组写作时常请丹麦同学帮我修改句子。这样我能很快从错误中吸取经验。}
\dialogturn{5}{Eksaminator}{Hvordan er forholdet mellem lærere og studerende i Danmark sammenlignet med dit hjemland?}{forholdet \ipa{ˈfɒːˌhʌlˀəð} · sammenlignet \ipa{ˈsamənˌliˀnəð}}{与你的祖国相比，丹麦师生之间的关系怎么样？}
\dialogturn{6}{Kursist}{Her er tonen meget mere uformel. Vi kalder lærerne ved fornavn og drikker kaffe sammen i pauserne. Det gør, at man tør stille dumme spørgsmål uden at være nervøs.}{uformel \ipa{ˈufʌˌmɛlˀ} · fornavn \ipa{ˈfɒːˌnɑu̯ˀn}}{这里的氛围要非正式得多。我们直呼老师名字，课间还一起喝咖啡，这让人敢于大胆提问而不觉得紧张。}
\dialogturn{7}{Eksaminator}{Hvilken slags job drømmer du om at få, når du er færdig med uddannelsen?}{drømmer \ipa{ˈdʁœmɐ} · færdig \ipa{ˈfɛːˀdi}}{毕业后你梦想找一份什么样的工作？}
\dialogturn{8}{Kursist}{Jeg håber at få et job i et mellemstort designbureau, hvor jeg både kan lave hjemmesider og brugervenlige apps --- og helst på en arbejdsplads, hvor vi taler dansk til hverdag.}{designbureau \ipa{diˈsɑi̯nbyˌʁo} · brugervenlige \ipa{ˈbʁuːɐˌvɛnliə}}{我希望在一家中型设计事务所工作，既做网站也做用户友好的应用——最好是在日常讲丹麦语的职场。}

\begin{transbox}
\small
\textbf{\color{NCDNavy}[Mindmap 3B 口语陈述范文：Mit frivillige arbejde｜157 词 · 约 1.5 分钟]}\\
\textit{Nu vil jeg gerne fortælle om mit frivillige arbejde \ipa{ˈfʁiˌviliə ˈɑːbɑi̯ˀdə}. Hver anden torsdag aften hjælper jeg til i et lokalt folkekøkken, som hedder Naboens Bord, og det ligger i et kulturhus i midtbyen. Jeg er der fra klokken halv fem til halv ni, så jeg bruger cirka otte timer om måneden på det. Mine arbejdsopgaver \ipa{ˈɑːbɑi̯ˀdsˌʌbˌɡæːvɐ} er ret varierede: Først skærer jeg grøntsager og hjælper kokken, og bagefter dækker jeg langborde og tager imod gæsterne ved døren. Til sidst rydder vi op i køkkenet sammen. Vi er et fast hold på ti frivillige i alle aldre --- både naboer og unge studerende --- og der er altid en rigtig god holdånd \ipa{ˈhʌlˌʌnˀ}. Det har givet mig et stort udbytte \ipa{ˈuðˌbydə}, for jeg øver hverdagssnak på dansk, har fået nye venner og lært meget om dansk madkultur. Selvom man bliver lidt træt i benene af at stå op længe, synes jeg, at arbejdet er utrolig meningsfuldt \ipa{ˈmeːneŋsˌfult} og giver mig ny energi.}
\end{transbox}

\paragraph*{Mindmap 3B Del 2 考官追问实战示范对话（4 题全解）：}
\dialogturn{1}{Eksaminator}{Hvordan fandt du frem til dette folkekøkken, og hvorfor valgte du at blive frivillig?}{folkekøkken \ipa{ˈfʌlɡəˌkʰøɡən} · valgte \ipa{ˈvalˀdə}}{你是怎么找到这家社区食堂的？为什么选择当志愿者？}
\dialogturn{2}{Kursist}{En kollega fortalte mig om stedet, og jeg havde lyst til at gøre noget socialt uden for mit arbejde. Mad samler jo folk, så det føltes som en helt naturlig måde at møde lokale danskere på.}{socialt \ipa{soˈɕæːˀld} · samler \ipa{ˈsamlɐ}}{一位同事告诉我的，我也想在工作之余做点公益社交。美食总能把大家聚在一起，所以感觉这是认识本地丹麦人非常自然的方式。}
\dialogturn{3}{Eksaminator}{Hvad kan du bedst lide af dine opgaver i køkkenet eller ved bordene?}{opgaver \ipa{ˈʌbˌɡæːvɐ} · bordene \ipa{ˈboːˀɐnə}}{在厨房或餐桌旁的各项任务中，你最喜欢哪一项？}
\dialogturn{4}{Kursist}{Jeg elsker mest at tage imod gæsterne og servere maden, fordi man så får mange små samtaler i gang. Det er lidt mindre sjovt at vaske de store gryder op til sidst, men når vi gør det sammen med musik, går det hurtigt.}{servere \ipa{sɛɐˈveːˀɐ} · gryder \ipa{ˈɡʁyːðɐ}}{我最喜欢迎接客人和上菜，因为那样能和不少人聊起来。最后刷大锅没那么有趣，不过大家放着音乐一起干就很快。}
\dialogturn{5}{Eksaminator}{Hvad har det frivillige arbejde betydet for dit danske sprog og dit netværk?}{betydet \ipa{beˈtˢyːˀðəð} · netværk \ipa{ˈnɛdˌvɛɐɡ}}{这段志愿经历对你的丹麦语和人脉圈有什么意义？}
\dialogturn{6}{Kursist}{Det har betydet rigtig meget! I begyndelsen var jeg genert, men nu tør jeg tale dansk om alt muligt ved bordene. Og to af de andre frivillige inviterede mig faktisk til fødselsdag i sidste måned.}{genert \ipa{ɕeˈneːˀɐd} · inviterede \ipa{enˈviˌtˢeːˀɐðə}}{意义非常重大！起初我很害羞，但现在在餐桌旁什么话题都敢用丹麦语聊。上个月还有两位志愿者朋友邀请我参加了生日聚会。}
\dialogturn{7}{Eksaminator}{Synes du, at det er let eller svært at passe frivilligt arbejde ved siden af hverdagen?}{passe \ipa{ˈpʰasə} · hverdagen \ipa{ˈvɛːˀˌdæːˀən}}{你觉得在日常工作生活之余兼顾志愿服务容易还是困难？}
\dialogturn{8}{Kursist}{Det er faktisk ret nemt, når det kun er hver anden torsdag. Fordi vi planlægger vagterne en måned i forvejen, kan man sagtens få det til at passe med familie og arbejde.}{vagterne \ipa{ˈvaɡdɐnə} · i forvejen \ipa{i ˈfɒːˌvɑi̯ˀən}}{隔周周四去一次其实挺容易的。因为我们提前一个月排好班次，完全能和家庭及工作协调好。}

\paragraph*{Mundtlig Opgave 2（EMNE: Indkøb og forbrug）Del 1 双人四图全真互动示范：}

\dialogturn{1}{Kursist A}{Skal vi kigge på det første billede med supermarkedet? Jeg handler selv i supermarked tre gange om ugen, fordi det ligger lige om hjørnet og har alle dagligvarer på ét sted. Hvad synes du om at købe ind i supermarkedet?}{supermarked \ipa{ˈsuːˀpʰɐˌmɑːɡəð} · dagligvarer \ipa{ˈdɑuliˌvɑːɐ}}{我们看看第一张超市图吧？我自己每周去超市买三次，因为它就在拐角而且一站买齐日用品。你觉得去超市买菜怎么样？}
\dialogturn{2}{Kursist B}{Jeg kan egentlig også godt lide det til hverdag, især når man vil vælge frisk frugt og kød selv. Men jeg synes tit, der er lange køer klokken fem om eftermiddagen. Foretrækker du i stedet at handle på nettet, ligesom på billede to?}{foretrækker \ipa{ˈfɒːˌtˢʁɛɡɐ} · nettet \ipa{ˈnɛdəð}}{平时我也挺喜欢的，尤其是想自己挑新鲜水果和肉的时候。不过下午五点经常排长队。你会更喜欢像图二那样网购吗？}
\dialogturn{3}{Kursist A}{Ja, nethandel bruger jeg rigtig meget til tunge ting som havregryn og vaskepulver, og også til tøj, fordi det sparer en masse tid. Til gengæld kan det være irriterende at sende tøj tilbage, hvis størrelsen ikke passer. Handler du ofte online?}{nethandel \ipa{ˈnɛdˌhanˀəl} · størrelsen \ipa{ˈsd̥œːɐlsən}}{是的，买燕麦片、洗衣粉等重物或衣服时我常用网购，因为非常省时间。但衣服尺码不合适还得退货挺烦的。你经常网购吗？}
\dialogturn{4}{Kursist B}{Kun nogle gange til elektronik eller bøger, for jeg bliver faktisk lidt træt af alle de papkasser. Men jeg elsker billede tre med loppemarked og genbrugsbutik! Jeg køber næsten alle mine tallerkener og lamper brugt, fordi det både er billigt og godt for klimaet. Hvad mener du om genbrug?}{loppemarked \ipa{ˈlʌbəˌmɑːɡəð} · klimaet \ipa{ˈkʰliːˌmæːˀəð}}{我只偶尔网购电子产品或书，因为收一堆纸箱挺烦的。但我超爱图三的跳蚤市场和二手店！我几乎所有盘子和台灯都买二手的，既便宜又环保。你怎么看二手购物？}
\dialogturn{5}{Kursist A}{Jeg er enig i, at genbrug er super til møbler og bøger, selvom det tager lidt længere tid at finde den rigtige ting. Hvad så med det sidste billede fra en gårdbutik eller et marked? Har du prøvet det?}{enig \ipa{ˈeːni} · gårdbutik \ipa{ˈɡɔːbuˌtˢiɡ}}{我同意买二手家具和书很棒，尽管挑到合适物件要多花点时间。那最后一张农场直营店或集市图呢？你去过吗？}
\dialogturn{6}{Kursist B}{Ja, om sommeren cykler vi tit forbi en lille gårdbutik på landet. Grøntsagerne smager meget bedre end i supermarkedet, og man støtter lokale landmænd. Dog er priserne lidt højere, så det er mest til weekenden. Er du også villig til at betale ekstra for friske varer?}{landmænd \ipa{ˈlanˌmɛnˀ} · villig \ipa{ˈvili}}{去过，夏天我们常骑车路过乡下小农场商店。蔬菜比超市好吃得多，还能支持本地农户。不过价格略贵，所以主要周末去。你也愿意为新鲜食材多花钱吗？}
\dialogturn{7}{Kursist A}{Ja, til gæster og i weekenden gør jeg gerne det, fordi kvaliteten bare er i top. Men til daglig er supermarkedet og et godt tilbud på nettet stadig det mest praktiske for mig.}{kvaliteten \ipa{kʰvaliˈtˢeːˀdən} · praktiske \ipa{ˈpʰʁɑɡtˢisɡə}}{是的，待客或周末我愿意，因为品质真是一流。但工作日还是超市和划算的网购最方便我。}

\paragraph*{Mundtlig Opgave 2 Del 2（考官个人深入追问 3 题高分示范）：}
\dialogturn{1}{Eksaminator (Q1)}{Hvordan kan du bedst lide at købe ind --- og hvorfor?}{hvordan \ipa{vʌˈdan} · bedst \ipa{ˈbɛsd}}{你最喜欢以哪种方式购物——为什么？}
\dialogturn{2}{Kursist}{Jeg kan bedst lide en blanding af supermarked i hverdagene og gårdbutik i weekenden. I supermarkedet går det hurtigt efter arbejde, mens jeg på et lokalt marked virkelig nyder den gode kvalitet og atmosfæren.}{blanding \ipa{ˈblaneŋ} · atmosfæren \ipa{admosˈfɛːˀɐn}}{我最喜欢工作日逛超市、周末去农场店结合。下班后在超市采购很快捷，而周末在本地集市能好好享受高品质食材和热闹氛围。}
\dialogturn{3}{Eksaminator (Q2)}{Hvad er vigtigst for dig, når du handler i hverdagen --- pris, kvalitet, tid, klima eller økologi?}{økologi \ipa{økʰoloˈɡiˀ} · vigtigst \ipa{ˈveɡtˢisd}}{平时买东西时对你最重要的是什么——价格、质量、时间、气候环保还是有机？}
\dialogturn{4}{Kursist}{Tid og kvalitet betyder mest for mig, men jeg tænker også meget på økologi og madspild. Derfor planlægger jeg ofte en madplan for hele ugen og køber danske, økologiske grøntsager, når de er i sæson.}{madspild \ipa{ˈmaðˌsbilˀ} · sæson \ipa{sɛˈsʌŋ}}{时间和质量对我最重要，但我也非常重视有机和减少食物浪费。因此我常制定一周食谱，专门买当季的丹麦有机蔬菜。}
\dialogturn{5}{Eksaminator (Q3)}{Kan du fortælle lidt om et godt eller dårligt køb, du har lavet for nylig?}{for nylig \ipa{fʌ ˈnyːli} · køb \ipa{ˈkʰøːˀb}}{能讲讲你最近一次很满意或很糟糕的购物经历吗？}
\dialogturn{6}{Kursist}{For to uger siden gjorde jeg et rigtig godt køb på et loppemarked i Aarhus. Jeg fandt et næsten nyt spisebord i massivt egetræ til kun fem hundrede kroner. Sælgeren hjalp endda med at bære det ned til min vens bil, så jeg sparede flere tusinde kroner.}{massivt \ipa{maˈsiu̯ˀd} · egetræ \ipa{ˈɑi̯əˌtˢʁæˀ}}{两周前我在奥胡斯跳蚤市场捡了个大漏。我只花五百克朗买到一张几乎全新的实木橡木餐桌。卖家还帮忙抬到我朋友车上，让我省下了好几千克朗。}

\ncdmodultestmthree{4}{Familie, By og Fremtid}{全真模考卷四：城乡搬迁、客运投诉与家庭分工}{城乡搬迁选词完形 / 周末改期与子女病假双聊天 / 首次企业实习补句 / 三人谈家务与家庭分工 / 客运大巴投诉与丹麦适应回信 / 家庭分工与丹麦语学习 Mindmap 及假期旅游四图双人互动}{B1 全核心语法总复习：-s 被动与 blive 被动辨析 · 介词 + 形式宾语 det + at/hv- 从句 · 分词作定语 · 无 hvis 条件倒装 · 让步对比 til trods for / selvom / mens... derimod · 篇章连接词 For / Men / Derfor / Desuden}

% =====================================================================
% Del 1: Læsning (50 minutter · 24 point)
% =====================================================================
\ncdsubsec{Del 1: Læsning（阅读理解 · 50 分钟 · 24 分）}

\begin{tipbox}{考场规则与配速提醒（Læsning Opgave 1, 2A, 2B, 3 og 4）}
Opgave 1, 2A, 2B, 3 og 4 udleveres samtidig. Samlet tid: \textbf{50 minutter}. Du må \textbf{ikke} bruge ordbog, noter eller mobiltelefon. Der gives 1 point for hvert korrekt svar; maksimum er \textbf{24 point}（及格线：\textbf{15--24 point = Bestået}；\textbf{0--14 point = Ikke bestået}）。建议时间分配：Opgave 1 (8 min) $\rightarrow$ Opgave 2A \& 2B (10 min) $\rightarrow$ Opgave 3 (14 min) $\rightarrow$ Opgave 4 (13 min)，最后留 5 分钟复查核对。
\end{tipbox}

% ---------------------------------------------------------------------
\subsubsection*{Opgave 1 (6 point): \textit{Fra storbyen til en lille landsby}}

\begin{emailbox}{Kursistark --- Instruktion（Læsning, Opgave 1）}
\begin{itemize}
  \item Læs teksten \textit{Fra storbyen til en lille landsby}.
  \item Skriv de ord (1--6), der mangler. Du skal kun bruge hvert ord \textbf{én gang}.
  \item Der er \textbf{seks ord, du ikke skal bruge}. Se eksemplet (0).
\end{itemize}

\vspace{0.35em}
\hrule
\vspace{0.45em}

\textbf{\color{NCDNavy}Fra storbyen til en lille landsby}

\vspace{0.3em}
Efter syv år i en dyr toværelses lejlighed midt i København besluttede Signe og Kasper, at de ville flytte ud på landet. De havde længe drømt om mere plads og en have, men det var svært at finde et hus til en \textbf{(0)} \underline{\textbf{rimelig}} pris i hovedstaden.

\vspace{0.35em}
En weekend i foråret besøgte de nogle venner i en lille landsby på Fyn, \textbf{(1)} \rule{2.0cm}{0.4pt} der kun bor omkring fire hundrede mennesker. De faldt straks for den rolige stemning og de grønne marker omkring byen. Selvom de var lidt bekymrede for den lange transporttid til arbejde, købte de et gammelt murstenshus med frugttræer i haven.

\vspace{0.35em}
Den første måned var Signe lidt usikker på, \textbf{(2)} \rule{2.0cm}{0.4pt} hun ville savne caféerne og kulturlivet i København. Om aftenen var gaderne helt stille, og i haven kunne de høre fuglene synge, \textbf{(3)} \rule{2.0cm}{0.4pt} de spiste morgenmad sammen. Signe havde egentlig troet, at folk på landet holdt sig lidt for sig selv, \textbf{(4)} \rule{2.0cm}{0.4pt} allerede den første weekend bankede naboerne på døren med friskbagt brød.

\vspace{0.35em}
I dag pendler Kasper med tog fire dage om ugen, og han har endelig vænnet sig \textbf{(5)} \rule{2.0cm}{0.4pt} den ekstra rejsetid, fordi han kan arbejde hjemmefra hver fredag. Parret er stadig meget glade for deres beslutning, \textbf{(6)} \rule{2.0cm}{0.4pt} turen til kontoret faktisk tager længere tid end før. Livet på landet har givet dem meget mere ro i hverdagen.

\vspace{0.55em}
\begin{center}
\fbox{\parbox{0.93\linewidth}{\centering\vspace{4pt}
\textbf{til} \quad\textbar\quad \textbf{hvis} \quad\textbar\quad \textbf{hvor} \quad\textbar\quad \textbf{ikke} \quad\textbar\quad \textbf{mens} \quad\textbar\quad \textbf{rimelig} \quad\textbar\quad \textbf{under}\\[3pt]
\textbf{om} \quad\textbar\quad \textbf{fordi} \quad\textbar\quad \textbf{selvom} \quad\textbar\quad \textbf{lige} \quad\textbar\quad \textbf{men} \quad\textbar\quad \textbf{så}%
\vspace{4pt}}}
\end{center}
\end{emailbox}

% ---------------------------------------------------------------------
\subsubsection*{Opgave 2A (NY · 3 point): \textit{Mette skriver til sin veninde Louise}}

\begin{emailbox}{Kursistark --- Instruktion（Læsning, Opgave 2A）}
\textbf{Situation:} Mette skriver til sin veninde Louise.
\begin{itemize}
  \item Læs Mettes og Louises chat. Der er 5 beskeder fra Mette.
  \item Find de svar fra Louise, der passer (A--F).
  \item Skriv bogstavet på det svar, der passer. Der er \textbf{tre svar, du ikke skal bruge}. Se eksemplet (0).
\end{itemize}
\end{emailbox}

\vspace{0.25em}
{\small
\renewcommand{\arraystretch}{1.35}
\begin{longtable}{@{} >{\raggedright\arraybackslash}p{8.25cm} @{\hspace{0.5cm}} >{\raggedright\arraybackslash}p{8.25cm} @{}}
\toprule
\textbf{\color{NCDNavy}Chat mellem Mette og Louise} & \textbf{\color{NCDNavy}Louises mulige svar (Z, A--F)} \\
\midrule
\endfirsthead
\toprule
\textbf{\color{NCDNavy}Chat mellem Mette og Louise (fortsat)} & \textbf{\color{NCDNavy}Louises mulige svar (fortsat)} \\
\midrule
\endhead
\bottomrule
\endfoot

\textbf{\color{NCDNavy}Mette:}\newline
Hej Louise! Glæder du dig til vores tur til Aarhus på lørdag med biograf og kunstmuseum?\par\vspace{4pt}
\textbf{0.} \textbf{\color{NCDTeal}Louise:} \textcolor{NCDRed}{\textbf{[ Eksempel: Z ]}} &
\textbf{\color{NCDRed}Z:}\ Ja, helt vildt! Jeg har set frem til det hele ugen. \\
\midrule

\textbf{\color{NCDNavy}Mette:}\newline
Også mig, men jeg har lige set på Rejseplanen, at der er sporarbejde på strækningen hele lørdagen. Så toget kører slet ikke.\par\vspace{4pt}
\textbf{1.} \textbf{\color{NCDTeal}Louise:} \rule{1.2cm}{0.4pt} &
\textbf{\color{NCDNavy}A:}\ Det lyder godt! Jeg elsker nemlig togbussen, for så kan man sove lidt ekstra undervejs.\par\vspace{6pt}
\textbf{\color{NCDNavy}B:}\ Søndag passer mig fint! Men jeg har desværre ikke tid til både filmen og udstillingen, fordi jeg skal besøge min bror sent om eftermiddagen. Skal vi droppe kunstmuseet? \\
\midrule

\textbf{\color{NCDNavy}Mette:}\newline
Det har du desværre ret i --- det tager næsten dobbelt så lang tid med togbus. Hvad hvis vi tager turen om søndagen i stedet?\par\vspace{4pt}
\textbf{2.} \textbf{\color{NCDTeal}Louise:} \rule{1.2cm}{0.4pt} &
\textbf{\color{NCDNavy}C:}\ Desværre, på søndag arbejder jeg hele dagen fra morgen til aften, så vi må vente til næste måned.\par\vspace{6pt}
\textbf{\color{NCDNavy}D:}\ Øv, hvor ærgerligt! Så skal vi vel med togbus, og det tager jo en evighed. \\
\midrule

\textbf{\color{NCDNavy}Mette:}\newline
Nej, så går vi jo glip af den nye udstilling, og det vil jeg nødigt. Kunne vi så ikke se kunstmuseet om formiddagen og droppe biografen?\par\vspace{4pt}
\textbf{3.} \textbf{\color{NCDTeal}Louise:} \rule{1.2cm}{0.4pt} &
\textbf{\color{NCDNavy}E:}\ Jeg tager gerne begge biografbilletter med på lørdag, så vi ikke kommer for sent til filmen.\par\vspace{6pt}
\textbf{\color{NCDNavy}F:}\ Jo, det lyder som en rigtig god plan! Så kan vi også nå en kop kaffe bagefter. Bestiller du billetter til museet? \\
\midrule

\textbf{\color{NCDNavy}Mette:}\newline
Super, så køber jeg museumsbilletter til os til kl. 10.30. Vi ses på stationen søndag morgen! &
\textit{\color{NCDGray}(Husk: Der er tre svar blandt A--F, som du ikke skal bruge.)} \\
\end{longtable}
}

% ---------------------------------------------------------------------
\subsubsection*{Opgave 2B (NY · 3 point): \textit{Anders skriver til sin chef Camilla}}

\begin{emailbox}{Kursistark --- Instruktion（Læsning, Opgave 2B）}
\textbf{Situation:} Anders skriver til sin chef Camilla.
\begin{itemize}
  \item Læs Anders' og Camillas chat. Der er 5 beskeder fra Anders.
  \item Find de svar fra Camilla, der passer (A--F).
  \item Skriv bogstavet på det svar, der passer. Der er \textbf{tre svar, du ikke skal bruge}. Se eksemplet (0).
\end{itemize}
\end{emailbox}

\vspace{0.25em}
{\small
\renewcommand{\arraystretch}{1.35}
\begin{longtable}{@{} >{\raggedright\arraybackslash}p{8.25cm} @{\hspace{0.5cm}} >{\raggedright\arraybackslash}p{8.25cm} @{}}
\toprule
\textbf{\color{NCDNavy}Chat mellem Anders og Camilla} & \textbf{\color{NCDNavy}Camillas mulige svar (Z, A--F)} \\
\midrule
\endfirsthead
\toprule
\textbf{\color{NCDNavy}Chat mellem Anders og Camilla (fortsat)} & \textbf{\color{NCDNavy}Camillas mulige svar (fortsat)} \\
\midrule
\endhead
\bottomrule
\endfoot

\textbf{\color{NCDNavy}Anders:}\newline
Godmorgen Camilla. Min søn Noah er desværre vågnet med høj feber, så jeg bliver nødt til at tage barns første sygedag i dag.\par\vspace{4pt}
\textbf{0.} \textbf{\color{NCDTeal}Camilla:} \textcolor{NCDRed}{\textbf{[ Eksempel: Z ]}} &
\textbf{\color{NCDRed}Z:}\ Godmorgen Anders. Det er jeg ked af at høre --- god bedring til ham! Kan du overhovedet arbejde lidt hjemmefra? \\
\midrule

\textbf{\color{NCDNavy}Anders:}\newline
Ja, min kone kommer hjem ved frokosttid og tager over, så jeg kan godt arbejde hjemmefra i eftermiddag. Men hvad med mit kundemøde kl. 9.30?\par\vspace{4pt}
\textbf{1.} \textbf{\color{NCDTeal}Camilla:} \rule{1.2cm}{0.4pt} &
\textbf{\color{NCDNavy}A:}\ Perfekt! Husk også lige at registrere barns første sygedag i vores system, ikke?\par\vspace{6pt}
\textbf{\color{NCDNavy}B:}\ Nej, Martin har desværre ferie hele denne uge, så du må selv ringe og aflyse mødet med kunden. \\
\midrule

\textbf{\color{NCDNavy}Anders:}\newline
Tusind tak, det er en kæmpe lettelse! Jeg sender ham straks alle noterne og præsentationen på mail.\par\vspace{4pt}
\textbf{2.} \textbf{\color{NCDTeal}Camilla:} \rule{1.2cm}{0.4pt} &
\textbf{\color{NCDNavy}C:}\ Ja, hvis du kan nå at gøre månedsbudgettet færdigt inden fyraften, ville det være en stor hjælp.\par\vspace{6pt}
\textbf{\color{NCDNavy}D:}\ Så synes jeg, du bare skal holde helt fri resten af dagen og først komme på arbejde igen på mandag. \\
\midrule

\textbf{\color{NCDNavy}Anders:}\newline
Ja, jeg har allerede registreret det i vores tidsystem, og jeg logger på Teams fra kl. 13. Er der noget særligt, jeg skal prioritere i eftermiddag?\par\vspace{4pt}
\textbf{3.} \textbf{\color{NCDTeal}Camilla:} \rule{1.2cm}{0.4pt} &
\textbf{\color{NCDNavy}E:}\ Bare rolig, jeg spørger Martin, om han kan tage kundemødet for dig. Kan du sende ham dit materiale?\par\vspace{6pt}
\textbf{\color{NCDNavy}F:}\ Det behøver du ikke, for jeg har allerede ringet til din kone og aftalt, at hun bliver hjemme hele dagen. \\
\midrule

\textbf{\color{NCDNavy}Anders:}\newline
Det skal jeg nok gøre færdigt før kl. 16. Vi snakkes ved på Teams i eftermiddag! &
\textit{\color{NCDGray}(Husk: Der er tre svar blandt A--F, som du ikke skal bruge.)} \\
\end{longtable}
}

% ---------------------------------------------------------------------
\subsubsection*{Opgave 3 (NY · 5 point): \textit{Fatimas første praktikplads}}

\begin{emailbox}{Kursistark --- Instruktion（Læsning, Opgave 3）}
\begin{itemize}
  \item Læs teksten \textit{Fatimas første praktikplads}. Der er fem afsnit.
  \item I hvert afsnit mangler der en sætning. Under hvert afsnit står der fire sætninger.
  \item Skriv bogstavet på den sætning, der passer i afsnittet. Der er \textbf{tre sætninger, du ikke skal bruge} i hvert afsnit.
\end{itemize}

\vspace{0.35em}
\hrule
\vspace{0.45em}

\textbf{1.} Fatima har boet i Odense i to år og har længe drømt om at arbejde med bogholderi og økonomi i Danmark. I sidste måned bestod hun sin danskprøve og fik endelig tilbudt en otte ugers praktikplads på et lokalt regnskabskontor. \rule{3.2cm}{0.4pt}. For selvom hun har arbejdet på kontor i sit hjemland, er alle danske regler og ord helt nye for hende.
\begin{enumerate}[label=\textbf{\Alph*.}, leftmargin=1.8em, itemsep=0.12em, topsep=0.25em]
  \item Hun er både meget glad og lidt nervøs før sin første dag.
  \item Hun har derfor besluttet at sige nej tak til praktikpladsen.
  \item Derfor kender hun allerede alle de danske skatteregler rigtig godt.
  \item Men hendes chef vil helst ikke have praktikanter på kontoret.
\end{enumerate}

\vspace{0.4em}
\textbf{2.} Mandag morgen bliver Fatima budt velkommen af sin kontaktperson, Mette, som viser hende rundt på kontoret og præsenterer hende for de andre ansatte. De første par dage skal Fatima læse gamle fakturaer og lære virksomhedens digitale system at kende. \rule{3.2cm}{0.4pt}. Heldigvis siger Mette, at hun altid må spørge, hvis der er et fagord eller en opgave, hun ikke forstår.
\begin{enumerate}[label=\textbf{\Alph*.}, leftmargin=1.8em, itemsep=0.12em, topsep=0.25em]
  \item Det er helt nemt for hende, fordi hun bruger systemet hjemmefra hver dag.
  \item Det tager lidt længere tid, end hun havde regnet med, fordi mange ord er svære.
  \item Derfor beder Mette hende om at arbejde helt alene uden at tale med nogen.
  \item Så Fatima bliver sendt hjem allerede før frokost den første dag.
\end{enumerate}

\vspace{0.4em}
\textbf{3.} Efter to uger har Fatima fået mere styr på computerprogrammet, men pauserne i kantinen føles stadig lidt svære. Når de otte kollegaer sidder sammen og spiser frokost, taler de nemlig hurtigt og griner af indforståede jokes. \rule{3.2cm}{0.4pt}. Derfor beslutter hun sig for at sætte sig ved siden af en venlig kollega og selv stille et par spørgsmål om weekendens planer.
\begin{enumerate}[label=\textbf{\Alph*.}, leftmargin=1.8em, itemsep=0.12em, topsep=0.25em]
  \item Fatima foretrækker faktisk at spise sin frokost alene ved sit skrivebord hver dag.
  \item Kollegaerne taler kun engelsk i kantinen, så Fatima kan forstå det hele.
  \item Det gør hende nogle gange usikker på, hvornår hun selv skal sige noget.
  \item Men det er kun, fordi hun ikke kan lide maden i kantinen.
\end{enumerate}

\vspace{0.4em}
\textbf{4.} Midt i praktikperioden får Fatima lov til selv at kontrollere månedens kvitteringer og taste tal ind i regnskabet. En eftermiddag opdager hun pludselig, at der mangler et stort beløb i en af mapperne. \rule{3.2cm}{0.4pt}. Mette bliver meget imponeret over, at Fatima har opdaget fejlen, og roser hende for at være så omhyggelig.
\begin{enumerate}[label=\textbf{\Alph*.}, leftmargin=1.8em, itemsep=0.12em, topsep=0.25em]
  \item Hun tænker, at det nok ikke er så vigtigt, og lægger bare mappen væk.
  \item Fatima bliver bange for, at hun selv har ødelagt det hele, og siger ingenting.
  \item Efter en halv times søgen finder hun heldigvis det forsvundne bilag og viser det til Mette.
  \item Derfor ringer hun straks til en anden virksomhed for at søge et nyt job.
\end{enumerate}

\vspace{0.4em}
\textbf{5.} Da de otte uger næsten er gået, bliver Fatima kaldt ind til en samtale på chefens kontor. Hun er spændt på, hvad chefen vil sige om hendes indsats i løbet af praktikken. \rule{3.2cm}{0.4pt}. Fatima går hjem med et kæmpe smil, for nu ved hun, at drømmen om et fast arbejde i Danmark faktisk kan blive til virkelighed.
\begin{enumerate}[label=\textbf{\Alph*.}, leftmargin=1.8em, itemsep=0.12em, topsep=0.25em]
  \item Chefen fortæller hende desværre, at kontoret ikke kan bruge hendes hjælp længere.
  \item Chefen tilbyder hende et vikariat på deltid fra næste måned, fordi teamet er så glade for hende.
  \item De bliver enige om, at Fatima slet ikke passer til at arbejde med regnskab.
  \item Men chefen har glemt, at Fatima overhovedet har været i praktik hos dem.
\end{enumerate}
\end{emailbox}

% ---------------------------------------------------------------------
\subsubsection*{Opgave 4 (NY · 7 point): \textit{Tre personer fortæller om arbejdsdeling i hjemmet}}

\begin{emailbox}{Kursistark --- Instruktion \& Svarskema（Læsning, Opgave 4）}
\begin{itemize}
  \item Læs de tre tekster, og læs spørgsmålene (1--7).
  \item Find den person, der passer til hvert af de syv spørgsmål.
  \item Sæt kryds (\textbf{X}) ud for den person (A, B eller C), der passer til spørgsmålet. Se eksemplet (0).
\end{itemize}

\vspace{0.3em}
{\small
\renewcommand{\arraystretch}{1.3}
\begin{tabularx}{\linewidth}{@{} >{\bfseries}p{0.6cm} >{\raggedright\arraybackslash}X >{\centering\arraybackslash}p{1.65cm} >{\centering\arraybackslash}p{1.65cm} >{\centering\arraybackslash}p{1.65cm} @{}}
\toprule
 & \textbf{Spørgsmål} & \textbf{Stine} & \textbf{Karim} & \textbf{Elena} \\
\midrule
0. & Hvem mener, at det er godt at fordele opgaverne efter evner og interesser? & A \underline{\hspace{0.45cm}} & B \underline{\hspace{0.45cm}} & C \underline{\textbf{X}} \\
\midrule
1. & Hvem får professionel hjælp udefra til rengøringen? & A \underline{\hspace{0.45cm}} & B \underline{\hspace{0.45cm}} & C \underline{\hspace{0.45cm}} \\
2. & Hvem bruger en fast skriftlig plan i køkkenet til at fordele pligterne? & A \underline{\hspace{0.45cm}} & B \underline{\hspace{0.45cm}} & C \underline{\hspace{0.45cm}} \\
3. & Hvem har overtaget al madlavningen, fordi partneren ikke var god til det? & A \underline{\hspace{0.45cm}} & B \underline{\hspace{0.45cm}} & C \underline{\hspace{0.45cm}} \\
4. & Hvem må ofte bede sine børn flere gange om at gøre deres pligter? & A \underline{\hspace{0.45cm}} & B \underline{\hspace{0.45cm}} & C \underline{\hspace{0.45cm}} \\
5. & Hvem synes, det er irriterende, at hjemmearbejde fører til ekstra husarbejde? & A \underline{\hspace{0.45cm}} & B \underline{\hspace{0.45cm}} & C \underline{\hspace{0.45cm}} \\
6. & Hvem køber ind til familien, mens partneren laver det meste af aftensmaden? & A \underline{\hspace{0.45cm}} & B \underline{\hspace{0.45cm}} & C \underline{\hspace{0.45cm}} \\
7. & Hvem overtog madlavningen i hverdagene, fordi partneren fik skiftende arbejdstider? & A \underline{\hspace{0.45cm}} & B \underline{\hspace{0.45cm}} & C \underline{\hspace{0.45cm}} \\
\bottomrule
\end{tabularx}
}
\end{emailbox}

\begin{emailbox}{Tre personer fortæller om arbejdsdeling i hjemmet (A--C)}
\textbf{\color{NCDNavy}A. Stine}\\
Da vores to børn var helt små, arbejdede jeg på deltid, så jeg lavede næsten alt husarbejde selv. Det førte dog til mange skænderier mellem min mand Nikolaj og mig, fordi jeg altid var træt om aftenen. I dag arbejder vi begge to fuldtid, og derfor har vi lavet et fast ugeskema på køleskabet. Nikolaj støvsuger og gør badeværelset rent, mens jeg tager mig af tøjvasken og indkøbene. Til gengæld elsker Nikolaj at stå i køkkenet, så han laver aftensmad fem dage om ugen, og jeg laver bare noget hurtigt i weekenden. Vores børn på seks og ni år har også faste pligter: De skal selv rydde op på deres værelser og tømme opvaskemaskinen et par gange om ugen. Jeg synes ikke, at ligestilling i hjemmet handler om at tælle hvert eneste minut, men om at begge forældre får lov til at slappe af om aftenen.

\vspace{0.45em}
\textbf{\color{NCDNavy}B. Karim}\\
I min barndom lavede min mor al mad og rengøring, så da jeg blev gift med Sarah, regnede jeg ærligt talt med, at det ville blive på samme måde. Det ændrede sig hurtigt, da Sarah begyndte som sygeplejerske med skiftende nattevagter. I dag handler jeg næsten altid ind og laver varm aftensmad alle hverdage, fordi hun ofte er helt udkørt efter arbejde. Ingen af os har lyst til at bruge vores frie weekender på at skrubbe gulve, så vi har valgt at betale et rengøringsfirma, der kommer hver anden fredag. Vores tre teenagere gider desværre ikke at lave mad, men til gengæld skiftes de til at slå græs, bære skrald ud og lægge tøj sammen --- selvom jeg stadig må minde dem om det flere gange. Efter min mening skal mænd tage lige så meget ansvar i familien som kvinder.

\vspace{0.45em}
\textbf{\color{NCDNavy}C. Elena}\\
Jeg bor sammen med min danske kæreste Frederik, og vi har ingen børn endnu. Da jeg arbejder hjemmefra tre dage om ugen, kommer jeg nogle gange automatisk til at sætte en vask over i frokostpausen, mens Frederik pendler to timer hver dag. Det irriterer mig faktisk lidt, for hjemmearbejde er jo stadig rigtigt arbejde! Tidligere prøvede vi at skiftes til at lave mad, men Frederik brændte maden på to gange og hader at følge opskrifter. Nu står jeg for al madlavningen, fordi jeg godt kan lide det. Til gengæld tager Frederik altid opvasken bagefter, pudser alle vinduer og kører affald på genbrugspladsen. Jeg tror, at god ligestilling handler om at dele opgaverne efter, hvad man hver især er bedst til --- ikke at begge skal lave præcis de samme ting. Hvis vi får børn en dag, skal de tidligt lære at hjælpe til.
\end{emailbox}

% =====================================================================
% Del 2: Skrivning (双任务：Opgave 1 投诉信 + Opgave 2 私人回信 >= 90 ord)
% =====================================================================
\clearpage
\ncdsubsec{Del 2: Skrivning（书面表达 · 双任务实战）}

\subsubsection*{Opgave 1（正式/半正式应用文 · 投诉信 Klage）}

\begin{emailbox}{Kursistark --- Skrivning, Opgave 1}
\textbf{Situation:} Du tog i går med langdistancebussen \textit{NordBus} fra København til Aalborg, men rejsen var meget dårlig. Bussen var kraftigt forsinket uden besked, dit sæde var ødelagt, og der var intet fungerende toilet. Da du klagede, var servicen også dårlig.\\
\textbf{Opgave:} Skriv en klage til \textit{NordBus}. Du skal fortælle:
\begin{itemize}[itemsep=0.05em, topsep=0.08em, parsep=0pt]
  \item hvor og hvornår du rejste
  \item hvad der gik galt under turen
  \item hvordan chaufføren eller kundeservice reagerede
  \item hvordan \textit{NordBus} kan kompensere dig.
\end{itemize}

\vspace{0.15em}
\hrule
\vspace{0.15em}
\textbf{Klage til NordBus}\\[2pt]
\rule{\linewidth}{0.4pt}\\[2pt]
\rule{\linewidth}{0.4pt}\\[2pt]
\rule{\linewidth}{0.4pt}
\end{emailbox}

\subsubsection*{Opgave 2（私人电子邮件回复 · Minimum 90 ord · Tid: 30 minutter）}

\begin{emailbox}{Kursistark --- Skrivning, Opgave 2}
\textbf{Situation:} Du har fået e-mailen her nedenfor fra din veninde Line. Læs e-mailen, og skriv et svar.\\
\textbf{Opgave:} Skriv e-mailen til Line. Du skal svare på alle Lines spørgsmål og indlede og afslutte e-mailen på en passende måde.\\
\textbf{Længde:} Minimum 90 ord. \qquad \textbf{Tid:} 30 minutter.

\vspace{0.2em}
\hrule
\vspace{0.2em}
\hfill\textit{Odense d. 19. november 2025}\\[2pt]
\textit{Kære \ldots}

\vspace{0.15em}
\textit{Hvordan har du det? Vi har ikke skrevet sammen i lang tid, og jeg tænkte på, hvordan det går med dit liv i Danmark nu sammenlignet med i begyndelsen. Hvad overrasker dig egentlig mest i Danmark? Hvad er blevet nemmere for dig i hverdagen? Er der stadig noget ved dansk kultur eller det danske sprog, som du synes er svært? Og hvilke planer har du for de næste måneder?}

\vspace{0.15em}
\textit{Jeg glæder mig rigtig meget til at høre fra dig!}

\vspace{0.15em}
\textit{Kærlig hilsen\\
Line}

\vspace{0.15em}
\hrule
\vspace{0.15em}
\textbf{Til:} Line \qquad \textbf{Emne:} Svar: Livet i Danmark nu\\[2pt]
\rule{\linewidth}{0.4pt}\\[2pt]
\rule{\linewidth}{0.4pt}\\[2pt]
\rule{\linewidth}{0.4pt}
\end{emailbox}

% =====================================================================
% Del 3: Mundtlig kommunikation
% =====================================================================
\clearpage
\ncdsubsec{Del 3: Mundtlig kommunikation（口语交际全真题卡）}

\subsubsection*{Mundtlig Opgave 1: \textit{Emneopgave med Mindmap}}

\begin{tipbox}{官方考场规则与红线提醒（Mundtlig Opgave 1）}
考生须自带 \textbf{2 张} 已填好的 Mindmap（须选自工作/实习/学习/语言学校或市民社会话题）。\textbf{考场铁律：Mindmap 上只能写关键词或短语（stikord），绝不能写完整句子！} 考官看你的空白/关键词卡记笔记：Del 1 脱稿陈述 \textbf{1--2 分钟}，Del 2 围绕你的内容追问 \textbf{3--4 分钟}（须简要说明理由，如 \textit{fordi / derfor / synes}）。\\
\textbf{\color{NCDRed}官方禁选红线：}\textit{Min bolig/lejlighed}、\textit{Min familie} 与 \textit{Min hverdag} 因复杂度不足 \textbf{禁止用于 Modul 3 Opgave 1}！下面的 \textit{Arbejdsdeling i min familie} 聚焦分工、子女责任与家庭平权，为 SIRI 官方推荐 B1 议题。
\end{tipbox}

\begin{grammatikbox}{Mindmap 4A：Arbejdsdeling i min familie（官方 6 辐射格 · 纯 Stikord 备考卡）}
\small
\renewcommand{\arraystretch}{1.32}
\begin{tabularx}{\linewidth}{@{} >{\bfseries\color{NCDNavy}\raggedright\arraybackslash}p{4.3cm} >{\raggedright\arraybackslash}X @{}}
\toprule
\textbf{Mindmap 官方分支} & \textbf{Stikord 关键词笔记（禁止整句！）} \\
\midrule
1. hvem bor sammen & mig, min partner (Thomas), datter (8 år) · begge fuldtidsarbejde \\
2. madlavning og indkøb & skiftes hver anden dag · ugentlig madplan · fælles indkøb lørdag \\
3. rengøring og tøjvask & fast lørdagsrutine · støvsuge, badeværelse · vaske og lægge tøj sammen \\
4. børns pligter/ansvar & rydde op på værelse · dække bord · tømme opvaskemaskine \\
5. økonomi/planlægning & fælles budgetkonto · dele faste udgifter · familiekalender på mobilen \\
6. mening om ligestilling i hjemmet & vigtigt med 50/50 · begge har fritid · godt forbillede for børn \\
\bottomrule
\end{tabularx}

\vspace{0.45em}
\textbf{\color{NCDRed}Del 2 考官追问预测（Eksaminators opfølgende spørgsmål）：}
\begin{enumerate}
  \item \textit{Du siger, at I skiftes til at lave mad hver anden dag. Hvad laver du helst, når det er din tur?}
  \item \textit{Hvordan reagerer jeres datter på sine pligter --- gør hun dem altid uden problemer?}
  \item \textit{Er der nogle huslige opgaver, som du eller din partner slet ikke kan lide at lave?}
  \item \textit{Hvordan har arbejdsdelingen i dit hjem ændret sig i forhold til, hvordan du selv voksede op?}
\end{enumerate}
\end{grammatikbox}

\begin{grammatikbox}{Mindmap 4B：Mine erfaringer med at lære dansk（官方 6 辐射格 · 纯 Stikord 备考卡）}
\small
\renewcommand{\arraystretch}{1.32}
\begin{tabularx}{\linewidth}{@{} >{\bfseries\color{NCDNavy}\raggedright\arraybackslash}p{4.3cm} >{\raggedright\arraybackslash}X @{}}
\toprule
\textbf{Mindmap 官方分支} & \textbf{Stikord 关键词笔记（禁止整句！）} \\
\midrule
1. sprogskole/modul & aftenhold to gange ugentligt · DU3 Modul 3 · 16 kursister i klassen \\
2. tid og metoder & 45 min. lektier dagligt · podcast i toget · skrive nye gloser op \\
3. udtale og grammatik & svære bløde d'er og vokaler · øve ordstilling og passiv \\
4. dansk i hverdagen/job & snakke med kollegaer i frokostpausen · bestille i butikker \\
5. største udfordring & forstå hurtig tale · genert over fejl · skifte fra engelsk \\
6. mening om udvikling & meget mere selvsikker nu · stolt over fremskridt · mål: PD3 \\
\bottomrule
\end{tabularx}

\vspace{0.45em}
\textbf{\color{NCDRed}Del 2 考官追问预测（Eksaminators opfølgende spørgsmål）：}
\begin{enumerate}
  \item \textit{Du nævner, at du lytter til podcasts i toget. Hvilke podcasts eller medier kan du anbefale?}
  \item \textit{Hvad gør du, når en dansker pludselig skifter til engelsk midt i en samtale?}
  \item \textit{Hvad synes du har været det vigtigste, I har lært på Modul 3 indtil nu?}
  \item \textit{Hvordan vil du fortsætte med at bruge og forbedre dit danske efter modultesten?}
\end{enumerate}
\end{grammatikbox}

\clearpage
\subsubsection*{Mundtlig Opgave 2: \textit{Paropgave \& Individuel samtale --- EMNE: Ferie og rejser}}

\begin{emailbox}{Kursistark --- Opgaveark: EMNE: Ferie og rejser（四图互动题卡）}
\small
\renewcommand{\arraystretch}{1.35}
\begin{tabularx}{\linewidth}{@{} >{\raggedright\arraybackslash}X @{\hspace{0.45cm}} >{\raggedright\arraybackslash}X @{}}
\toprule
\textbf{\color{NCDNavy}Billede 1} & \textbf{\color{NCDNavy}Billede 2} \\
\textbf{[图片 1: 丹麦海边度假屋 Sommerhus]}\newline
\textit{En familie eller vennegruppe ved et hyggeligt dansk sommerhus tæt på kysten og sandstranden.} &
\textbf{[图片 2: 欧洲大城市文化观光]}\newline
\textit{Storbyferie i en europæisk hovedstad med museer, caféer, historiske gader og sightseeing.} \\[6pt]
\midrule
\textbf{\color{NCDNavy}Billede 3} & \textbf{\color{NCDNavy}Billede 4} \\
\textbf{[图片 3: 山林自然徒步露营]}\newline
\textit{Aktiv naturferie med vandrestøvler, rygsæk, telt eller træhytte i bjerge og skove.} &
\textbf{[图片 4: 居家休息 Staycation / 一日短途游]}\newline
\textit{Ferie hjemme med afslapning på altanen eller i haven samt små dagsture i lokalområdet med tog eller cykel.} \\
\bottomrule
\end{tabularx}
\end{emailbox}

\begin{grammatikbox}{Testerark --- 官方考官引导脚本与个人深挖三问（Opgave 2）}
\textbf{\color{NCDNavy}Del 1 (Paropgave --- ca. 4 min.)}

Opgavearket lægges foran kursisterne, der præsenteres for emnet på følgende måde:
\begin{quote}\itshape\small
Emnet er ”Ferie og rejser”. Her er nogle billeder, der viser forskellige måder at holde ferie på. I skal tale med hinanden om, hvad I synes om de forskellige måder at holde ferie på, og hvorfor. I skal stille spørgsmål til hinanden og svare på hinandens spørgsmål.
\end{quote}

Hvis kursisterne ikke udtrykker deres præferencer i forhold til \textbf{alle fire billeder}, sørger tester for, at dette sker ved at spørge:
\begin{quote}\itshape\small
Kan I lide at holde ferie i et sommerhus i Danmark / tage på storbyferie i udlandet / tage på aktiv naturferie med vandring eller camping / holde ferie hjemme med korte dagsture? Hvorfor? / Hvorfor ikke?
\end{quote}

\vspace{0.35em}
\hrule
\vspace{0.35em}

\textbf{\color{NCDRed}Del 2 (Individuel opgave --- ca. 3 min. pr. kursist)}

Tester stiller hver kursist nedenstående spørgsmål:
\begin{enumerate}
  \item \textit{Hvordan kan du bedst lide at holde ferie --- og hvorfor?}
  \item \textit{Hvad er vigtigst for dig på en ferie (fx pris, vejr, afslapning eller nye oplevelser)?}
  \item \textit{Kan du fortælle lidt om den bedste ferie, du har haft for nylig? (hvor, hvornår, med hvem, hvad lavede I)}
\end{enumerate}
{\small\textit{Hvis kursisten har besvaret et spørgsmål i Del 1, stilles uddybende spørgsmål, fx: ”Du siger, at \ldots{} vil du fortælle lidt mere om \ldots?”}}
\end{grammatikbox}

% =====================================================================
% Del 4: Retteark, Detaljeret Analyse & Højscore-Modeller
% =====================================================================
\clearpage
\ncdsubsec{Del 4: Retteark, Detaljeret Analyse \& Højscore-Modeller}

\subsubsection*{4.1 官方 \textit{Retteark} 评分卡与考前 30 秒终极自查清单}

\begin{grammatikbox}{DU3 Modul 3 --- Modultest 4 官方标准评分卡（Retteark）}
\small
\renewcommand{\arraystretch}{1.24}
\begin{tabularx}{\linewidth}{@{} >{\bfseries\color{NCDNavy}\raggedright\arraybackslash}p{3.45cm} >{\raggedright\arraybackslash}X >{\centering\arraybackslash}p{1.65cm} >{\centering\arraybackslash}p{2.55cm} @{}}
\toprule
\textbf{考核题型} & \textbf{标准答案（Rettenøgle）} & \textbf{满分} & \textbf{及格标准} \\
\midrule
Læsning Opgave 1 & \textbf{1.} hvor \quad \textbf{2.} om \quad \textbf{3.} mens \quad \textbf{4.} men \quad \textbf{5.} til \quad \textbf{6.} selvom\newline（干扰词：fordi, hvis, ikke, lige, så, under） & 6 point & 分项计分 \\
\cmidrule{1-4}
Læsning Opgave 2A & \textbf{1.} D \quad \textbf{2.} B \quad \textbf{3.} F \quad（干扰项：A, C, E） & 3 point & 分项计分 \\
\cmidrule{1-4}
Læsning Opgave 2B & \textbf{1.} E \quad \textbf{2.} A \quad \textbf{3.} C \quad（干扰项：B, D, F） & 3 point & 分项计分 \\
\cmidrule{1-4}
Læsning Opgave 3 & \textbf{1.} A \quad \textbf{2.} B \quad \textbf{3.} C \quad \textbf{4.} C \quad \textbf{5.} B & 5 point & 分项计分 \\
\cmidrule{1-4}
Læsning Opgave 4 & \textbf{1.} B \quad \textbf{2.} A \quad \textbf{3.} C \quad \textbf{4.} B \quad \textbf{5.} C \quad \textbf{6.} A \quad \textbf{7.} B & 7 point & 分项计分 \\
\cmidrule{1-4}
\textbf{\color{NCDRed}Læsning 总计} & \textbf{Opgave 1, 2A, 2B, 3 og 4 合订总评}（50 分钟 · 禁用词典） & \textbf{24 point} & \textbf{\color{NCDRed}15--24 = Bestået} \\
\midrule
Skrivning Opgave 1 & 向 \textit{NordBus} 投诉：覆盖全部 4 要点 + 语体得体（用 \textit{I / jer / jeres}）+ 句法流畅 & 合格制 & 建议目标 105--145 词 \\
\midrule
Skrivning Opgave 2 & 回复 \textit{Line} 邮件：完整回答全部 4 个问题 + 得体称呼与落款 + \textbf{minimum 90 ord}（30 min.） & 合格制 & \textbf{$\ge$ 90 词}（建议 105--140 词） \\
\midrule
Mundtlig Opgave 1 & 2 张纯 \textit{stikord} Mindmap 抽签其一：Del 1 陈述 1--2 min + Del 2 追问 3--4 min & 合格制 & 整体评估 \textit{Bestået} \\
\midrule
Mundtlig Opgave 2 & \textit{Ferie og rejser}：Del 1 双人围绕全部 4 图互动约 4 min + Del 2 个人问答约 3 min & 合格制 & 整体评估 \textit{Bestået} \\
\bottomrule
\end{tabularx}
\end{grammatikbox}

\begin{tipbox}{考前 30 秒终极自查清单（DU3 Modul 3 结业实战总检）}
作为本册四套 DU3 Modul 3 全真模考的收官卷，进考场前请逐条默念以下 5 条结业保命规则：
\begin{enumerate}
  \item \textbf{阅读 50 分钟禁用词典与 15/24 及格底线}：五部分一次发放（Opg 1 6 分 + 2A/2B 各 3 分 + Opg 3 5 分 + Opg 4 7 分），\textbf{15 分即及格}。遇到 Opg 4 千万别靠“原词匹配”秒选——一定要比对\textbf{同义转述（\textit{Parafrasering}）}与否定反转陷阱！
  \item \textbf{Opg 2A/2B 与 Opg 3 语篇衔接密码}：做聊天补全紧盯\textbf{跨回合指代与语气呼应}（如 \textit{Det har du ret i / ham / sådan noget}）；做段内补句必须“前看一句、后看信号词”（\textit{For} 引出前句原因，\textit{Men} 构成逻辑反转，\textit{Derfor/Så} 顺接结果）。
  \item \textbf{两大 B1 语法引擎与被动语态区分}：\ding{172} 一般规则、习惯或公告用 \textbf{\textit{-s} 被动}（\textit{aftaler planlægges}），具体一次性已发生事件用 \textbf{\textit{blive} 被动}（\textit{bussen blev forsinket / passagerer blev ikke informeret}）；\ding{173} 介词搭配 \textit{at/hv-} 从句时自查衔接（\textit{usikker på, om... / regne med, at...}），并确保从句内 \textit{ikke/altid/faktisk} 严格位于\textbf{动词前}。
  \item \textbf{写作双任务硬核底线}：Opgave 2 私人回信\textbf{严格要求 $\ge$ 90 词}（建议 105--135 词），必须包含礼貌称呼、落款并逐一回答朋友的全部问题；Opgave 1 投诉信或推荐信务必不漏 \textbf{4 个指定要点}，对公司机构统一称呼 \textbf{\textit{I / jer / jeres}}。
  \item \textbf{口试两条不可触碰的红线}：Opgave 1 \textbf{严禁}使用 \textit{Min bolig}、\textit{Min familie}、\textit{Min hverdag} 等低复杂度话题，Mindmap 上只写纯 \textbf{\textit{stikord}}（禁整句）；Opgave 2 双人环节必须对\textbf{全部 4 张图}均明确表达偏好并用 \textit{fordi / derfor / synes} 简述理由，同时主动向搭档发问！
\end{enumerate}
\end{tipbox}

\subsubsection*{4.2 阅读 Opgave 1--4 逐题解析、干扰项拆解、全文中文精译与核心词表}

\begin{answerbox}
\textbf{\color{NCDNavy}【Opgave 1 逐空语法/语境解析、干扰词排雷与中文精译】}
\begin{itemize}
  \item \textbf{(1) hvor}：先行词为地点 \textit{en lille landsby på Fyn}，后接存在句 \textit{der kun bor...}（“在那里只住着大约四百人”）。此处从句副词 \textit{kun} 置于动词 \textit{bor} 之前（SAV 从句语序），直接排除需接主句语序的并列连词 \textit{men / så}。
  \item \textbf{(2) om}：固定搭配 \textit{usikker på, om...}（“不确定是否……”），引导间接是非疑问从句。词库中没有 \textit{at}，且语义表不确定性，故只能选 \textbf{om}。
  \item \textbf{(3) mens}：连接两个同时进行的动作，“一边听鸟儿歌唱，\textbf{一边}（\textit{mens}）一起吃早饭”。其他连词（\textit{fordi, hvis}）均无法构成合理的伴随时间关系。
  \item \textbf{(4) men}：语义急转折，“Signe 原以为乡下人各顾各的，\textbf{但是}（\textit{men}）第一个周末邻居就带着刚出炉的面包来敲门了”，且 \textit{allerede den første weekend bankede naboerne...} 为标准主句 V2 倒装，故排除从属连词 \textit{selvom / fordi}。
  \item \textbf{(5) til}：固定动词短语 \textit{vænne sig til noget}（“习惯于某事”）。\textit{under} 不与 \textit{vænne sig} 搭配。
  \item \textbf{(6) selvom}：让步转折，“尽管去办公室路程其实比以前更久，但他们依然很满意这个决定”。注意 \textit{turen til kontoret \underline{faktisk} tager...} 中 \textit{faktisk} 位于动词 \textit{tager} 之前（标准从句语序），既在语法上封死了 \textit{men}，又在逻辑上排除了表示因果的 \textit{fordi}（没人会“因为通勤变长而高兴”）。
  \item \textbf{6 个干扰词排雷}：\textit{fordi}（(4)(6) 逻辑均不成立）、\textit{hvis}（全文为已发生事实，无需假设条件）、\textit{ikke}（无法连接从句或搭配介词）、\textit{lige}（副词无合适空位）、\textit{så}（表“所以”，与 (4) 期望反转矛盾，且不能接 (6) 从句语序）、\textit{under}（与 \textit{vænne sig} 不搭配）。
\end{itemize}

\textbf{\color{NCDTeal}[Opgave 1 中文精译]}\\
在哥本哈根市中心一间昂贵的两居室公寓住了七年后，西涅和卡斯珀决定搬到乡下去。他们早就梦想能有更大的空间和一座花园，但在首都很难以合理的价格买到房子。春天的一个周末，他们去菲英岛一个小村庄看望朋友，那里只住着大约四百人。他们立刻被宁静的气氛和镇子周围绿油油的田野迷住了。尽管两人对漫长的通勤路程有些担心，他们还是买下了一栋院子里种着果树的旧砖房。第一个月西涅还有些拿不准自己是否会想念哥本哈根的咖啡馆和文化生活。夜里街道静悄悄的，早晨在花园里一边一起吃早饭，一边就能听见鸟儿歌唱。西涅原本以为乡下人会比较各顾各的，但才到第一个周末，邻居就捧着刚烤好的面包来敲门了。如今卡斯珀每周四天坐火车通勤，他也终于习惯了多出来的这段路程，因为他每周五都能居家办公。这对伴侣至今依然很庆幸自己的决定，尽管去办公室的路程其实比以前更费时间。乡村生活给他们的日常带来了更多的宁静。

\vspace{0.45em}
\hrule
\vspace{0.45em}
\textbf{\color{NCDNavy}【Opgave 2A \& 2B 逐空语篇指代解析与中文精译】}
\begin{itemize}
  \item \textbf{Opgave 2A (Mette \& Louise)}：\textbf{1. D}（Mette 说铁路施工火车停运 $\rightarrow$ Louise 感叹“真扫兴，那得坐铁路替代巴士，慢得要命” $\rightarrow$ 紧接 Mette 下一句 \textit{Det har du desværre ret i --- det tager næsten dobbelt så lang tid med togbus} 精准回扣）；\textbf{2. B}（Mette 提议改到周日 $\rightarrow$ Louise 说周日可以但傍晚要去看哥哥来不及既看电影又逛美术馆，提议“要不放弃美术馆？” $\rightarrow$ 下句 Mette 回答 \textit{Nej, så går vi jo glip af den nye udstilling...}）；\textbf{3. F}（Mette 提议上午看美术馆、放弃看电影 $\rightarrow$ Louise 同意并反问“你来订博物馆门票吗？” $\rightarrow$ 末句 Mette 回应“太好了，我来买 10:30 的博物馆票”）。\\
  \textbf{排雷}：\textbf{A}（说“喜欢替代巴士”与下句 \textit{Det har du desværre ret i} 矛盾）；\textbf{C}（周日全天上班与后续周日约定矛盾）；\textbf{E}（还停留在周六拿电影票，与改期周日去博物馆矛盾）。
  \item \textbf{Opgave 2B (Anders \& Camilla)}：\textbf{1. E}（Anders 问早上 9:30 的客户会议怎么办 $\rightarrow$ Camilla 说让 Martin 顶替并让 Anders 把资料发给\textbf{他} $\rightarrow$ 紧接下句 \textit{Jeg sender \underline{ham} straks alle noterne...} 精准回扣）；\textbf{2. A}（Camilla 提醒在考勤系统登记“儿童首日病假” $\rightarrow$ 下句 Anders 回答 \textit{Ja, jeg har allerede registreret det i vores tidsystem}）；\textbf{3. C}（Anders 问下午有没有特别优先事项 $\rightarrow$ Camilla 拜托他在下班前完成月度预算 \textit{månedsbudgettet} $\rightarrow$ 末句 Anders 回应 \textit{Det skal jeg nok gøre færdigt før kl. 16}）。\\
  \textbf{排雷}：\textbf{B}（说 Martin 在休假与下句发资料给“他”矛盾）；\textbf{D}（让 Anders 直接休息到周一与下句下午一点上 Teams 矛盾）；\textbf{F}（擅自给 Anders 妻子打电话与妻子中午回来接班的事实冲突）。
\end{itemize}

\textbf{\color{NCDTeal}[Opgave 2A \& 2B 中文精译]}\\
\textbf{2A（梅特与露易丝）}：梅特：“嗨露易丝！期待周六去奥胡斯看电影和逛美术馆吗？”｜露易丝(Z)：“超期待！我盼了一整周了。”｜梅特：“我也是，但我刚在 Rejseplanen 查到周六全天铁路施工，火车根本不开。”｜露易丝(1-D)：“唉真可惜！那我们估计得坐替代大巴了，那简直慢得像一个世纪。”｜梅特：“你说得太对了，坐替代大巴时间几乎翻倍。要不我们改成周日去？”｜露易丝(2-B)：“周日我倒是合适！但傍晚我得去看我哥，来不及同时看电影和逛展览了，要不把美术馆取消？”｜梅特：“不行，那样就错过新展览了，我可不想错过。那我们上午去美术馆、干脆放弃电影行吗？”｜露易丝(3-F)：“行啊，这计划听起来很棒！之后还能喝杯咖啡。你订博物馆门票吗？”｜梅特：“太好了，那我来买十点半的博物馆票，周日早上车站见！”\\
\textbf{2B（安德斯与卡米拉）}：安德斯：“早安卡米拉。我儿子诺亚早上发高烧了，我今天必须请‘孩子生病第一天’假。”｜卡米拉(Z)：“早安安德斯。真遗憾，祝他早日康复！你今天还能稍微居家办公吗？”｜安德斯：“可以，我太太中午回家接替照顾，所以下午我可以居家办公。不过我 9:30 的客户会议怎么办？”｜卡米拉(1-E)：“别担心，我问问马丁能不能替你开这个会。你能把资料发给他吗？”｜安德斯：“太感谢了，真是松了一口气！我马上把所有笔记和演示文稿发到他邮箱。”｜卡米拉(2-A)：“完美！记得在我们的系统里也登记一下孩子第一天病假，好吗？”｜安德斯：“好的，我已经在考勤系统里登记了，下午一点我会登录 Teams。下午有什么需要优先处理的事吗？”｜卡米拉(3-C)：“有的，如果你能在下班前把月度预算做完，那就帮大忙了。”｜安德斯：“我一定在四点前搞定。下午 Teams 见！”

\vspace{0.45em}
\hrule
\vspace{0.45em}
\textbf{\color{NCDNavy}【Opgave 3 逐段前瞻后顾逻辑与中文精译】}
\begin{itemize}
  \item \textbf{1. Svar: A}：空前说终于拿到 8 周会计事务所实习，空后用 \textbf{\textit{For selvom...}} 解释“因为虽然她在本国做过办公室工作，但所有丹麦法规和词汇对她来说都是全新的”，故空位必须同时表达“既开心又有一点紧张”的矛盾心理（A）。B/C/D 均与拿到实习或语言陌生矛盾。
  \item \textbf{2. Svar: B}：空前介绍头几天要看旧发票、学软件，空后用 \textit{Heldigvis siger Mette, at hun altid må spørge, hvis der er et fagord...} 安慰她，说明空位必须交待“由于专业词多，比她预想的费时间”（B）。A 说“很简单”与后文安慰逻辑冲突。
  \item \textbf{3. Svar: C}：空前讲八位同事午休说话很快并开内部玩笑，空后用 \textbf{\textit{Derfor beslutter hun sig for...}} 引出她决定主动坐到友善同事身边发问，故空位须是导致她“不知何时插话”的心理原因（C）。A“喜欢独自吃饭”与后句坐到同事旁冲突。
  \item \textbf{4. Svar: C}：空前发现文件夹里少了一大笔钱，空后 Mette 称赞她细心发现了错误，说明中间缺失的正是“找了半小时终于找到丢失单据并拿给 Mette 看”（C）。A/B 隐瞒不报则不可能得到 Mette 表扬。
  \item \textbf{5. Svar: B}：空前说被叫进老板办公室、心情忐忑，空后说她“带着大大的笑容回家，觉得稳定工作的梦想真的有望实现”，故老板宣布的必然是“团队很满意，下个月给她一份兼职代班合同”（B）。
\end{itemize}

\textbf{\color{NCDTeal}[Opgave 3 中文精译]}\\
(1) 法蒂玛在欧登塞住了两年，一直梦想在丹麦从事会计和财务工作。上个月她通过了丹麦语考试，终于在当地一家会计师事务所拿到了为期八周的实习岗位。第一天上班前，她既非常开心又有点紧张。因为虽然她在家乡做过办公室工作，但所有丹麦的规章制度和专业词汇对她来说都是全新的。(2) 周一早上，她的联系人梅特热情迎接了她，带她参观办公室并把她介绍给其他员工。头几天法蒂玛需要阅读旧发票并熟悉公司的数字系统。因为很多专业词汇很难，这件事比她预想的要花更多时间。幸好梅特说，只要遇到不懂的术语或任务，她随时都可以问。(3) 两周后法蒂玛更熟练地掌握了电脑软件，但食堂里的午休时间依然让人觉得有点吃力。八位同事围坐吃饭时语速很快，还会笑一些圈内梗。这有时让她拿不准自己什么时候该开口插话。因此她决定坐到一位友善的同事旁边，主动问几个关于周末计划的问题。(4) 实习期过半时，法蒂玛获准独立核对当月收据并把数字录入账目。一天下午她突然发现其中一个文件夹里少了一大笔金额。找了半小时后，她幸好找到了丢失的凭证并拿给梅特看。梅特对法蒂玛发现这个错误赞不绝口，并表扬她非常细心。(5) 八周快结束时，法蒂玛被叫进老板办公室谈话。她很紧张老板会对她实习期间的表现作何评价。老板告诉她，因为团队非常喜欢她，决定从下个月起给她提供一份兼职替班合同。法蒂玛面带大大的笑容回到家，因为现在她知道，在丹麦找到稳定工作的梦想真的可以实现。

\vspace{0.45em}
\hrule
\vspace{0.45em}
\textbf{\color{NCDNavy}【Opgave 4 逐题同义定位与三人全文中文精译】}
\begin{itemize}
  \item \textbf{1. Svar: B (Karim)}：\textit{professionel hjælp udefra til rengøringen} $\leftrightarrow$ Karim：\textit{betale et rengøringsfirma, der kommer hver anden fredag}。
  \item \textbf{2. Svar: A (Stine)}：\textit{en fast skriftlig plan i køkkenet} $\leftrightarrow$ Stine：\textit{et fast ugeskema på køleskabet}（冰箱上的固定周计划表）。
  \item \textbf{3. Svar: C (Elena)}：\textit{overtaget al madlavningen, fordi partneren ikke var god til det} $\leftrightarrow$ Elena：Frederik 两次把菜烧糊且讨厌看食谱，现在由 Elena \textit{står for al madlavningen}。
  \item \textbf{4. Svar: B (Karim)}：\textit{bede sine børn flere gange} $\leftrightarrow$ Karim：\textit{minde dem om det flere gange}（不得不多次提醒三个青少年）。Stine 的孩子虽有固定家务但未提需反复催促。
  \item \textbf{5. Svar: C (Elena)}：\textit{irriterende, at hjemmearbejde fører til ekstra husarbejde} $\leftrightarrow$ Elena：在家办公三天时常顺手洗衣服，\textit{Det irriterer mig faktisk lidt, for hjemmearbejde er jo stadig rigtigt arbejde!}
  \item \textbf{6. Svar: A (Stine)}：\textit{køber ind..., mens partneren laver det meste af aftensmaden} $\leftrightarrow$ Stine：她负责洗衣和采购（\textit{indkøbene}），而 Nikolaj 每周做五天晚饭。Karim 自己既买菜又做工作日全套晚饭，不符“伴侣做大部分晚饭”。
  \item \textbf{7. Svar: B (Karim)}：\textit{overtog madlavningen i hverdagene, fordi partneren fik skiftende arbejdstider} $\leftrightarrow$ Karim：妻子 Sarah 当护士轮值夜班（\textit{skiftende nattevagter}）后，他开始负责工作日买菜和做热晚饭。
\end{itemize}

\textbf{\color{NCDTeal}[Opgave 4 中文精译]}\\
\textbf{A. 斯蒂娜}：以前两个孩子还很小的时候我做兼职，几乎包揽了所有家务。可那导致我和丈夫尼古拉晚上经常吵架，因为我总是累坏了。现在我们俩都全职工作，因此在冰箱上贴了一份固定的周计划表。尼古拉负责吸尘和打扫卫生间，我则负责洗衣服和采购。作为平衡，尼古拉特别喜欢下厨，所以他每周做五天晚饭，我只在周末简单做点吃的。我们六岁和九岁的孩子也有固定家务：必须自己收拾房间，并每周帮忙清空几次洗碗机。我觉得家庭平权不是去斤斤计较每一分钟，而是让父母双方晚上都能有时间喘口气放松一下。\\
\textbf{B. 卡里姆}：在我童年时都是妈妈包办所有做饭和打扫，所以刚和萨拉结婚时，我老实说也以为家里会是这样。但萨拉开始做护士、要倒班上夜班后，情况很快变了。如今几乎总是由我买菜并在每个工作日做热晚饭，因为她下班后常常精疲力竭。我们俩都不想把宝贵的周末花在擦地板上，所以选择花钱请保洁公司隔周五来打扫一次。我们的三个青少年孩子不爱做饭，不过他们轮流割草坪、倒垃圾和叠衣服——虽然我还是得提醒他们好几次。在我看来，男人在家庭里应当承担和女人一样多的责任。\\
\textbf{C. 埃琳娜}：我和丹麦男友弗雷德里克住在一起，还没有孩子。因为我每周有三天居家办公，午休时不知不觉就会去洗一桶衣服，而弗雷德里克每天要通勤两小时。这其实让我有点烦，因为居家办公也是正经上班啊！以前我们试过轮流做饭，但弗雷德里克两次把菜烧糊，而且讨厌照着食谱做。现在所有做饭都归我，因为我本来就喜欢下厨。作为交换，弗雷德里克事后总是负责洗碗、擦所有窗户，并把垃圾送去回收站。我认为好的平权是按各自最擅长的事情来分工，而不是两个人非得干一模一样的活儿。将来有了孩子，也要从小教他们搭把手。
\end{answerbox}

\subsubsection*{本卷核心生词精表（Gloser --- Modultest 4 · 全 15 词 B1 精编）}

\begin{glosertable}
名词 & en landsby & \textit{(landsbyen, landsbyer, landsbyerne)} & 村庄，乡村小镇 \\
名词 & et foreningsliv & \textit{(foreningslivet, ingen flertal)} & 社团生活，协会活动 \\
动词 & at pendle & \textit{(pendler, pendlede, har pendlet)} & 跨城通勤，往返上下班 \\
名词 & en afgang & \textit{(afgangen, afgange, afgangene)} & 发车班次，出发时间 \\
名词 & en forsinkelse & \textit{(forsinkelsen, forsinkelser, forsinkelserne)} & 延误，晚点 \\
名词 & en erstatning & \textit{(erstatningen, erstatninger, erstatningerne)} & 赔偿，补偿金 \\
名词 & en refusion & \textit{(refusionen, refusioner, refusionerne)} & 退款，退费 \\
名词 & en praktikplads & \textit{(praktikpladsen, praktikpladser, praktikpladserne)} & 实习岗位 \\
名词 & et regnskab & \textit{(regnskabet, regnskaber, regnskaberne)} & 账目，财务报表 \\
名词 & en arbejdsdeling & \textit{(arbejdsdelingen, arbejdsdelinger, arbejdsdelingerne)} & 劳动分工，家务分配 \\
名词 & ligestilling & \textit{(ligestillingen, ingen flertal)} & 平等，平权（常指性别平等） \\
名词 & en pligt & \textit{(pligten, pligter, pligterne)} & 义务，分内家务/职责 \\
形容词 & omhyggelig & \textit{(omhyggeligt, omhyggelige)} & 细心的，认真仔细的 \\
动词 & at tilpasse sig & \textit{(tilpasser sig, tilpassede sig, har tilpasset sig)} & 适应，主动调整适应 \\
名词 & en ferieform & \textit{(ferieformen, ferieformer, ferieformerne)} & 度假方式，旅游形式 \\
\end{glosertable}

\subsubsection*{4.3 写作双任务双档范文（Opgave 1 投诉信 \& Opgave 2 私人回信）}

\begin{answerbox}
\textbf{\color{NCDNavy}【Opgave 1 范文一：B1 稳过版（105 ord · 四要点清晰齐全）】}\\
\textit{Til: Kundeservice, NordBus \quad|\quad Emne: Klage over busrejse fra København til Aalborg}

\vspace{0.25em}
\textit{Kære NordBus}

\vspace{0.2em}
\textit{Jeg skriver til jer for at klage over en meget dårlig busrejse i går den 18. november. Jeg rejste med jeres bus fra København til Aalborg klokken 09.15.}

\vspace{0.2em}
\textit{For det første var bussen næsten to timer forsinket, og vi fik desværre ingen besked. Derudover var mit sæde ødelagt, og toilettet i bussen virkede ikke under turen.}

\vspace{0.2em}
\textit{Da jeg talte med chaufføren i pausen, sagde han, at han ikke kunne gøre noget. Bagefter ringede jeg til jeres kundeservice, men de ville heller ikke hjælpe mig.}

\vspace{0.2em}
\textit{Derfor håber jeg, at I vil give mig en undskyldning og betale billetten tilbage.}

\vspace{0.2em}
\textit{Med venlig hilsen\\
Sofie Jensen}

\vspace{0.25em}
\small\textbf{点评（正文 105 词）}：句式稳健零硬伤，精准兑现全部 4 个要点（乘车时间与起终点 / 延误无通知 + 坏座椅 + 故障厕所 / 司机与客服推诿反应 / 道歉与全额退票诉求），全程对客运公司统一使用 \textit{I / jer / jeres}，正确运用从句否定前置（\textit{at han ikke kunne}）。

\vspace{0.45em}
\hrule
\vspace{0.45em}

\textbf{\color{NCDRed}【Opgave 1 范文二：B1+ 高分示范版（137 ord · 亮点句法全调度）】}\\
\textit{Til: Kundeservice, NordBus \quad|\quad Emne: Klage over forsinkelse og dårlig service}

\vspace{0.25em}
\textit{Kære NordBus}

\vspace{0.2em}
\textit{Jeg vil hermed klage over en meget utilfredsstillende rejse mellem København og Aalborg i går den 18. november med afgang kl. 09.15.}

\vspace{0.2em}
\textit{Allerede fra starten opstod der problemer. Bussen \da{blev} nemlig \da{forsinket} med næsten to timer, uden at passagererne \da{blev informeret} undervejs. Da vi endelig kørte, opdagede jeg, at mit sæde var knækket, og at toilettet var låst på hele den lange tur.}

\vspace{0.2em}
\textit{Da jeg gjorde chaufføren opmærksom på fejlene, reagerede han desværre uhøfligt og bad mig om at sætte mig igen. Bagefter kontaktede jeg jeres kundeservice, men der \da{blev ikke taget} ansvar for situationen, \dabold{hvilket} jeg er meget skuffet over.}

\vspace{0.2em}
\textit{\dabold{Har man betalt} fuld pris for en behagelig rejse, \dabold{forventer man} ordentlige forhold. Jeg håber derfor, at I vil give mig en fuld refusion af billetten samt et gavekort.}

\vspace{0.2em}
\textit{Med venlig hilsen\\
Sofie Jensen}

\vspace{0.25em}
\small\textbf{高分点评（正文 137 词）}：全面调度 B1 高阶句法：\ding{172} 连续精准使用 \textit{blive} 被动（\textit{blev forsinket, blev informeret, blev ikke taget ansvar}）客观陈述具体投诉事件；\ding{173} \textit{uden at} 否定伴随从句与 \textit{bad mig om at} 介词搭配；\ding{174} 整句关系从句 \textit{..., hvilket jeg er meget skuffet over} 提升逻辑连贯度；\ding{175} 末段以“无 \textit{hvis} 条件倒装”（\textit{Har man betalt..., forventer man...}）收尾索赔，正式得体且说服力极强。

\vspace{0.45em}
\hrule
\vspace{0.45em}

\textbf{\color{NCDNavy}【Opgave 2 范文一：B1 稳过版（111 ord · 远超 90 词底线）】}\\
\textit{Til: Line \quad|\quad Emne: Svar: Livet i Danmark nu}

\vspace{0.25em}
\textit{Kære Line}

\vspace{0.2em}
\textit{Tak for din mail! Det er dejligt at høre fra dig.}

\vspace{0.2em}
\textit{Det, der overrasker mig mest i Danmark, er, hvor meget folk stoler på hinanden, og at man bruger fornavn til sin chef.}

\vspace{0.2em}
\textit{Heldigvis er mange ting blevet nemmere nu. Jeg forstår endelig MitID og e-Boks, og jeg tør tale dansk, når jeg handler ind eller snakker med naboerne.}

\vspace{0.2em}
\textit{Til gengæld synes jeg stadig, at den hurtige udtale er svær, når flere danskere taler sammen. Det er også svært at få danske venner uden for skolen.}

\vspace{0.2em}
\textit{I de næste måneder skal jeg til modultest i dansk, og bagefter vil jeg søge en praktikplads på et kontor.}

\vspace{0.2em}
\textit{Kærlig hilsen\\
Maya}

\vspace{0.25em}
\small\textbf{点评（正文 111 词）}：自然呼应 Line 的全部 4 个问题（最惊讶之处 / 变容易的事 / 仍觉得困难之处 / 未来数月计划），字数稳超 90 词要求（不含抬头标题已达 111 词），称呼与落款温馨得体。

\vspace{0.45em}
\hrule
\vspace{0.45em}

\textbf{\color{NCDRed}【Opgave 2 范文二：B1+ 高分示范版（132 ord · 文化观察与句法示范）】}\\
\textit{Til: Line \quad|\quad Emne: Svar: Livet i Danmark nu}

\vspace{0.25em}
\textit{Kære Line}

\vspace{0.2em}
\textit{Hvor er det hyggeligt at høre fra dig! Mit liv i Danmark føles meget anderledes nu end i begyndelsen.}

\vspace{0.2em}
\textit{\da{Det, der} stadig overrasker mig mest her, er den store tillid og det flade hierarki. Man bruger fornavn til chefen, og forældre lader trygt barnevogne stå ude på gaden.}

\vspace{0.2em}
\textit{Samtidig er hverdagen blevet væsentlig nemmere. Jeg behøver ikke længere at oversætte hver besked i e-Boks, og fordi jeg læser danske nyheder, føler jeg mig langt mere hjemme.}

\vspace{0.2em}
\textit{Alligevel synes jeg, at det sociale liv er svært. \dabold{Til trods for at} danskerne er venlige på jobbet, \da{planlægges} næsten alle aftaler uger i forvejen, og de bløde d'er driller mig stadigvæk!}

\vspace{0.2em}
\textit{I de næste måneder skal jeg bestå Modul 3, og derefter \da{regner jeg med, at} jeg skal i virksomhedspraktik.}

\vspace{0.2em}
\textit{Kærlig hilsen\\
Maya}

\vspace{0.25em}
\small\textbf{高分点评（正文 132 词）}：紧扣 NCD 3 文化深读亮点：\ding{172} 精准点出 \textit{tillid} 与 \textit{det flade hierarki}；\ding{173} 运用习惯性 \textit{-s} 被动 \textit{planlægges} 对比丹麦人的社交预约习惯；\ding{174} 亮出高级让步连词 \textit{Til trods for at} 与“介词 + \textit{at}-从句”（\textit{regner jeg med, at...}），自然流畅地达标 B1+ 水准。
\end{answerbox}

\subsubsection*{4.4 口语示范稿（Mindmap 4A \& 4B 约 150 词陈述 + 语音 + Del 2 问答与 Opgave 2 四图互动对话）}

\begin{answerbox}
\textbf{\color{NCDNavy}【Mindmap 4A：Arbejdsdeling i min familie（Del 1 约 150 词脱稿陈述示范 + 核心音标）】}

\vspace{0.18em}
\textit{Jeg vil gerne fortælle om arbejdsdelingen \ipa{ˈɑːbɑjdsˌdeːleŋən} i min familie. Jeg bor i en lejlighed sammen med min partner Thomas og vores datter på otte år. Da vi begge arbejder på fuldtid, er vi nødt til at samarbejde i hverdagen.}

\vspace{0.18em}
\textit{Med hensyn til madlavning og indkøb \ipa{ˈenˌkʰøːˀb} skiftes vi til at lave aftensmad hver anden dag. Hver søndag laver vi en madplan, og om lørdagen handler vi sammen. Rengøringen klarer vi med en fast lørdagsrutine: Thomas støvsuger og gør badeværelset rent, mens jeg vasker og lægger tøj sammen.}

\vspace{0.18em}
\textit{Vores datter har også faste pligter. Hun skal rydde op på sit værelse, dække bord og tømme opvaskemaskinen \ipa{ˈʌbvasgəmaˌsgiːnən}, før hun ser fjernsyn. Økonomien \ipa{økonoˈmiːˀən} styrer vi via en fælles budgetkonto til faste udgifter, og vi deler en digital familiekalender på mobilen.}

\vspace{0.18em}
\textit{Jeg synes, ligestilling \ipa{ˈliːəˌsdeleŋ} i hjemmet er meget vigtigt. Det handler om, at begge voksne har fritid, og at vi giver vores datter et godt forbillede.}

\vspace{0.2em}
\small\textbf{[参考译文]}：我想谈谈我家的家务分工。我和伴侣托马斯以及我们八岁的女儿住在一套公寓里。因为我们俩都全职工作，平日里必须互相配合。在做饭和采购方面，我们隔天轮流做晚饭；每周日一起制定食谱，周六一起去采购。打扫卫生我们也有固定的周六分工流程：托马斯负责吸尘和清理卫生间，我负责洗衣服和叠衣服。我们的女儿也有固定任务：看电视前必须整理房间、摆好餐桌并清空洗碗机。财务方面我们用共同的预算账户管理固定开支，并在手机上共享数字家庭日历。我认为家庭平权非常重要，这意味着两位成年人都拥有空闲时间，同时也能给女儿树立好榜样。

\vspace{0.25em}
\textbf{\color{NCDRed}Del 2 考官追问与示范回答：}
\begin{enumerate}
  \item \textbf{Q:} \textit{Du siger, at I skiftes til at lave mad hver anden dag. Hvad laver du helst, når det er din tur?}\\
  \textbf{A:} \textit{Når det er min tur, laver jeg helst nemme, varme retter som kylling i karry eller pastaretter med masser af grøntsager, fordi vi alle tre kommer trætte hjem.}
  \item \textbf{Q:} \textit{Hvordan reagerer jeres datter på sine pligter --- gør hun dem altid uden problemer?}\\
  \textbf{A:} \textit{Ikke altid! Nogle gange må vi minde hende om opvaskemaskinen et par gange, men som regel gør hun det gerne, især når vi sætter glad musik på.}
  \item \textbf{Q:} \textit{Er der nogle huslige opgaver, som du eller din partner slet ikke kan lide at lave?}\\
  \textbf{A:} \textit{Ja, ingen af os kan lide at pudse vinduer eller stryge skjorter, så vi stryger faktisk næsten aldrig tøj i hverdagen.}
  \item \textbf{Q:} \textit{Hvordan har arbejdsdelingen i dit hjem ændret sig i forhold til, hvordan du selv voksede op?}\\
  \textbf{A:} \textit{Da jeg var barn, lavede min mor næsten alt i køkkenet alene. I dag er det heldigvis meget mere naturligt, at mænd og kvinder deler både ansvaret og opgaverne ligeligt.}
\end{enumerate}

\vspace{0.4em}
\hrule
\vspace{0.4em}

\textbf{\color{NCDNavy}【Mindmap 4B：Mine erfaringer med at lære dansk（Del 1 约 150 词脱稿陈述示范 + 核心音标）】}

\vspace{0.18em}
\textit{Jeg vil fortælle om mine erfaringer \ipa{ɐˈfɑːˀeŋɐ} med at lære dansk. Jeg går på et aftenhold på sprogskolen to gange om ugen på DU3, Modul 3, hvor vi er seksten kursister i klassen fra forskellige lande.}

\vspace{0.18em}
\textit{Ud over undervisningen bruger jeg cirka femogfyrre minutter dagligt på lektier. Jeg lytter til danske podcasts i toget, læser korte danske nyheder online og skriver nye gloser ned i en notesbog. Det sværeste ved sproget er helt klart udtalen \ipa{ˈuðˌtˢɑːlən} --- især de bløde d'er og vokalerne. I grammatikken øver jeg meget på ordstilling \ipa{ˈoɐ̯ˌsdeleŋ} i ledsætninger og passivformer.}

\vspace{0.18em}
\textit{I hverdagen bruger jeg dansk så meget som muligt, når jeg handler i butikker eller snakker med kollegaer i frokostpausen. Min største udfordring \ipa{ˈuðˌfɒːdʁeŋ} er, når folk taler hurtigt, fordi jeg nogle gange bliver genert og skifter til engelsk. Alligevel er jeg stolt af min udvikling \ipa{ˈuðˌveɡleŋ}, for jeg føler mig mere selvsikker nu og håber at nå PD3.}

\vspace{0.2em}
\small\textbf{[参考译文]}：我想分享一下自己学习丹麦语的经历。我在语言学校每周上两次夜校班，目前在 DU3 Modul 3，班里有十六位来自不同国家的同学。除了课堂学习，我每天大约花四十五分钟做作业。我会在火车上听丹麦语播客、在网上读简短的丹麦语新闻，并把新单词记在笔记本里。这门语言最难的部分无疑是发音——尤其是软 d 和元音。语法上我重点练习从句语序和被动语态。日常生活中，我在商店购物或午休和同事聊天时尽量多用丹麦语。最大的挑战是别人语速快的时候，因为有时我会害羞并切回英语。尽管如此，我对自己的进步感到自豪，现在自信多了，希望能顺利达到 PD3。

\vspace{0.25em}
\textbf{\color{NCDRed}Del 2 考官追问与示范回答：}
\begin{enumerate}
  \item \textbf{Q:} \textit{Du nævner, at du lytter til podcasts i toget. Hvilke podcasts eller medier kan du anbefale?}\\
  \textbf{A:} \textit{Jeg kan varmt anbefale DR's nyheder på let dansk samt podcasten ”Dansk i ørerne”, fordi værten taler tydeligt og forklarer spændende ting om dansk kultur.}
  \item \textbf{Q:} \textit{Hvad gør du, når en dansker pludselig skifter til engelsk midt i en samtale?}\\
  \textbf{A:} \textit{Før i tiden skiftede jeg bare med, men nu smiler jeg og siger pænt: ”Jeg vil faktisk rigtig gerne øve mit dansk, så lad os fortsætte på dansk!” Det respekterer folk altid.}
  \item \textbf{Q:} \textit{Hvad synes du har været det vigtigste, I har lært på Modul 3 indtil nu?}\\
  \textbf{A:} \textit{Det vigtigste for mig har været at kunne skrive formelle klager og fortælle sammenhængende om samfundsemner i stedet for bare at bruge korte hovedsætninger.}
  \item \textbf{Q:} \textit{Hvordan vil du fortsætte med at bruge og forbedre dit danske efter modultesten?}\\
  \textbf{A:} \textit{Jeg vil gerne starte som frivillig i en lokal idrætsklub eller læsekreds, så jeg får endnu flere danske bekendte at snakke med i min fritid.}
\end{enumerate}
\end{answerbox}

\vspace{0.35em}
\textbf{\color{NCDNavy}\large 【Mundtlig Opgave 2：双人四图互动示范逐字稿（EMNE: Ferie og rejser）】}

\dialogturn{1}{Kursist A}{Skal vi starte med billede et? Jeg elsker virkelig at holde ferie i et sommerhus ved kysten, fordi der er så roligt, og man kan hygge sig med familien uden stress. Hvad synes du om sommerhusferie?}{ˈsʌmɐˌhuːˀs / ˈkʰøsdən / ˈhyɡə sɑj}{我们先从图一开始好吗？我真的很喜欢去海边度假屋度假，因为那里很安静，可以毫无压力地和家人共度温馨时光。你觉得度假屋旅行怎么样？}

\dialogturn{2}{Kursist B}{Jeg kan også rigtig godt lide sommerhuse, især om sommeren, når man kan gå ned til stranden og grille om aftenen. Men hvis det regner hele ugen, kan det godt blive lidt kedeligt. Hvad med billede to --- kan du lide at tage på storbyferie?}{ˈsdoɐ̯ˌbyːˌfeɐ̯jə / ˈʁɑjnɐ}{我也很喜欢度假屋，特别是夏天可以去海边散步、晚上烧烤。不过如果连下一整周雨，可能会有点无聊。那图二呢——你喜欢大城市度假吗？}

\dialogturn{3}{Kursist A}{Ja, absolut! Jeg foretrækker faktisk storbyferie i Europa, fordi jeg interesserer mig meget for kunst, museer og god mad. Selvom det kan være lidt dyrt og larmende, oplever man altid noget nyt. Foretrækker du storbyferie eller naturferie som på billede tre?}{foɐˈtˢʁɛɡɐ / muˈseːɐ / nɑˈtˢuːɐ̯ˌfeɐ̯jə}{是的，非常喜欢！其实我更偏爱欧洲大城市旅行，因为我对艺术、博物馆和美食很感兴趣。虽然可能有点贵又有点吵，但总能体验到新鲜事物。你更喜欢大城市度假还是像图三那样的自然度假？}

\dialogturn{4}{Kursist B}{Jeg er faktisk helt vild med aktiv naturferie som på billede tre! Når man vandrer i bjergene eller sover i telt i skoven, får man masser af frisk luft og glemmer arbejdet fuldstændig. Det kræver selvfølgelig godt udstyr, og man kan blive øm i benene, men jeg synes, det er det hele værd. Har du nogensinde prøvet camping?}{ˈvɑndʁɐ / ˈuðˌsdyːɐ̯ˀ / ˈømˀ}{其实我超爱图三这种户外自然度假！在山里徒步或在森林里露营时，能呼吸大量新鲜空气，把工作彻底抛在脑后。当然这需要好装备，腿也可能会酸，但我觉得完全值得。你试过露营吗？}

\dialogturn{5}{Kursist A}{Kun én gang, og ærligt talt savnede jeg en blød seng og et varmt bad! Hvad synes du egentlig om det sidste billede med staycation --- altså at holde ferie derhjemme og tage på korte dagsture?}{ˈæɐ̯lid ˈtˢɑːˀld / ˈdɑwsˌtˢuːɐ}{只试过一次，说实话我很想念软床和热水澡！那你对最后一张图——居家度假，也就是待在家里然后去短途一日游——怎么看？}

\dialogturn{6}{Kursist B}{Jeg synes faktisk, det er en ret smart og klimavenlig løsning, hvis man ikke har så mange penge eller bare trænger til at sove længe. Man sparer meget transporttid, men risikoen er, at man bare kommer til at ordne vasketøj derhjemme. Hvad tænker du om at holde ferie hjemme?}{ˈkʰliːmaˌvɛnli / ˈvɑsgəˌtˢʌj}{其实我觉得如果预算不多或者只是想睡个懒觉，这是个既聪明又环保的选择。能省下很多路途时间，但风险是一不小心就在家顺手洗起衣服做起家务了。你觉得在家度假怎么样？}

\dialogturn{7}{Kursist A}{Det er jeg helt enig i! Hvis jeg bliver hjemme hele ferien, føles det ikke rigtig som en pause. Derfor vil jeg helst rejse et par dage ud af byen --- selv bare til en anden landsdel i Danmark.}{ˈeːni / ˈlɑnsˌdeːˀl}{我完全同意！如果整个假期都待在家里，感觉不太像真正放假。所以我还是更愿意离开城市出去玩几天——哪怕只是去丹麦的另一个地区也好。}

\begin{answerbox}
\textbf{\color{NCDRed}【Opgave 2 --- Del 2 个人深入三问示范回答（Individuel opgave）】}
\begin{enumerate}
  \item \textbf{Q:} \textit{Hvordan kan du bedst lide at holde ferie --- og hvorfor?}\\
  \textbf{A:} \textit{Jeg kan bedst lide en blanding af storbyferie og lidt afslapning ved havet. Jeg elsker at se seværdigheder og prøve lokale retter, men efter et par travle dage har jeg også brug for at sidde stille med en god bog.}
  \item \textbf{Q:} \textit{Hvad er vigtigst for dig på en ferie (fx pris, vejr, afslapning eller nye oplevelser)?}\\
  \textbf{A:} \textit{For mig er nye oplevelser og godt selskab allervigtigst. Vejret kan man jo ikke styre --- især ikke i Danmark --- men hvis man rejser sammen med gode venner eller familie og har et rimeligt budget, får man næsten altid en god tur.}
  \item \textbf{Q:} \textit{Kan du fortælle lidt om den bedste ferie, du har haft for nylig? (hvor, hvornår, med hvem, hvad lavede I)}\\
  \textbf{A:} \textit{Sidste efterår tog min mand og jeg fire dage til Lissabon i Portugal. Vi havde fantastisk solskin på tyve grader i oktober, og vi gik rundt i de gamle gader med sporvogn, lyttede til fado-musik og spiste frisk fisk ved havnen. Det bedste var, at vi slet ikke havde travlt og bare nød stemningen sammen.}
\end{enumerate}
\end{answerbox}


\end{document}
